% Probabilités, variables aléatoires et loi binomiale — Mathématiques Tle D — Cours signé OnBuch+
% Programme officiel MINESEC. Compiler avec : tectonic M11-probabilites.tex
\newcommand{\DOCMATIERE}{Mathématiques}
\newcommand{\DOCNIVEAU}{Tle D}
\newcommand{\DOCMODULE}{Organisation et gestion des données}
\newcommand{\DOCLECON}{11}
\newcommand{\DOCTITRE}{Probabilités, variables aléatoires et loi binomiale}
% OnBuch+ — préambule « cours signés » ENRICHI (compilé avec tectonic / XeLaTeX)
% Système visuel « Pop » d'OnBuch+ : fond crème (jamais blanc), bordures encre
% épaisses, ombres « dures » décalées, titres Archivo Black en capitales.
% Version MATHÉMATIQUES : graphes (pgfplots/tikz), boîtes pédagogiques
% (définition, propriété, méthode, attention, le savais-tu, exercice résolu,
% exercice, TP, bilan), page de garde et signature OnBuch+ ancrée en bas.
%
% Macros attendues dans main.tex AVANT \input{preamble} :
%   \DOCMATIERE  \DOCNIVEAU  \DOCMODULE  \DOCLECON (numéro)  \DOCTITRE
\documentclass[11pt]{article}

\usepackage{fontspec}
\usepackage[french]{babel}
\usepackage[a4paper,margin=2cm,top=3.4cm,bottom=2.8cm,headheight=1.4cm,footskip=1.3cm]{geometry}
\usepackage{amsmath,amssymb,amsfonts}
\usepackage{mathtools} % \xmapsto, \coloneqq... souvent utilisés par les modèles
\usepackage{cancel} % simplifications barrées
% \abs{z} (module, valeur absolue) et \norm{u} (norme), utilisés par les modèles
\providecommand{\abs}{}\let\abs\relax\DeclarePairedDelimiter{\abs}{\lvert}{\rvert}
\providecommand{\norm}{}\let\norm\relax\DeclarePairedDelimiter{\norm}{\lVert}{\rVert}
\usepackage[table]{xcolor} % [table] : fournit aussi \rowcolors (alternance de lignes)
\usepackage{pagecolor}
\usepackage{tikz}
\usetikzlibrary{arrows.meta,positioning,calc,shapes.geometric,shapes.misc,shapes.arrows,decorations.pathmorphing,decorations.markings,patterns,shadows,mindmap,trees,fit,backgrounds,matrix,intersections}
\usepackage{pgfplots}
\usepackage{pgf-pie} % diagrammes circulaires \pie (statistiques)
\pgfplotsset{compat=1.18}
\usepgfplotslibrary{fillbetween,groupplots} % aires entre courbes (calcul intégral)
\usepackage{enumitem}
\usepackage{fancyhdr}
% [most] charge toutes les bibliothèques tcolorbox (listings, minted, raster,
% documentation...) et gonfle inutilement la mémoire TeX sur un document long
% (~300 boîtes) : on ne charge que ce qui est réellement utilisé.
\usepackage[skins,breakable]{tcolorbox}
\usepackage{graphicx}
\usepackage{colortbl}
\usepackage{booktabs}
\usepackage{tabularx}
\usepackage{diagbox} % case d'en-tête diagonale des tableaux à double entrée
% Colonnes extensibles centrée / gauche / droite (C, L, R), souvent utilisées par les modèles
\newcolumntype{C}{>{\centering\arraybackslash}X}
\newcolumntype{L}{>{\raggedright\arraybackslash}X}
\newcolumntype{R}{>{\raggedleft\arraybackslash}X}
\usepackage{array}
\usepackage{multicol}
\usepackage{siunitx}
% parse-numbers=false : accepte aussi \SI{2,5 \times 10^{-3}}{...}
\sisetup{locale=FR,output-decimal-marker={,},group-minimum-digits=4,parse-numbers=false}
\DeclareSIUnit\ohm{\ensuremath{\symup{\Omega}}} % U+2126 absent de la police
\usepackage{qrcode}
% \degree : commande fréquemment utilisée par les modèles pour un angle ou une
% température en dehors de \SI{}{} (non fournie par les paquets chargés).
\providecommand{\degree}{^{\circ}}
% \qed (fin de démonstration), utilisé par les modèles sans amsthm
\providecommand{\qed}{\hfill$\square$}
% \barre : double barre des valeurs interdites dans un tableau de variations (tkz-tab)
\providecommand{\barre}{\ensuremath{\|}}
% environnement proof (amsthm non chargé) utilisé par les modèles
\ifdefined\proof\else\newenvironment{proof}[1][Démonstration]{\par\noindent\textit{#1.}\ }{\qed\par}\fi
\usepackage{lastpage}
\usepackage{hyperref}
\hypersetup{hidelinks}

% ============================================================
% Palette « Pop » OnBuch+ — mêmes hex que la classe P (plus_ui.dart)
% ============================================================
\definecolor{poporange}{HTML}{F59321}
\definecolor{popdark}{HTML}{E07A0C}
\definecolor{popcream}{HTML}{FFF6E8}
\definecolor{popink}{HTML}{1C1714}
\definecolor{poppurple}{HTML}{7A5AE0}
\definecolor{popgreen}{HTML}{2E9E6A}
\definecolor{poppink}{HTML}{E0457B}
\definecolor{popblue}{HTML}{2F7DE1}
\definecolor{popgold}{HTML}{C98A12}
\definecolor{popmuted}{HTML}{8A7F76}
\definecolor{popline}{HTML}{EADFD2}
\definecolor{popgray}{HTML}{8A7F76}
\colorlet{popgrey}{popgray}
% Alias tolérants (le modèle nomme parfois les couleurs différemment)
\colorlet{popwhite}{white}
\colorlet{popblack}{popink}
\colorlet{popred}{poppink}
\colorlet{popyellow}{popgold}
% Teintes claires pour fonds de tableaux / zones
\colorlet{popblueL}{popblue!10!white}
\colorlet{poporangeL}{poporange!14!white}
\colorlet{popgreenL}{popgreen!12!white}
\colorlet{poppurpleL}{poppurple!12!white}
\colorlet{poppinkL}{poppink!12!white}
\colorlet{popgoldL}{popgold!14!white}

\pagecolor{popcream}

% ============================================================
% Typographie
% ============================================================
\defaultfontfeatures{Ligatures=TeX}
\setmainfont{PlusJakartaSans-Medium.ttf}[Path=../../../fonts/,BoldFont=PlusJakartaSans-Bold.ttf]
% Maths en sans-serif (Fira Math) avec chiffres et lettres droites de Plus
% Jakarta, pour que les formules se fondent dans le texte.
\usepackage[math-style=ISO,bold-style=ISO]{unicode-math}
\setmathfont{FiraMath-Regular.otf}[Path=../../../fonts/]
\setmathfont{PlusJakartaSans-Medium.ttf}[Path=../../../fonts/,range={up/num,up/{Latin,latin},\mathup}]
\usepackage{icomma} % virgule décimale sans espace en mode math
% Caractères mathématiques Unicode absents des polices de texte (Plus Jakarta,
% Archivo Black) : rendus par leur équivalent mathématique, en texte comme en
% formule — sinon ils s'affichent en carré vide (ex. « Arithmétique dans ℤ »).
\usepackage{newunicodechar}
\newunicodechar{ℝ}{\ensuremath{\mathbb{R}}}
\newunicodechar{ℤ}{\ensuremath{\mathbb{Z}}}
\newunicodechar{ℂ}{\ensuremath{\mathbb{C}}}
\newunicodechar{ℕ}{\ensuremath{\mathbb{N}}}
\newunicodechar{ℚ}{\ensuremath{\mathbb{Q}}}
\newunicodechar{𝔻}{\ensuremath{\mathbb{D}}}
\newunicodechar{α}{\ensuremath{\alpha}}
\newunicodechar{β}{\ensuremath{\beta}}
\newunicodechar{γ}{\ensuremath{\gamma}}
\newunicodechar{δ}{\ensuremath{\delta}}
\newunicodechar{ε}{\ensuremath{\varepsilon}}
\newunicodechar{θ}{\ensuremath{\theta}}
\newunicodechar{λ}{\ensuremath{\lambda}}
\newunicodechar{μ}{\ensuremath{\mu}}
\newunicodechar{σ}{\ensuremath{\sigma}}
\newunicodechar{τ}{\ensuremath{\tau}}
\newunicodechar{φ}{\ensuremath{\varphi}}
\newunicodechar{ψ}{\ensuremath{\psi}}
\newunicodechar{ω}{\ensuremath{\omega}}
\newunicodechar{Σ}{\ensuremath{\Sigma}}
% majuscules produites par \MakeUppercase dans les titres (α-aminés -> Α)
\newunicodechar{Α}{\ensuremath{\alpha}}
\newunicodechar{Β}{\ensuremath{\beta}}
\newunicodechar{Γ}{\ensuremath{\Gamma}}
\newunicodechar{Δ}{\ensuremath{\Delta}}
\newunicodechar{Ε}{\ensuremath{\varepsilon}}
\newunicodechar{Θ}{\ensuremath{\Theta}}
\newunicodechar{Λ}{\ensuremath{\Lambda}}
\newunicodechar{Μ}{\ensuremath{\mu}}
\newunicodechar{Τ}{\ensuremath{\tau}}
\newunicodechar{Φ}{\ensuremath{\Phi}}
\newunicodechar{Ψ}{\ensuremath{\Psi}}
\newunicodechar{Ω}{\ensuremath{\Omega}}
\newunicodechar{↦}{\ensuremath{\mapsto}}
\newunicodechar{∈}{\ensuremath{\in}}
\newunicodechar{∉}{\ensuremath{\notin}}
\newunicodechar{∧}{\ensuremath{\wedge}}
\newunicodechar{⇔}{\ensuremath{\Leftrightarrow}}
\newunicodechar{⟺}{\ensuremath{\Longleftrightarrow}}
\newunicodechar{⇒}{\ensuremath{\Rightarrow}}
\newunicodechar{⟹}{\ensuremath{\Longrightarrow}}
\newunicodechar{∘}{\ensuremath{\circ}}
\newunicodechar{∪}{\ensuremath{\cup}}
\newunicodechar{∩}{\ensuremath{\cap}}
\newunicodechar{≡}{\ensuremath{\equiv}}
\newunicodechar{⊂}{\ensuremath{\subset}}
\newunicodechar{⊥}{\ensuremath{\perp}}
\newunicodechar{∃}{\ensuremath{\exists}}
\newunicodechar{∀}{\ensuremath{\forall}}
\newunicodechar{∛}{\ensuremath{\sqrt[3]{\,}}}
\newunicodechar{∜}{\ensuremath{\sqrt[4]{\,}}}
\newunicodechar{⃗}{\ensuremath{{}^{\to}}}
\newunicodechar{ⁿ}{\ifmmode^{n}\else\textsuperscript{\ensuremath{n}}\fi}
\newunicodechar{ˣ}{\ifmmode^{x}\else\textsuperscript{\ensuremath{x}}\fi}
\newunicodechar{ᵇ}{\ifmmode^{b}\else\textsuperscript{\ensuremath{b}}\fi}
\newunicodechar{ʳ}{\ifmmode^{r}\else\textsuperscript{\ensuremath{r}}\fi}
\newunicodechar{ᵘ}{\ifmmode^{u}\else\textsuperscript{\ensuremath{u}}\fi}
\newunicodechar{ʸ}{\ifmmode^{y}\else\textsuperscript{\ensuremath{y}}\fi}
\newunicodechar{ᵃ}{\ifmmode^{a}\else\textsuperscript{\ensuremath{a}}\fi}
\newunicodechar{ᵉ}{\ifmmode^{e}\else\textsuperscript{\ensuremath{e}}\fi}
\newunicodechar{ᵏ}{\ifmmode^{k}\else\textsuperscript{\ensuremath{k}}\fi}
\newunicodechar{ᵖ}{\ifmmode^{p}\else\textsuperscript{\ensuremath{p}}\fi}
\newunicodechar{ᴺ}{\ifmmode^{N}\else\textsuperscript{\ensuremath{N}}\fi}
\newunicodechar{ᶜ}{\ifmmode^{c}\else\textsuperscript{\ensuremath{c}}\fi}
\newunicodechar{ʰ}{\ifmmode^{h}\else\textsuperscript{\ensuremath{h}}\fi}
\newunicodechar{ᵠ}{\ifmmode^{q}\else\textsuperscript{\ensuremath{q}}\fi}
\newunicodechar{ₙ}{\ifmmode_{n}\else\textsubscript{\ensuremath{n}}\fi}
\newunicodechar{ₐ}{\ifmmode_{a}\else\textsubscript{\ensuremath{a}}\fi}
\newunicodechar{ᵢ}{\ifmmode_{i}\else\textsubscript{\ensuremath{i}}\fi}
\newunicodechar{ᵣ}{\ifmmode_{r}\else\textsubscript{\ensuremath{r}}\fi}
\newunicodechar{ₖ}{\ifmmode_{k}\else\textsubscript{\ensuremath{k}}\fi}
\newunicodechar{ᵦ}{\ifmmode_{\beta}\else\textsubscript{\ensuremath{\beta}}\fi}
\newunicodechar{ₓ}{\ifmmode_{x}\else\textsubscript{\ensuremath{x}}\fi}
\newfontfamily\popdisplay{ArchivoBlack-Regular.ttf}[Path=../../../fonts/]
\newfontfamily\poplogo{SpaceGrotesk-Bold.ttf}[Path=../../../fonts/]
\newfontfamily\popsemi{PlusJakartaSans-SemiBold.ttf}[Path=../../../fonts/]
\newfontfamily\popxbold{PlusJakartaSans-ExtraBold.ttf}[Path=../../../fonts/]
\newfontfamily\popxboldls{PlusJakartaSans-ExtraBold.ttf}[Path=../../../fonts/,LetterSpace=3.0]

\setlist[enumerate]{leftmargin=1.7em,itemsep=5pt,topsep=4pt}
\setlist[itemize]{leftmargin=1.7em,itemsep=4pt,topsep=4pt,label={\color{poporange}\textbullet}}
\setlength{\parindent}{0pt}
\setlength{\parskip}{5pt}
\renewcommand{\arraystretch}{1.25}

% pgfplots : style Pop par défaut
\pgfplotsset{
  every axis/.append style={axis line style={popink,line width=1pt},tick style={popink},
    grid style={popline,line width=0.5pt},label style={font=\small},tick label style={font=\footnotesize}},
  popcurve/.style={poporange,line width=1.8pt,no markers,smooth},
}
% Même style disponible dans un \draw TikZ simple (hors pgfplots)
\tikzset{popcurve/.style={poporange,line width=1.8pt}}

% ============================================================
% Pill / squiggle
% ============================================================
\newcommand{\poppill}[3]{%
  \tikz[baseline=-0.55ex]\node[draw=popink,line width=1.2pt,rounded corners=2.6pt,%
    fill=#2,inner xsep=6.5pt,inner ysep=3pt]{%
    \popxboldls\fontsize{7}{7}\selectfont\color{#3}\MakeUppercase{#1}};%
}
\newcommand{\popsquiggle}[1]{%
  \tikz[baseline=-0.4ex]{
    \draw[#1,line width=2.6pt,line cap=round]
      plot [smooth] coordinates {(0,0.09) (0.34,0.01) (0.68,0.21) (1.02,0.01) (1.36,0.10)};
  }%
}
% Mot-clé surligné (marqueur orange sous le texte)
\newcommand{\cle}[1]{{\popxbold\color{popdark}#1}}

% ============================================================
% En-tête / pied
% ============================================================
\pagestyle{fancy}
\fancyhf{}
\renewcommand{\headrulewidth}{0pt}
\renewcommand{\footrulewidth}{0pt}
\lhead{\raisebox{-0.15ex}{\poplogo\fontsize{13}{13}\selectfont\color{popink}OnBuch\color{poporange}+}}
\rhead{\poppill{\DOCMATIERE\ \textbullet\ \DOCNIVEAU\ \textbullet\ Leçon \DOCLECON}{white}{popink}}
\lfoot{\popxbold\fontsize{7.6}{7.6}\selectfont\color{popmuted}\MakeUppercase{Cours signé OnBuch+}\ \textbullet\ {\popsemi\DOCTITRE}}
\rfoot{\tikz[baseline=-0.6ex]\node[rounded corners=8.5pt,fill=popink,minimum height=17pt,minimum width=17pt,inner xsep=5pt,inner sep=0]{\popxbold\fontsize{8}{8}\selectfont\color{popcream}\thepage\,/\,\pageref*{LastPage}};}

% ============================================================
% Titres
% ============================================================
\newcommand{\chaptitre}[2]{%
  \begin{tcolorbox}[enhanced,colback=poporange,colframe=popink,boxrule=2.6pt,arc=11pt,
      drop shadow=popink,left=15pt,right=15pt,top=12pt,bottom=11pt,before skip=3pt,after skip=18pt]
    \poppill{#2}{popink}{popcream}\par\vspace{9pt}
    {\raggedright\hyphenpenalty=10000\exhyphenpenalty=10000\popdisplay\fontsize{22}{24}\selectfont\color{popcream}\MakeUppercase{#1}\par}
  \end{tcolorbox}
}
% Section numérotée : pastille orange avec le numéro + titre Archivo
\newcounter{popsec}
\newcommand{\coursec}[1]{%
  \stepcounter{popsec}\setcounter{popsub}{0}%
  \par\vspace{16pt}\noindent
  \begin{minipage}{\linewidth}
  \tikz[baseline=-0.7ex]\node[draw=popink,line width=1.6pt,fill=poporange,rounded corners=4pt,minimum size=22pt,inner sep=0]{\popdisplay\fontsize{12}{12}\selectfont\color{popcream}\thepopsec};%
  \hspace{8pt}{\popdisplay\fontsize{15}{17}\selectfont\color{popink}\MakeUppercase{#1}}\par
  \vspace{4pt}\noindent{\color{poporange}\rule{36pt}{3.6pt}}
  \end{minipage}\par\nopagebreak\vspace{8pt}
}
\newcounter{popsub}
\newcommand{\courssub}[1]{%
  \stepcounter{popsub}%
  \par\vspace{9pt}\noindent{\popxbold\fontsize{11.5}{13}\selectfont\color{popink}{\color{poporange}\thepopsec.\thepopsub}\hspace{6pt}#1}\par\nopagebreak\vspace{4pt}
}

% ============================================================
% Boîtes pédagogiques — même recette : fond blanc, bordure épaisse colorée,
% ombre dure encre, badge plein. Titre optionnel [#1].
% ============================================================
\tcbset{popbase/.style={breakable,enhanced,colback=white,boxrule=2.3pt,arc=9pt,
  shadow={3pt}{-3pt}{0pt}{popink},left=14pt,right=14pt,top=11pt,bottom=11pt,before skip=11pt,after skip=15pt,
  fontupper=\popsemi\fontsize{10.6}{14.5}\selectfont\color{popink},
  fontlower=\popsemi\fontsize{10.6}{14.5}\selectfont\color{popink}}}
% Coche dessinée (le glyphe ✓ n'existe pas dans les polices du document)
\newcommand{\popcheck}[1]{\tikz[baseline=-0.5ex]\draw[#1,line width=1.6pt,line cap=round,line join=round](-0.13,0.0)--(-0.04,-0.09)--(0.13,0.1);}
\newcommand{\popboxhead}[3]{\poppill{#1}{#2}{white}\ifx\relax#3\relax\else\hspace{6pt}{\popxbold\fontsize{10.6}{12}\selectfont\color{popink}#3}\fi\par\vspace{7pt}}

\newtcolorbox{aretenir}[1][]{popbase,colframe=popgreen,before upper={\popboxhead{À retenir}{popgreen}{#1}}}
\newtcolorbox{exemplebox}[1][]{popbase,colframe=poporange,before upper={\popboxhead{Exemple}{poporange}{#1}}}
\newtcolorbox{definition}[1][]{popbase,colframe=popblue,before upper={\popboxhead{Définition}{popblue}{#1}}}
\newtcolorbox{propriete}[1][]{popbase,colframe=poppurple,before upper={\popboxhead{Propriété}{poppurple}{#1}}}
\newtcolorbox{methode}[1][]{popbase,colframe=popgold,before upper={\popboxhead{Méthode}{popgold}{#1}}}
\newtcolorbox{attention}[1][]{popbase,colframe=poppink,before upper={\popboxhead{Attention}{poppink}{#1}}}
\newtcolorbox{experience}[1][]{popbase,colframe=popblue,colback=popblueL,before upper={\popboxhead{Expérience}{popblue}{#1}}}
% « Le savais-tu ? » : carte violette pleine (contexte, culture, Cameroun)
\newtcolorbox{savaistu}[1][]{popbase,colback=poppurple,colframe=popink,
  fontupper=\popsemi\fontsize{10.4}{14.2}\selectfont\color{white},
  before upper={\poppill{Le savais-tu\,?}{white}{poppurple}\ifx\relax#1\relax\else\hspace{6pt}{\popxbold\color{white}#1}\fi\par\vspace{7pt}}}
% Exercice résolu : énoncé en haut, \tcblower puis la solution sur fond crème
\newcounter{popexo}\newcounter{popexr}
\newtcolorbox{exoresolu}[1][]{popbase,colframe=poporange,colbacklower=popcream,
  segmentation style={popink,line width=1.2pt,dashed},
  before upper={\stepcounter{popexr}\popboxhead{Exercice résolu \thepopexr}{poporange}{#1}},
  before lower={\poppill{Solution}{popink}{popcream}\par\vspace{6pt}}}
% Exercice (non résolu en place) — la correction est dans la partie Corrigés
\newtcolorbox{exercice}[1][]{popbase,colframe=popink,
  before upper={\stepcounter{popexo}\popboxhead{Exercice \thepopexo}{popink}{#1}}}
% Corrigé
\newtcolorbox{corrige}[1][]{popbase,colframe=popgreen,colback=popgreenL,
  before upper={\popboxhead{Corrigé}{popgreen}{#1}}}
% Objectifs / prérequis / bilan
\newtcolorbox{objectifs}{popbase,colframe=popink,colback=poporangeL,
  before upper={\popboxhead{Objectifs de la leçon}{popink}{}}}
\newtcolorbox{prerequis}{popbase,colframe=popink,colback=popblueL,
  before upper={\popboxhead{Prérequis}{popblue}{}}}
\newtcolorbox{bilan}{popbase,colframe=popink,colback=white,boxrule=2.8pt,
  before upper={\popboxhead{Fiche bilan}{poporange}{L'essentiel à connaître pour l'examen}}}
% Figure encadrée avec légende
\newtcolorbox{popfigure}[1][]{enhanced,breakable=false,colback=white,colframe=popline,boxrule=1.4pt,arc=8pt,
  left=10pt,right=10pt,top=10pt,bottom=8pt,before skip=10pt,after skip=12pt,halign=center,
  lower separated=false,halign lower=center,
  fontlower=\popsemi\fontsize{9}{11.5}\selectfont\color{popmuted},
  #1}
\newcommand{\legende}[1]{\par\vspace{6pt}{\popsemi\fontsize{9}{11.5}\selectfont\color{popmuted}#1}}

% ============================================================
% Signature OnBuch+
% ============================================================
\newcommand{\poplogoinline}[1]{{\poplogo\fontsize{#1}{#1}\selectfont\color{popink}OnBuch\color{poporange}+}}
\newcommand{\signatureonbuch}{%
  \begin{tcolorbox}[enhanced,colback=popink,colframe=popink,boxrule=2.2pt,arc=10pt,
      drop shadow=poporange,left=14pt,right=14pt,top=10pt,bottom=10pt,before skip=0pt,after skip=0pt]
    \tikz[baseline=-0.7ex]\node[circle,fill=popgreen,draw=popcream,line width=1.3pt,minimum size=22pt,inner sep=0]{\popcheck{white}};%
    \hspace{9pt}\begin{minipage}[c]{0.62\linewidth}
      {\popxbold\fontsize{12}{13}\selectfont\color{popcream}Signé\ \textbullet\ {\poplogo OnBuch\color{poporange}+}}\par\vspace{2pt}
      {\raggedright\popsemi\fontsize{8}{10.5}\selectfont\color{popline}\MakeUppercase{Équipe pédagogique OnBuch+}\\\MakeUppercase{Conforme au programme officiel MINESEC}\par}
    \end{minipage}\hfill
    {\poplogo\fontsize{22}{22}\selectfont\color{popcream}OnBuch\color{poporange}+}
  \end{tcolorbox}%
}

% ============================================================
% Page de garde
% #1 titre · #2 matière · #3 niveau · #4 module · #5 n° leçon · #6 durée ·
% #7 description / objectifs (texte ou itemize)
% ============================================================
\newcommand{\pagegarde}[7]{%
  \thispagestyle{empty}
  \newgeometry{a4paper,margin=1.7cm,top=1.6cm,bottom=1.4cm}%
  \begin{tikzpicture}[remember picture,overlay]
    \fill[poporange!18] ([xshift=-3.2cm,yshift=-3.6cm]current page.north east) circle (4.6cm);
    \fill[poppurple!14] ([xshift=2.4cm,yshift=7.6cm]current page.south west) circle (3.8cm);
  \end{tikzpicture}%
  {\poplogo\fontsize{30}{30}\selectfont\color{popink}OnBuch\color{poporange}+}\hfill
  \raisebox{0.6ex}{\poppill{Cours signé \textbullet\ Premium}{popink}{popcream}}\par
  \vspace{1pt}\popsquiggle{poppurple}\par
  \vspace{8pt}
  \begin{tcolorbox}[enhanced,colback=poporange,colframe=popink,boxrule=2.8pt,arc=13pt,
      drop shadow=popink,left=18pt,right=18pt,top=12pt,bottom=13pt,after skip=10pt]
    \poppill{#2 \textbullet\ #3}{popink}{popcream}\hspace{5pt}\poppill{#4}{white}{popink}\par\vspace{8pt}
    {\popdisplay\fontsize{14}{14}\selectfont\color{popink}LEÇON #5}\par\vspace{4pt}
    {\raggedright\hyphenpenalty=10000\exhyphenpenalty=10000\popdisplay\fontsize{25}{27}\selectfont\color{popcream}\MakeUppercase{#1}\par}
    \vspace{8pt}
    {\popxbold\fontsize{9.5}{9.5}\selectfont\color{popink}\MakeUppercase{Durée conseillée : #6}}
  \end{tcolorbox}
  \begin{tcolorbox}[enhanced,colback=white,colframe=popink,boxrule=2.2pt,arc=10pt,
      drop shadow=popink,left=16pt,right=16pt,top=10pt,bottom=10pt,after skip=8pt]
    \poppill{Dans cette leçon}{poppurple}{white}\par\vspace{5pt}
    \popsemi\fontsize{9.3}{12.6}\selectfont\color{popink}#7
  \end{tcolorbox}
  \noindent\begin{minipage}[t]{0.487\linewidth}
    \begin{tcolorbox}[enhanced,colback=popgreen,colframe=popink,boxrule=2.2pt,arc=10pt,
        drop shadow=popink,left=11pt,right=11pt,top=10pt,bottom=10pt,height=3.1cm,valign=center]
      \tcbox[colback=white,colframe=popink,boxrule=1.3pt,arc=5pt,boxsep=0pt,nobeforeafter,
        left=3pt,right=3pt,top=3pt,bottom=3pt,valign=center]{\qrcode[height=1.9cm]{https://play.google.com/store/apps/details?id=com.luvvix.onbuch&hl=fr}}%
      \hspace{8pt}\begin{minipage}[c]{0.5\linewidth}
        {\popdisplay\fontsize{11}{12.5}\selectfont\color{white}\MakeUppercase{Télécharge OnBuch}}\par\vspace{4pt}
        {\raggedright\popsemi\fontsize{8.4}{11}\selectfont\color{white}Gratuit \textbullet\ Google Play\\Tuteur IA, annales, concours, résultats\par}
      \end{minipage}
    \end{tcolorbox}
  \end{minipage}\hfill
  \begin{minipage}[t]{0.487\linewidth}
    \begin{tcolorbox}[enhanced,colback=poppurple,colframe=popink,boxrule=2.2pt,arc=10pt,
        drop shadow=popink,left=11pt,right=11pt,top=10pt,bottom=10pt,height=3.1cm,valign=center]
      \tikz[baseline=-0.7ex]\node[circle,fill=white,draw=popink,line width=1.3pt,minimum size=1.6cm,inner sep=0]{\popdisplay\fontsize{24}{24}\selectfont\color{poppurple}+};%
      \hspace{8pt}\begin{minipage}[c]{0.6\linewidth}
        {\popdisplay\fontsize{11}{12.5}\selectfont\color{white}\MakeUppercase{Passe à OnBuch+}}\par\vspace{4pt}
        {\raggedright\popsemi\fontsize{8.4}{11}\selectfont\color{white}Tous les cours signés, Tuteur IA illimité, fiches PDF — onglet OnBuch+ de l'app\par}
      \end{minipage}
    \end{tcolorbox}
  \end{minipage}
  \vspace{8pt}
  \popcontenu
  \vfill
  \signatureonbuch
  \restoregeometry
  \newpage
}

% Rangée « Ce cours contient » (page de garde)
\newcommand{\poptile}[3]{% #1 couleur #2 gros texte #3 légende
  \begin{tcolorbox}[enhanced,colback=#1,colframe=popink,boxrule=2pt,arc=9pt,shadow={2.5pt}{-2.5pt}{0pt}{popink},
      left=6pt,right=6pt,top=9pt,bottom=9pt,halign=center,valign=center,height=1.75cm,width=\linewidth,nobeforeafter]
    {\popdisplay\fontsize{12.5}{13}\selectfont\color{white}\MakeUppercase{#2}}\par\vspace{3pt}
    {\centering\popsemi\fontsize{7.8}{9.5}\selectfont\color{white}#3\par}
  \end{tcolorbox}}
\newcommand{\popcontenu}{%
  \par\noindent{\popxboldls\fontsize{7.5}{7.5}\selectfont\color{popmuted}CE COURS SIGNÉ CONTIENT}\par\vspace{7pt}
  \noindent\begin{minipage}[t]{0.235\linewidth}\poptile{popblue}{Cours}{complet, illustré, pas à pas}\end{minipage}\hfill
  \begin{minipage}[t]{0.235\linewidth}\poptile{poporange}{Méthodes}{et exercices résolus}\end{minipage}\hfill
  \begin{minipage}[t]{0.235\linewidth}\poptile{poppink}{Exercices}{type examen + corrigés}\end{minipage}\hfill
  \begin{minipage}[t]{0.235\linewidth}\poptile{popgreen}{Fiche bilan}{l'essentiel à retenir}\end{minipage}\par}

% Sommaire
\newtcolorbox{sommaire}{enhanced,colback=white,colframe=popink,boxrule=2.2pt,arc=10pt,
  drop shadow=popink,left=18pt,right=18pt,top=13pt,bottom=13pt,before skip=6pt,after skip=20pt,
  before upper={\poppill{Sommaire}{poppurple}{white}\par\vspace{9pt}\popsemi\fontsize{10.6}{14.5}\selectfont\color{popink}}}

% Signature de fin de cours — poussée tout en bas de la dernière page
\newcommand{\findecours}{%
  \par\vfill\nopagebreak
  \begin{center}\popsquiggle{poporange}\end{center}\vspace{4pt}
  \signatureonbuch
}
\AtBeginDocument{\def\llbracket{\mathopen{[\mkern-3.5mu[}}\def\rrbracket{\mathclose{]\mkern-3.5mu]}}} % intervalles d'entiers (glyphe ⟦ absent de la police)
% Glyphes absents de Fira Math : redéfinis à partir de glyphes disponibles
\AtBeginDocument{%
  \def\setminus{\mathbin{\backslash}}\let\smallsetminus\setminus
  \def\vdots{\mathord{\vbox{\baselineskip3.2pt\lineskiplimit0pt\kern4pt\hbox{.}\hbox{.}\hbox{.}}}}%
  \def\ddots{\mathinner{\mkern1mu\raise6.5pt\hbox{.}\mkern2mu\raise3.6pt\hbox{.}\mkern2mu\raise0.7pt\hbox{.}\mkern1mu}}%
  \def\triangle{\mathord{\scalebox{0.9}{$\symup{\Delta}$}}}%
  \def\nabla{\mathord{\raisebox{\depth}{\scalebox{1}[-1]{$\symup{\Delta}$}}}}%
  \def\checkmark{\popcheck{popdark}}%
  \def\star{\mathbin{\ast}}\def\bigstar{\mathord{\ast}}%
  \def\diamond{\mathbin{\scalebox{0.6}{$\Diamond$}}}%
  \def\bigcup{\mathop{\mathchoice{\raisebox{-0.3em}{\scalebox{1.7}{$\cup$}}}{\scalebox{1.25}{$\cup$}}{\cup}{\cup}}\displaylimits}%
  \def\bigcap{\mathop{\mathchoice{\raisebox{-0.3em}{\scalebox{1.7}{$\cap$}}}{\scalebox{1.25}{$\cap$}}{\cap}{\cap}}\displaylimits}%
  \def\lesseqqgtr{\mathrel{\lessgtr}}\def\bigoplus{\mathop{\oplus}\displaylimits}%
}
\providecommand{\ssi}{si et seulement si} % abréviation fréquente

%% FIGFIX-BEGIN v5 — correctifs globaux des figures (docs/onbuch-generation/tools/figfix)
\usepackage{adjustbox}\usepackage{etoolbox}\tcbuselibrary{raster}
% toute figure TikZ/pgfplots plus large que la ligne est réduite à la largeur disponible
\BeforeBeginEnvironment{tikzpicture}{\begin{adjustbox}{max width=\linewidth}}
\AfterEndEnvironment{tikzpicture}{\end{adjustbox}}
% légendes pgfplots lisibles par-dessus les courbes
\pgfplotsset{every axis/.append style={legend style={fill=white,fill opacity=0.92,text opacity=1,draw=black!15,inner sep=3pt}}}
% étiquettes d'arêtes (syntaxe edge["..."]) avec fond blanc
\tikzset{every edge quotes/.append style={fill=white,inner sep=1.2pt}}
%% FIGFIX-END
\begin{document}

\pagegarde{Probabilités, variables aléatoires et loi binomiale}{Mathématiques}{Tle D}{Organisation et gestion des données}{11}{7 h}{Cette leçon introduit le langage et les outils du calcul des probabilités : événements, équiprobabilité, formule de la réunion, puis les variables aléatoires (loi, espérance, variance) et enfin le schéma de Bernoulli et la loi binomiale. Elle fournit les méthodes essentielles pour modéliser et résoudre les problèmes de hasard rencontrés au Bac et dans la vie courante.
\vspace{4pt}
\begin{multicols}{2}\small
\begin{itemize}
\item À la fin de cette leçon, je sais décrire une expérience aléatoire par son univers et traduire une situation par des événements
\item À la fin de cette leçon, je sais calculer une probabilité dans un cas d'équiprobabilité par dénombrement
\item À la fin de cette leçon, je sais utiliser P(Ā) = 1 - P(A) et P(A ∪ B) = P(A) + P(B) - P(A ∩ B)
\item À la fin de cette leçon, je sais construire un arbre de probabilités pour modéliser une expérience à plusieurs épreuves
\item À la fin de cette leçon, je sais déterminer la loi de probabilité d'une variable aléatoire et la présenter dans un tableau
\item À la fin de cette leçon, je sais calculer l'espérance, la variance et l'écart-type d'une variable aléatoire et interpréter l'espérance
\end{itemize}\end{multicols}}

\begin{sommaire}
\begin{multicols}{2}
\begin{enumerate}
\item Vocabulaire des probabilités et équiprobabilité
\item Opérations sur les événements et arbres de probabilités
\item Variables aléatoires et loi de probabilité
\item Espérance, variance et écart-type
\item Épreuve de Bernoulli et schéma de Bernoulli
\item Loi binomiale et applications
\item Probabilités conditionnelles et probabilités totales
\item Activité d'intégration
\item Méthodes et exercices résolus
\item Exercices
\item Corrigés des exercices
\item Fiche bilan
\end{enumerate}
\end{multicols}
\end{sommaire}

\begin{objectifs}
\begin{itemize}
\item À la fin de cette leçon, je sais décrire une expérience aléatoire par son univers et traduire une situation par des événements
\item À la fin de cette leçon, je sais calculer une probabilité dans un cas d'équiprobabilité par dénombrement
\item À la fin de cette leçon, je sais utiliser P(Ā) = 1 - P(A) et P(A ∪ B) = P(A) + P(B) - P(A ∩ B)
\item À la fin de cette leçon, je sais construire un arbre de probabilités pour modéliser une expérience à plusieurs épreuves
\item À la fin de cette leçon, je sais déterminer la loi de probabilité d'une variable aléatoire et la présenter dans un tableau
\item À la fin de cette leçon, je sais calculer l'espérance, la variance et l'écart-type d'une variable aléatoire et interpréter l'espérance
\item À la fin de cette leçon, je sais reconnaître un schéma de Bernoulli et justifier l'utilisation de la loi binomiale
\item À la fin de cette leçon, je sais calculer P(X = k), E(X) et V(X) pour une loi binomiale B(n ; p)
\item À la fin de cette leçon, je sais résoudre un problème concret (jeu, contrôle qualité, sondage) en mobilisant ces outils
\end{itemize}
\end{objectifs}

\begin{prerequis}
\begin{itemize}
\item Dénombrement : arrangements, combinaisons, coefficient binomial C(n,k) (leçon sur l'analyse combinatoire)
\item Calculs sur les fractions et les puissances
\item Notion d'ensemble : réunion, intersection, complémentaire
\item Lecture et construction de tableaux et de diagrammes en barres
\item Pourcentages et proportionnalité
\end{itemize}
\end{prerequis}

\newpage

\coursec{Vocabulaire des probabilités et équiprobabilité}

\textbf{Situation d'accroche.} À Yaoundé, un opérateur de téléphonie mobile lance une grande tombola promotionnelle : chaque client qui recharge au moins \SI{1000}{FCFA} reçoit un code unique et a, annonce la publicité, « une chance sur 100 » de gagner un smartphone. Un élève de Terminale, client fidèle, se pose deux questions très concrètes : s'il recharge 5 fois dans le mois, quelle est sa probabilité de gagner \emph{au moins une fois} ? Et de son côté, l'opérateur, qui attend \num{100000} participants, doit budgéter les lots : combien de smartphones doit-il prévoir \emph{en moyenne} ?

Ces deux questions semblent simples, mais aucune ne peut être résolue avec les seuls outils du collège. La première exige de savoir combiner des probabilités (et fera appel, dans les sections suivantes, à la loi binomiale) ; la seconde repose sur la notion d'\emph{espérance}. Toute cette leçon est construite pour y répondre rigoureusement. Mais avant de calculer, il faut apprendre à \emph{parler} le langage du hasard : c'est l'objet de cette première section.

\textbf{Problématique.} Comment modéliser mathématiquement une situation dépendant du hasard, et comment calculer la probabilité d'un événement lorsque toutes les issues ont la même chance de se réaliser ?

\courssub{Expérience aléatoire et univers}

Certaines expériences sont \emph{déterministes} : si tu laisses tomber une pierre, elle tombe, toujours. D'autres ne le sont pas : quand on lance un dé bien équilibré, on connaît la liste des résultats possibles, mais on ne peut pas prévoir lequel va sortir. C'est cette seconde catégorie qui nous intéresse.

\begin{definition}[Expérience aléatoire]
Une \cle{expérience aléatoire} est une expérience dont le résultat dépend du hasard : on connaît l'ensemble des résultats possibles, mais on ne peut pas prévoir avec certitude celui qui sera obtenu.
\end{definition}

\begin{exemplebox}[Expériences aléatoires de la vie courante]
\begin{itemize}
\item Lancer un dé à six faces et noter le numéro obtenu.
\item Tirer une carte au hasard dans un jeu de 32 cartes.
\item Tirer au sort un lot dans une tombola de quartier à Douala.
\item Interroger un élève au hasard dans une classe de Terminale C et noter son sexe.
\end{itemize}
\end{exemplebox}

Chaque résultat possible d'une expérience aléatoire s'appelle une \cle{issue} (ou \emph{éventualité}). L'ensemble de toutes les issues joue un rôle central : c'est lui qui porte tous les calculs.

\begin{definition}[Univers]
L'\cle{univers} d'une expérience aléatoire, noté $\Omega$ (lire « oméga »), est l'ensemble de toutes les issues possibles de cette expérience.
\end{definition}

\begin{exemplebox}[Trois univers à connaître par cœur]
\begin{itemize}
\item Lancer d'un dé : $\Omega = \{1\,;\,2\,;\,3\,;\,4\,;\,5\,;\,6\}$, donc $\mathrm{Card}(\Omega) = 6$.
\item Lancer d'une pièce : $\Omega = \{\text{Pile}\,;\,\text{Face}\}$.
\item Tirage d'une carte dans un jeu de 32 : $\Omega = \{\text{32 cartes}\}$, $\mathrm{Card}(\Omega) = 32$.
\end{itemize}
\end{exemplebox}

\begin{attention}[Bien choisir son univers]
Le choix de l'univers dépend de ce qu'on observe. Si on lance \emph{deux} dés et qu'on s'intéresse aux deux faces obtenues, l'univers est l'ensemble des 36 couples $(i\,;\,j)$ avec $i, j \in \{1, \dots, 6\}$, et \emph{non pas} l'ensemble des sommes $\{2, 3, \dots, 11\}$. Nous verrons pourquoi ce choix est décisif : les 36 couples sont équiprobables, les 11 sommes ne le sont pas.
\end{attention}

\courssub{Événements : le vocabulaire fondamental}

Un événement, c'est une \emph{question} que l'on pose sur le résultat de l'expérience : « le numéro obtenu est-il pair ? », « la carte tirée est-elle un cœur ? ». Mathématiquement, une telle question se traduit par un sous-ensemble de $\Omega$.

\begin{definition}[Événement, événement élémentaire]
Soit $\Omega$ l'univers d'une expérience aléatoire.
\begin{itemize}
\item Un \cle{événement} est une partie (un sous-ensemble) de $\Omega$.
\item Un \cle{événement élémentaire} est un événement réduit à une seule issue, c'est-à-dire un singleton $\{\omega\}$ avec $\omega \in \Omega$.
\item L'événement $\Omega$ tout entier est appelé \cle{événement certain} : il se réalise toujours.
\item L'ensemble vide $\varnothing$ est appelé \cle{événement impossible} : il ne se réalise jamais.
\end{itemize}
\end{definition}

\begin{exemplebox}[Avec le lancer d'un dé]
Soit $\Omega = \{1\,;\,2\,;\,3\,;\,4\,;\,5\,;\,6\}$.
\begin{itemize}
\item $A = $ « obtenir un nombre pair » $= \{2\,;\,4\,;\,6\}$ est un événement (non élémentaire).
\item $B = $ « obtenir 5 » $= \{5\}$ est un événement élémentaire.
\item « Obtenir un nombre entre 1 et 6 » est l'événement certain $\Omega$.
\item « Obtenir 7 » est l'événement impossible $\varnothing$.
\end{itemize}
\end{exemplebox}

Les événements se combinent grâce aux opérations ensemblistes, que tu as étudiées en Seconde. La correspondance entre le langage courant et le langage des ensembles est essentielle :

\begin{center}
\begin{tabularx}{\linewidth}{|l|X|}
\hline
\rowcolor{popblueL} \textbf{Langage courant} & \textbf{Écriture ensembliste} \\
\hline
« $A$ \textbf{ou} $B$ se réalise » (au moins l'un des deux) & $A \cup B$ (réunion) \\
\hline
« $A$ \textbf{et} $B$ se réalisent tous les deux » & $A \cap B$ (intersection) \\
\hline
« $A$ ne se réalise pas » & $\bar{A}$ (événement contraire) \\
\hline
\end{tabularx}
\end{center}

\begin{definition}[Événement contraire, événements incompatibles]
Soit $A$ un événement d'un univers $\Omega$.
\begin{itemize}
\item L'\cle{événement contraire} de $A$, noté $\bar{A}$, est l'ensemble des issues de $\Omega$ qui n'appartiennent pas à $A$ : $\bar{A} = \Omega \setminus A$.
\item Deux événements $A$ et $B$ sont dits \cle{incompatibles} (ou disjoints) s'ils ne peuvent pas se réaliser en même temps, c'est-à-dire si $A \cap B = \varnothing$.
\end{itemize}
\end{definition}

\begin{attention}[Contraire $\neq$ incompatible]
Ne confonds pas ces deux notions ! $A$ et $\bar{A}$ sont toujours incompatibles, mais deux événements incompatibles ne sont pas forcément contraires l'un de l'autre. Avec le dé, $A = \{1\}$ et $B = \{2\}$ sont incompatibles, pourtant $\bar{A} = \{2\,;\,3\,;\,4\,;\,5\,;\,6\} \neq B$.
\end{attention}

La figure suivante résume visuellement toutes ces notions.

\begin{popfigure}
\begin{center}
\begin{tikzpicture}[scale=1.05]
% Univers
\draw[thick, popdark, fill=popcream] (0,0) rectangle (8.6,5.4);
\node[popdark, font=\bfseries] at (0.75,5.0) {$\Omega$};
% A
\draw[thick, popblue, fill=popblue!25] (2.9,2.5) circle (1.7);
\node[popblue, font=\bfseries] at (2.0,3.6) {$A$};
% B
\draw[thick, poporange, fill=poporange!25] (5.1,2.5) circle (1.7);
\node[poporange, font=\bfseries] at (6.0,3.6) {$B$};
% A inter B
\begin{scope}
\clip (2.9,2.5) circle (1.7);
\fill[poppurple!55] (5.1,2.5) circle (1.7);
\end{scope}
\node[white, font=\footnotesize\bfseries] at (4.0,2.5) {$A\cap B$};
% bar A
\node[popgreen, font=\footnotesize\bfseries] at (7.6,0.55) {$\bar{A}$};
\draw[popgreen, thick, ->] (7.3,0.75) -- (6.9,1.4);
% issues
\fill[popdark] (2.2,2.2) circle (1.5pt) node[below left, font=\footnotesize] {$\omega_1$};
\fill[popdark] (7.4,4.4) circle (1.5pt) node[above, font=\footnotesize] {$\omega_2 \in \bar{A}$};
\end{tikzpicture}
\end{center}
\legende{Diagramme de Venn : l'univers $\Omega$, deux événements $A$ et $B$, leur intersection $A \cap B$ (zone violette) et l'événement contraire $\bar{A}$ (tout ce qui est hors du disque $A$).}
\end{popfigure}

\courssub{Probabilité d'un événement}

Revenons à la tombola de Yaoundé. Dire qu'un client a « une chance sur 100 » de gagner, c'est associer à l'événement « gagner » le nombre $\frac{1}{100} = 0{,}01$. Cette idée — mesurer la vraisemblance d'un événement par un nombre — est le cœur de la leçon.

\begin{definition}[Probabilité d'un événement]
Soit $\Omega = \{\omega_1\,;\,\omega_2\,;\,\dots\,;\,\omega_n\}$ un univers fini. Définir une \cle{probabilité} sur $\Omega$, c'est associer à chaque issue $\omega_i$ un nombre $p_i \in [0\,;\,1]$, appelé probabilité de l'événement élémentaire $\{\omega_i\}$, tel que
\[ p_1 + p_2 + \dots + p_n = 1. \]
La probabilité d'un événement $A$, notée $P(A)$, est alors la \emph{somme} des probabilités des issues qui composent $A$.
\end{definition}

\begin{propriete}[Conséquences immédiates]
Pour tout événement $A$ d'un univers fini $\Omega$ :
\begin{itemize}
\item $0 \le P(A) \le 1$ ;
\item $P(\Omega) = 1$ (l'événement certain a une probabilité de 1) ;
\item $P(\varnothing) = 0$ (l'événement impossible a une probabilité de 0).
\end{itemize}
\end{propriete}

\begin{experience}[Pourquoi la somme vaut-elle 1 ?]
Lance un dé 600 fois et note la fréquence de chaque face. Tu constateras que chaque fréquence se rapproche de $\frac{1}{6}$, et que la somme des six fréquences vaut exactement 1 à chaque série de lancers, puisque chaque lancer produit \emph{une et une seule} face. La probabilité est le modèle théorique de ces fréquences : la contrainte $p_1 + \dots + p_n = 1$ traduit simplement le fait qu'à chaque expérience, \emph{une} issue se réalise obligatoirement.
\end{experience}

\courssub{Le cas de l'équiprobabilité}

La situation la plus simple — et la plus fréquente au Baccalauréat — est celle où toutes les issues ont la même chance d'apparaître : dé non truqué, pièce équilibrée, tirage au sort honnête.

\begin{definition}[Équiprobabilité]
On dit qu'il y a \cle{équiprobabilité} lorsque toutes les issues de $\Omega$ ont la même probabilité. Si $\mathrm{Card}(\Omega) = n$, chaque événement élémentaire a alors pour probabilité $\dfrac{1}{n}$.
\end{definition}

\begin{propriete}[Formule des cas favorables]
En situation d'équiprobabilité, pour tout événement $A$ :
\[ P(A) = \frac{\mathrm{Card}(A)}{\mathrm{Card}(\Omega)} = \frac{\text{nombre de cas favorables}}{\text{nombre de cas possibles}}. \]
\end{propriete}

\begin{experience}[Démonstration (formatrice)]
Notons $n = \mathrm{Card}(\Omega)$. Par équiprobabilité, chaque issue a pour probabilité $\frac{1}{n}$. Par définition, $P(A)$ est la somme des probabilités des issues de $A$. Or $A$ contient $\mathrm{Card}(A)$ issues, chacune de probabilité $\frac{1}{n}$. Donc
\[ P(A) = \underbrace{\frac{1}{n} + \frac{1}{n} + \dots + \frac{1}{n}}_{\mathrm{Card}(A) \text{ termes}} = \frac{\mathrm{Card}(A)}{n} = \frac{\mathrm{Card}(A)}{\mathrm{Card}(\Omega)}. \qquad \blacksquare \]
\end{experience}

\begin{propriete}[Probabilité de l'événement contraire]
Pour tout événement $A$ : $\; P(\bar{A}) = 1 - P(A)$.
\end{propriete}

\begin{experience}[Démonstration exigible]
En situation d'équiprobabilité avec $\mathrm{Card}(\Omega) = n$ : comme $A$ et $\bar{A}$ forment une partition de $\Omega$, on a $\mathrm{Card}(A) + \mathrm{Card}(\bar{A}) = n$, donc $\mathrm{Card}(\bar{A}) = n - \mathrm{Card}(A)$. Alors
\[ P(\bar{A}) = \frac{n - \mathrm{Card}(A)}{n} = 1 - \frac{\mathrm{Card}(A)}{n} = 1 - P(A). \qquad \blacksquare \]
Cette propriété reste vraie même hors équiprobabilité : elle découle de $P(A) + P(\bar{A}) = P(\Omega) = 1$.
\end{experience}

\begin{methode}[Calculer une probabilité par équiprobabilité]
\begin{enumerate}
\item \textbf{Vérifier l'équiprobabilité} : l'énoncé doit préciser « dé équilibré », « tirage au hasard », « pièce non truquée », etc. Sans cette hypothèse, la formule des cas favorables est interdite.
\item \textbf{Dénombrer l'univers} : calculer $\mathrm{Card}(\Omega)$, éventuellement avec les outils de dénombrement (arrangements, combinaisons).
\item \textbf{Dénombrer l'événement} : calculer $\mathrm{Card}(A)$ avec le même soin.
\item \textbf{Conclure} : $P(A) = \dfrac{\mathrm{Card}(A)}{\mathrm{Card}(\Omega)}$, en simplifiant la fraction et en rédigeant une phrase de conclusion.
\end{enumerate}
\end{methode}

\begin{attention}[Le piège classique : les sommes de deux dés]
On lance deux dés équilibrés et on s'intéresse à la \emph{somme} des faces. Les sommes possibles sont $2, 3, \dots, 12$ : onze valeurs. Erreur fréquente : écrire $P(\text{somme} = 7) = \frac{1}{11}$. C'est \textbf{faux}, car les onze sommes ne sont \emph{pas équiprobables} ! Il n'y a qu'une façon d'obtenir 2 (le couple $(1\,;\,1)$), mais six façons d'obtenir 7. Les issues équiprobables sont les 36 \emph{couples} $(i\,;\,j)$, pas les sommes. Le tableau ci-dessous montre que $P(\text{somme} = 7) = \frac{6}{36} = \frac{1}{6}$.
\end{attention}

\begin{popfigure}
\begin{center}
\renewcommand{\arraystretch}{1.15}
\begin{tabular}{|c|*{6}{>{\centering\arraybackslash}p{0.9cm}|}}
\hline
\rowcolor{popblueL} \textbf{Dé 1 $\backslash$ Dé 2} & \textbf{1} & \textbf{2} & \textbf{3} & \textbf{4} & \textbf{5} & \textbf{6} \\
\hline
\textbf{1} & 2 & 3 & 4 & 5 & 6 & \cellcolor{popgoldL} 7 \\
\hline
\textbf{2} & 3 & 4 & 5 & 6 & \cellcolor{popgoldL} 7 & 8 \\
\hline
\textbf{3} & 4 & 5 & 6 & \cellcolor{popgoldL} 7 & 8 & 9 \\
\hline
\textbf{4} & 5 & 6 & \cellcolor{popgoldL} 7 & 8 & 9 & 10 \\
\hline
\textbf{5} & 6 & \cellcolor{popgoldL} 7 & 8 & 9 & 10 & 11 \\
\hline
\textbf{6} & \cellcolor{popgoldL} 7 & 8 & 9 & 10 & 11 & 12 \\
\hline
\end{tabular}
\end{center}
\legende{Les 36 issues équiprobables du lancer de deux dés, représentées par la somme obtenue. Les six cases dorées correspondent à la somme 7 : $P(\text{somme} = 7) = \frac{6}{36} = \frac{1}{6}$, et non $\frac{1}{11}$.}
\end{popfigure}

\courssub{Dénombrement et probabilités : exemples résolus}

Dès que les ensembles sont grands, on ne compte plus « à la main » : on utilise les combinaisons (tirages simultanés, sans ordre) et les arrangements (tirages ordonnés). Rappel : le nombre de façons de choisir $k$ objets parmi $n$ sans tenir compte de l'ordre est $\dbinom{n}{k} = \dfrac{n!}{k!(n-k)!}$.

\begin{exoresolu}[Tirage d'une carte dans un jeu de 32]
On tire une carte au hasard dans un jeu de 32 cartes (8 cartes dans chacune des 4 couleurs : cœur, carreau, trèfle, pique ; il y a 4 rois). On considère les événements $C$ : « tirer un cœur » et $R$ : « tirer un roi ». Calculer $P(C)$, $P(R)$, $P(C \cap R)$ et $P(\bar{C})$.
\tcblower
\textbf{Solution détaillée.} Le tirage est fait au hasard : il y a équiprobabilité, et $\mathrm{Card}(\Omega) = 32$.
\begin{itemize}
\item Il y a 8 cœurs : $P(C) = \dfrac{8}{32} = \dfrac{1}{4}$.
\item Il y a 4 rois : $P(R) = \dfrac{4}{32} = \dfrac{1}{8}$.
\item Un seul roi de cœur : $P(C \cap R) = \dfrac{1}{32}$.
\item Événement contraire : $P(\bar{C}) = 1 - P(C) = 1 - \dfrac{1}{4} = \dfrac{3}{4}$. (Vérification directe : 24 cartes ne sont pas des cœurs, et $\frac{24}{32} = \frac{3}{4}$.)
\end{itemize}
La probabilité de tirer « un cœur \emph{ou} un roi » sera calculée dans la section 2 avec la formule de la réunion : $P(C \cup R) = P(C) + P(R) - P(C \cap R) = \frac{1}{4} + \frac{1}{8} - \frac{1}{32} = \frac{11}{32}$.
\end{exoresolu}

\begin{exoresolu}[Tombola du lycée]
Pour la fête du lycée de Ngaoundéré, on vend 200 billets de tombola numérotés de 1 à 200. On tire au sort un billet gagnant. Calculer la probabilité que le numéro gagnant soit : (a) un multiple de 10 ; (b) un numéro à deux chiffres.
\tcblower
\textbf{Solution détaillée.} Tirage au sort : équiprobabilité, $\mathrm{Card}(\Omega) = 200$.
\begin{itemize}
\item[(a)] Les multiples de 10 entre 1 et 200 sont $10, 20, \dots, 200$ : il y en a 20. Donc $P(A) = \dfrac{20}{200} = \dfrac{1}{10}$.
\item[(b)] Les numéros à deux chiffres vont de 10 à 99 : il y en a $99 - 10 + 1 = 90$. Donc $P(B) = \dfrac{90}{200} = \dfrac{9}{20}$.
\end{itemize}
\end{exoresolu}

\begin{exoresolu}[Délégation de classe et combinaisons]
Une classe de Terminale D compte 12 filles et 8 garçons. On choisit au hasard une délégation de 3 élèves pour représenter la classe. Quelle est la probabilité que la délégation soit composée exactement de 2 filles et 1 garçon ?
\tcblower
\textbf{Solution détaillée.} Le choix est simultané et sans ordre : on utilise les combinaisons, avec équiprobabilité de tous les groupes de 3 élèves.
\begin{itemize}
\item Nombre de délégations possibles : $\mathrm{Card}(\Omega) = \dbinom{20}{3} = \dfrac{20 \times 19 \times 18}{3 \times 2 \times 1} = \num{1140}$.
\item Cas favorables : choisir 2 filles parmi 12 \emph{et} 1 garçon parmi 8 :
\[ \mathrm{Card}(A) = \binom{12}{2} \times \binom{8}{1} = 66 \times 8 = 528. \]
\item Conclusion : $P(A) = \dfrac{528}{\num{1140}} = \dfrac{44}{95} \approx 0{,}46$.
\end{itemize}
\end{exoresolu}

\begin{attention}[Rédaction : les trois réflexes du candidat]
\begin{itemize}
\item Toujours \textbf{annoncer l'équiprobabilité} (« le dé est équilibré, donc les issues sont équiprobables ») avant d'utiliser la formule des cas favorables.
\item Toujours \textbf{préciser l'univers} et son cardinal dans la rédaction.
\item Une probabilité est un nombre de $[0\,;\,1]$ : si tu trouves $P(A) = 1{,}4$ ou $P(A) = -0{,}2$, il y a forcément une erreur de calcul ou de dénombrement.
\end{itemize}
\end{attention}

\begin{savaistu}[Pourquoi l'organisateur gagne toujours]
Les jeux de hasard — Loto, paris sportifs, tombolas promotionnelles — reposent sur l'équiprobabilité : chaque billet, chaque grille a la même chance de gagner. Pourtant, à long terme, c'est toujours l'organisateur qui s'enrichit. Le secret ? Les gains distribués sont calculés pour être, \emph{en moyenne}, inférieurs aux mises encaissées. Cette « moyenne » est précisément l'\emph{espérance} de la variable aléatoire « gain », que nous étudierons dans la section 4. Pour la tombola de notre accroche : si chacun des \num{100000} clients a une chance sur 100 de gagner, l'opérateur doit prévoir en moyenne $\num{100000} \times \frac{1}{100} = \num{1000}$ smartphones — un coût prévisible, largement couvert par les recharges. Les mathématiques ne prédisent pas \emph{qui} gagnera, mais elles prédisent \emph{combien} gagneront : c'est toute la puissance des probabilités.
\end{savaistu}

\begin{exercice}[Pour s'entraîner]
On lance deux dés équilibrés à six faces et on note la somme des faces obtenues.
\begin{enumerate}
\item Calculer la probabilité d'obtenir une somme égale à 9.
\item Calculer la probabilité d'obtenir une somme strictement supérieure à 9.
\item Calculer la probabilité d'obtenir une somme inférieure ou égale à 9. Que remarque-t-on ?
\end{enumerate}
\end{exercice}

\begin{corrige}[Exercice : sommes de deux dés]
Les 36 couples $(i\,;\,j)$ sont équiprobables, $\mathrm{Card}(\Omega) = 36$.
\begin{enumerate}
\item Somme égale à 9 : couples $(3\,;\,6)$, $(4\,;\,5)$, $(5\,;\,4)$, $(6\,;\,3)$, soit 4 cas favorables : $P = \dfrac{4}{36} = \dfrac{1}{9}$.
\item Somme strictement supérieure à 9, c'est-à-dire 10, 11 ou 12 : 3 cas pour 10, 2 cas pour 11, 1 cas pour 12, soit 6 cas : $P = \dfrac{6}{36} = \dfrac{1}{6}$.
\item L'événement « somme $\le 9$ » est le contraire de « somme $> 9$ » : $P = 1 - \dfrac{1}{6} = \dfrac{5}{6}$. On remarque que le passage au contraire évite de compter 30 cas un par un : c'est souvent la méthode la plus rapide.
\end{enumerate}
\end{corrige}

\begin{aretenir}[L'essentiel de la section]
\begin{itemize}
\item Une \cle{expérience aléatoire} a des issues imprévisibles mais connues ; leur ensemble est l'\cle{univers} $\Omega$.
\item Un \cle{événement} est une partie de $\Omega$ ; $\bar{A}$ est son contraire ; $A$ et $B$ sont incompatibles si $A \cap B = \varnothing$.
\item Une probabilité associe à chaque événement un nombre de $[0\,;\,1]$, avec $P(\Omega) = 1$ et $P(\varnothing) = 0$.
\item En \cle{équiprobabilité} : $P(A) = \dfrac{\mathrm{Card}(A)}{\mathrm{Card}(\Omega)}$ — à condition d'avoir vérifié que les issues sont équiprobables.
\item Toujours utile : $P(\bar{A}) = 1 - P(A)$.
\end{itemize}
\end{aretenir}

\coursec{Opérations sur les événements et arbres de probabilités}

Dans la section précédente, nous avons appris à décrire une expérience aléatoire (univers, événements, événement contraire) et à calculer une probabilité en situation d'\cle{équiprobabilité} grâce au quotient $\dfrac{\text{nombre de cas favorables}}{\text{nombre de cas possibles}}$. Mais dans la vie courante, on s'intéresse rarement à un seul événement isolé : un commerçant de Mokolo veut connaître la probabilité qu'un client achète du riz \emph{ou} de l'huile ; un élève veut la probabilité d'obtenir \emph{au moins} une mention sur deux épreuves ; un technicien d'Orange veut la probabilité que deux relais tombent en panne \emph{en même temps}. Pour répondre à ces questions, il faut savoir \emph{combiner} les événements. C'est l'objet de cette section : réunion, intersection, événements incompatibles, stratégie du contraire, puis l'outil roi du chapitre : l'\cle{arbre de probabilités}.

\courssub{Réunion et intersection d'événements}

\begin{definition}[Réunion et intersection]
Soient $A$ et $B$ deux événements d'un univers $\Omega$.
\begin{itemize}
\item La \cle{réunion} $A \cup B$ (lire « $A$ \textbf{ou} $B$ ») est l'événement réalisé lorsque $A$ est réalisé, \emph{ou} $B$ est réalisé, \emph{ou les deux}. Le « ou » des mathématiques est \textbf{inclusif}.
\item L'\cle{intersection} $A \cap B$ (lire « $A$ \textbf{et} $B$ ») est l'événement réalisé lorsque $A$ et $B$ sont réalisés \emph{simultanément}.
\item $A$ et $B$ sont dits \cle{incompatibles} (ou disjoints) lorsqu'ils ne peuvent pas se réaliser en même temps, c'est-à-dire $A \cap B = \emptyset$.
\end{itemize}
\end{definition}

\begin{exemplebox}[Lecture concrète]
On lance un dé équilibré à six faces : $\Omega = \{1 ; 2 ; 3 ; 4 ; 5 ; 6\}$. Soient $A =$ « obtenir un nombre pair » $= \{2 ; 4 ; 6\}$ et $B =$ « obtenir un multiple de $3$ » $= \{3 ; 6\}$. Alors :
\[ A \cup B = \{2 ; 3 ; 4 ; 6\} \qquad \text{et} \qquad A \cap B = \{6\}. \]
$A$ et $B$ ne sont \textbf{pas} incompatibles, car le résultat $6$ réalise les deux à la fois. En revanche, $A$ et $C =$ « obtenir $5$ » sont incompatibles : $A \cap C = \emptyset$.
\end{exemplebox}

\begin{attention}[Le « ou » n'est pas le « ou exclusif » du langage courant]
Au restaurant, « fromage \emph{ou} dessert » signifie souvent « l'un ou l'autre, mais pas les deux ». En probabilités, $A \cup B$ inclut toujours la possibilité que $A$ et $B$ se réalisent \emph{ensemble}. C'est précisément cette subtilité qui explique la formule du paragraphe suivant.
\end{attention}

\courssub{Probabilité de la réunion de deux événements}

Imaginons que l'on compte les éléments de $A \cup B$ en additionnant naïvement le nombre d'éléments de $A$ et celui de $B$ : les éléments qui appartiennent aux deux ensembles (ceux de $A \cap B$) sont alors comptés \textbf{deux fois}. Pour corriger, il faut retirer une fois ce double comptage. Cette idée de dénombrement donne immédiatement le théorème fondamental.

\begin{propriete}[Théorème — probabilité de la réunion]
Pour tous événements $A$ et $B$ d'un univers $\Omega$ :
\[ P(A \cup B) = P(A) + P(B) - P(A \cap B). \]
En particulier, si $A$ et $B$ sont \textbf{incompatibles} ($A \cap B = \emptyset$) :
\[ P(A \cup B) = P(A) + P(B). \]
\emph{La démonstration en situation d'équiprobabilité est exigible au Baccalauréat.}
\end{propriete}

\begin{experience}[Démonstration guidée du théorème]
On se place en situation d'\textbf{équiprobabilité} sur un univers fini $\Omega$ contenant $n$ issues. On note $\mathrm{Card}(E)$ le nombre d'éléments d'un événement $E$.
\begin{enumerate}
\item \textbf{Compter $A \cup B$.} Quand on additionne $\mathrm{Card}(A) + \mathrm{Card}(B)$, chaque élément de $A \cap B$ est compté deux fois (une fois dans $A$, une fois dans $B$). Pour obtenir le bon total, on retire une fois $\mathrm{Card}(A \cap B)$ :
\[ \mathrm{Card}(A \cup B) = \mathrm{Card}(A) + \mathrm{Card}(B) - \mathrm{Card}(A \cap B). \]
\item \textbf{Diviser par $n$.} Par équiprobabilité, $P(E) = \dfrac{\mathrm{Card}(E)}{n}$. En divisant l'égalité précédente par $n$ :
\[ \frac{\mathrm{Card}(A \cup B)}{n} = \frac{\mathrm{Card}(A)}{n} + \frac{\mathrm{Card}(B)}{n} - \frac{\mathrm{Card}(A \cap B)}{n}. \]
\item \textbf{Conclure.} On reconnaît dans chaque terme une probabilité :
\[ P(A \cup B) = P(A) + P(B) - P(A \cap B). \]
\item \textbf{Cas incompatible.} Si $A \cap B = \emptyset$, alors $P(A \cap B) = 0$ et la formule devient $P(A \cup B) = P(A) + P(B)$.
\end{enumerate}
\end{experience}

\begin{popfigure}
\begin{center}
\begin{tikzpicture}[scale=1.05]
% Disque A
\fill[popblue!35] (0,0) circle (1.5);
% Disque B
\fill[poporange!40] (1.8,0) circle (1.5);
% Intersection re-colorée
\begin{scope}
\clip (0,0) circle (1.5);
\fill[popgreen!45] (1.8,0) circle (1.5);
\end{scope}
\draw[popblue, thick] (0,0) circle (1.5);
\draw[poporange, thick] (1.8,0) circle (1.5);
\node[popblue, font=\bfseries] at (-1,0.9) {$A$};
\node[poporange, font=\bfseries] at (2.8,0.9) {$B$};
\node at (0.9,0.35) {$A \cap B$};
\node at (-0.9,-0.5) {\small compté 1 fois};
\node at (2.7,-0.5) {\small compté 1 fois};
\node at (0.9,-0.55) {\small compté 2 fois};
\node at (0.9,-1.05) {\small $\Rightarrow$ on retire 1 fois};
% Univers
\draw[popdark, rounded corners] (-2.2,-1.9) rectangle (4.1,2);
\node[popdark] at (3.6,1.7) {$\Omega$};
\end{tikzpicture}
\end{center}
\legende{Diagramme de Venn : en additionnant $P(A)$ et $P(B)$, la zone $A \cap B$ est comptée deux fois ; il faut donc soustraire $P(A \cap B)$ une fois.}
\end{popfigure}

\begin{exemplebox}[Application directe]
Dans une classe de Terminale C de 50 élèves, 28 élèves pratiquent le football, 20 pratiquent la course, et 10 pratiquent les deux. On choisit un élève au hasard. Soient $F$ « l'élève joue au football » et $C$ « l'élève fait de la course ».
\[ P(F \cup C) = P(F) + P(C) - P(F \cap C) = \frac{28}{50} + \frac{20}{50} - \frac{10}{50} = \frac{38}{50} = 0{,}76. \]
La probabilité que l'élève pratique au moins l'un des deux sports est $0{,}76$.
\end{exemplebox}

\begin{attention}[Erreur fréquente : le double comptage]
Ne jamais écrire $P(A \cup B) = P(A) + P(B)$ sans avoir \textbf{vérifié} que $A$ et $B$ sont incompatibles ! Dans l'exemple précédent, additionner directement donnerait $\frac{28}{50} + \frac{20}{50} = \frac{48}{50} = 0{,}96$, ce qui est faux : les 10 élèves pratiquant les deux sports auraient été comptés deux fois. Au Bac, justifie toujours : « les événements sont incompatibles, donc... » ou applique la formule complète.
\end{attention}

\courssub{La stratégie du contraire : le réflexe « au moins un »}

L'événement contraire $\bar{A}$ vérifie $P(\bar{A}) = 1 - P(A)$, donc $P(A) = 1 - P(\bar{A})$. Cette égalité, toute simple, devient une \emph{stratégie} redoutable dès que l'énoncé contient les mots « \cle{au moins un} ». En effet, « obtenir au moins un succès » se réalise de plusieurs façons (exactement 1, exactement 2, ...), alors que son contraire « obtenir \textbf{aucun} succès » ne se réalise que d'une seule façon, souvent très facile à calculer.

\begin{methode}[Traiter un événement « au moins un »]
\begin{enumerate}
\item Repérer l'expression « au moins un » (ou « au moins une ») dans la question.
\item Nommer $A$ l'événement cherché et identifier son contraire : $\bar{A} =$ « \textbf{aucun} » (ou « aucune »).
\item Calculer $P(\bar{A})$, généralement plus simple (un seul cas à traiter, souvent un produit de probabilités).
\item Conclure : $P(A) = 1 - P(\bar{A})$.
\end{enumerate}
\end{methode}

\begin{exemplebox}[Intuition rapide]
On lance deux fois une pièce équilibrée. Probabilité d'obtenir \emph{au moins un} Pile ? Le contraire est « aucun Pile », c'est-à-dire Face puis Face, de probabilité $\frac12 \times \frac12 = \frac14$. Donc $P(\text{au moins un Pile}) = 1 - \frac14 = \frac34$. Vérification par dénombrement : sur les 4 issues $\{PP, PF, FP, FF\}$, trois contiennent au moins un Pile.
\end{exemplebox}

\courssub{Arbres de probabilités : construire et exploiter}

Beaucoup d'expériences se déroulent \textbf{en plusieurs étapes} : tirer deux boules successivement, répéter un lancer, subir deux tests médicaux. L'\cle{arbre pondéré} est l'outil qui permet d'organiser ces situations sans rien oublier.

\begin{definition}[Arbre pondéré]
Un arbre pondéré représente les étapes successives d'une expérience : chaque nœud se ramifie en branches correspondant aux issues possibles de l'étape, et chaque branche porte la probabilité de cette issue \emph{sachant le chemin déjà parcouru}.
\end{definition}

\begin{propriete}[Règles de calcul sur un arbre — admises]
\begin{itemize}
\item \textbf{Règle des produits} : la probabilité d'un chemin (une succession d'issues) est le \emph{produit} des probabilités rencontrées le long de ce chemin.
\item \textbf{Règle de la somme} : la probabilité d'un événement réalisé par plusieurs chemins est la \emph{somme} des probabilités de ces chemins.
\item \textbf{Contrôle} : la somme des probabilités des branches issues d'un même nœud vaut toujours $1$.
\end{itemize}
\end{propriete}

\begin{methode}[Construire et exploiter un arbre pondéré]
\begin{enumerate}
\item Identifier les étapes de l'expérience et, à chaque étape, les issues possibles.
\item Calculer la probabilité de chaque branche, en tenant compte de ce qui s'est passé avant (crucial en cas de tirage \textbf{sans remise} : la composition de l'urne change).
\item Vérifier qu'à chaque nœud, la somme des probabilités vaut $1$.
\item Pour un événement donné : repérer les chemins favorables, calculer la probabilité de chacun (produits), puis les additionner.
\item Si la question contient « au moins un », envisager le passage au contraire.
\end{enumerate}
\end{methode}

\courssub{Exemple complet : deux tirages sans remise}

\begin{exoresolu}[Urne : 3 boules rouges, 2 noires — deux tirages sans remise]
Une urne contient 3 boules rouges et 2 boules noires, indiscernables au toucher. On tire une première boule, on ne la remet pas, puis on tire une deuxième boule.
\begin{enumerate}
\item Construire l'arbre pondéré de l'expérience.
\item Calculer la probabilité d'obtenir deux boules rouges.
\item Calculer la probabilité d'obtenir au moins une boule rouge.
\end{enumerate}
\tcblower
\textbf{1. Arbre pondéré.} Au premier tirage, $P(R_1) = \frac35$ et $P(N_1) = \frac25$. Au deuxième tirage, la composition de l'urne dépend du premier résultat (tirage \emph{sans remise}) :
\begin{itemize}
\item si la première boule est rouge, il reste 2 rouges et 2 noires : $P(R_2) = \frac24 = \frac12$ et $P(N_2) = \frac12$ ;
\item si la première boule est noire, il reste 3 rouges et 1 noire : $P(R_2) = \frac34$ et $P(N_2) = \frac14$.
\end{itemize}
L'arbre complet est représenté sur la figure ci-dessous.

\textbf{2. Deux boules rouges.} Un seul chemin convient : $R_1$ puis $R_2$. Par la règle des produits :
\[ P(R_1 \cap R_2) = \frac35 \times \frac12 = \frac{3}{10} = 0{,}3. \]

\textbf{3. Au moins une rouge.} Trois chemins sont favorables ($RR$, $RN$, $NR$) : le calcul direct est possible mais long. Appliquons la \textbf{stratégie du contraire} : le contraire de « au moins une rouge » est « aucune rouge », c'est-à-dire deux noires (chemin $N_1 N_2$) :
\[ P(\bar{A}) = P(N_1 \cap N_2) = \frac25 \times \frac14 = \frac{2}{20} = \frac{1}{10}. \]
Donc :
\[ P(A) = 1 - \frac{1}{10} = \frac{9}{10} = 0{,}9. \]
\emph{Vérification par la règle de la somme} : $P(RR) + P(RN) + P(NR) = \frac{3}{10} + \frac{3}{10} + \frac{3}{10} = \frac{9}{10}$. Les deux méthodes concordent.
\end{exoresolu}

\begin{popfigure}
\begin{center}
\begin{tikzpicture}[
grow=right,
level distance=3.2cm,
level 1/.style={sibling distance=3.4cm},
level 2/.style={sibling distance=1.7cm},
every node/.style={font=\small},
edge from parent/.style={draw=popmuted, thick}
]
\node[popdark, font=\bfseries] {Départ}
child {
node[circle, fill=poporange!25, draw=poporange, inner sep=2pt] (N1) {$N_1$}
child {
node[circle, fill=poporange!25, draw=poporange, inner sep=2pt] (N1N2) {$N_2$}
edge from parent node[above, poporange] {$\frac14$}
}
child {
node[circle, fill=popblue!25, draw=popblue, inner sep=2pt] (N1R2) {$R_2$}
edge from parent node[above, popblue] {$\frac34$}
}
edge from parent node[above, poporange] {$\frac25$}
}
child {
node[circle, fill=popblue!25, draw=popblue, inner sep=2pt] (R1) {$R_1$}
child {
node[circle, fill=poporange!25, draw=poporange, inner sep=2pt] (R1N2) {$N_2$}
edge from parent node[above, poporange] {$\frac12$}
}
child {
node[circle, fill=popblue!25, draw=popblue, inner sep=2pt] (R1R2) {$R_2$}
edge from parent node[above, popblue] {$\frac12$}
}
edge from parent node[above, popblue] {$\frac35$}
};
% Probabilités des chemins, alignées à droite des feuilles
\node[right=0.25cm of R1R2, popdark, anchor=west] {$P(R_1 \cap R_2) = \frac35 \times \frac12 = \frac{3}{10}$};
\node[right=0.25cm of R1N2, popdark, anchor=west] {$P(R_1 \cap N_2) = \frac35 \times \frac12 = \frac{3}{10}$};
\node[right=0.25cm of N1R2, popdark, anchor=west] {$P(N_1 \cap R_2) = \frac25 \times \frac34 = \frac{3}{10}$};
\node[right=0.25cm of N1N2, popdark, anchor=west] {$P(N_1 \cap N_2) = \frac25 \times \frac14 = \frac{1}{10}$};
\end{tikzpicture}
\end{center}
\legende{Arbre pondéré des deux tirages sans remise. Les probabilités du deuxième niveau dépendent du premier tirage : c'est la signature du « sans remise ». La somme des quatre chemins vaut $\frac{3}{10} + \frac{3}{10} + \frac{3}{10} + \frac{1}{10} = 1$.}
\end{popfigure}

\begin{attention}[Sans remise : les probabilités changent d'un niveau à l'autre]
L'erreur classique consiste à recopier $\frac35$ et $\frac25$ au deuxième niveau de l'arbre. Après un tirage \textbf{sans remise}, l'urne ne contient plus que 4 boules, et sa composition dépend de la boule déjà tirée. Prends le réflexe de réécrire, à chaque nœud, le contenu de l'urne avant de calculer les probabilités des branches. À l'inverse, si le tirage est \textbf{avec remise}, l'urne est reconstituée à l'identique : les probabilités se recopient alors telles quelles d'un niveau à l'autre.
\end{attention}

\begin{savaistu}[Les arbres, outils de décision]
Les arbres de probabilités ne servent pas qu'au Bac : en médecine, ils modélisent les tests de dépistage (test positif ou négatif sachant que le patient est malade ou sain) ; dans les télécommunications, des entreprises comme MTN ou Orange les utilisent pour estimer la fiabilité d'un réseau où plusieurs équipements se succèdent. Le mathématicien anglais Thomas Bayes (XVIII\textsuperscript{e} siècle) a bâti sur ces idées une théorie — les probabilités « à l'envers » — qui est aujourd'hui au cœur de l'intelligence artificielle.
\end{savaistu}

\begin{exercice}[À toi de jouer]
Une urne contient 4 boules vertes et 1 boule blanche. On effectue deux tirages successifs sans remise.
\begin{enumerate}
\item Construire l'arbre pondéré de l'expérience.
\item Calculer la probabilité d'obtenir deux boules vertes.
\item Calculer la probabilité d'obtenir au moins une boule blanche.
\end{enumerate}
\end{exercice}

\begin{corrige}[Exercice 1]
\textbf{1.} Premier tirage : $P(V_1) = \frac45$, $P(B_1) = \frac15$. Si $V_1$ est réalisé, il reste 3 vertes et 1 blanche : branches $\frac34$ et $\frac14$. Si $B_1$ est réalisé, il ne reste que 4 vertes : $P(V_2) = 1$ et $P(B_2) = 0$.

\textbf{2.} $P(V_1 \cap V_2) = \frac45 \times \frac34 = \frac{12}{20} = \frac35$.

\textbf{3.} Contraire de « au moins une blanche » : « aucune blanche », c'est-à-dire deux vertes. Donc :
\[ P(\text{au moins une blanche}) = 1 - \frac35 = \frac25. \]
Vérification par les chemins : $P(V_1 \cap B_2) + P(B_1 \cap V_2) = \frac45 \times \frac14 + \frac15 \times 1 = \frac15 + \frac15 = \frac25$.
\end{corrige}

\begin{aretenir}[L'essentiel de la section]
\begin{itemize}
\item $A \cup B$ : « $A$ ou $B$ ou les deux » ; $A \cap B$ : « $A$ et $B$ simultanément ».
\item Pour tous événements : $P(A \cup B) = P(A) + P(B) - P(A \cap B)$ ; si $A$ et $B$ sont incompatibles, $P(A \cup B) = P(A) + P(B)$.
\item « Au moins un » $\Rightarrow$ passer au contraire « aucun » : $P(A) = 1 - P(\bar{A})$.
\item Sur un arbre : produits le long d'un chemin, somme des chemins favorables ; à chaque nœud, la somme des probabilités des branches vaut $1$.
\item Sans remise : les probabilités du niveau suivant dépendent du résultat précédent ; avec remise : elles se recopient à l'identique.
\end{itemize}
\end{aretenir}

\coursec{Variables aléatoires et loi de probabilité}

Dans la section précédente, nous avons appris à décrire les événements d'une expérience aléatoire et à calculer leurs probabilités à l'aide d'arbres pondérés. Mais dans la vie courante, on s'intéresse rarement à un événement en soi : on s'intéresse à un \emph{nombre} qui dépend du hasard. Le transporteur de la ligne Douala--Yaoundé veut connaître le \emph{nombre de passagers} par voyage ; l'opérateur téléphonique compte le \emph{nombre d'appels} reçus par heure ; le joueur de loterie veut savoir \emph{combien il gagne ou perd}. Dans chaque cas, le hasard produit une valeur numérique : c'est exactement l'idée de \cle{variable aléatoire}. Cette section va nous apprendre à décrire complètement un tel phénomène grâce à sa \cle{loi de probabilité} et à sa \cle{fonction de répartition}.

\courssub{Notion de variable aléatoire}

\textbf{Intuition.} Lançons deux dés et intéressons-nous à la somme des faces obtenues. Le résultat du lancer est un couple $(1\,;\,4)$, $(3\,;\,6)$, etc. : ce n'est pas un nombre unique. Mais à chaque couple, on peut \emph{associer} un nombre : la somme des deux faces. Ainsi, au couple $(1\,;\,4)$ on associe $5$, au couple $(3\,;\,6)$ on associe $9$. On a construit une \emph{fonction} qui, à chaque issue de l'expérience, associe un nombre réel. Cette fonction s'appelle une variable aléatoire.

\begin{definition}[Variable aléatoire]
Soit $\Omega$ l'univers (fini) d'une expérience aléatoire muni d'une probabilité $P$. On appelle \cle{variable aléatoire} sur $\Omega$ toute fonction $X$ définie sur $\Omega$ et à valeurs dans $\mathbb{R}$ :
\[ X : \Omega \longrightarrow \mathbb{R}, \qquad \omega \longmapsto X(\omega). \]
L'ensemble des valeurs prises par $X$ est noté $X(\Omega) = \{x_1\,;\,x_2\,;\,\dots\,;\,x_n\}$.
\end{definition}

\begin{attention}[Ce que « variable aléatoire » ne veut pas dire]
Une variable aléatoire n'est \textbf{ni une variable au sens algébrique} (comme le $x$ d'une équation), \textbf{ni un nombre} : c'est une \emph{fonction}. Le mot « aléatoire » vient du fait que la valeur $X(\omega)$ dépend de l'issue $\omega$, laquelle relève du hasard. Avant l'expérience, on ne sait pas quelle valeur $X$ prendra ; après l'expérience, la valeur est parfaitement déterminée.
\end{attention}

Pour chaque valeur $x_i$ prise par $X$, on peut considérer l'événement « $X$ prend la valeur $x_i$ », que l'on note $(X = x_i)$. C'est un événement au sens des sections précédentes :
\[ (X = x_i) = \{\omega \in \Omega \; ; \; X(\omega) = x_i\}. \]
On peut donc calculer sa probabilité $P(X = x_i)$. De même, $(X \leq x)$ désigne l'événement « $X$ prend une valeur inférieure ou égale à $x$ », et $(X > x_i)$ est l'événement contraire de $(X \leq x_i)$.

\begin{exemplebox}[Somme de deux dés]
On lance deux dés équilibrés (un rouge, un bleu). L'univers est l'ensemble des $36$ couples $(i\,;\,j)$ avec $i, j \in \{1\,;\,2\,;\,\dots\,;\,6\}$, toutes les issues étant équiprobables de probabilité $\dfrac{1}{36}$. On définit la variable aléatoire $S$ = somme des deux faces. Alors $S(\Omega) = \{2\,;\,3\,;\,\dots\,;\,12\}$. Par exemple, l'événement $(S = 5)$ est réalisé par les couples $(1\,;\,4)$, $(2\,;\,3)$, $(3\,;\,2)$, $(4\,;\,1)$, donc $P(S = 5) = \dfrac{4}{36} = \dfrac{1}{9}$.
\end{exemplebox}

\courssub{Loi de probabilité d'une variable aléatoire}

\textbf{Intuition.} Connaître une variable aléatoire, c'est savoir répondre à deux questions : \emph{quelles valeurs peut-elle prendre ?} et \emph{avec quelles probabilités ?} La réponse complète à ces deux questions s'appelle la loi de probabilité de $X$.

\begin{definition}[Loi de probabilité]
Soit $X$ une variable aléatoire prenant les valeurs $x_1, x_2, \dots, x_n$. On appelle \cle{loi de probabilité} de $X$ la donnée de tous les couples $\big(x_i\,;\, p_i\big)$ où $p_i = P(X = x_i)$. On la présente sous forme de tableau :
\begin{center}
\begin{tabularx}{\linewidth}{|l|*{4}{>{\centering\arraybackslash}X|}}
\hline
\rowcolor{popblueL} Valeurs $x_i$ & $x_1$ & $x_2$ & $\dots$ & $x_n$ \\
\hline
$p_i = P(X = x_i)$ & $p_1$ & $p_2$ & $\dots$ & $p_n$ \\
\hline
\end{tabularx}
\end{center}
\end{definition}

\begin{propriete}[Condition fondamentale]
Les nombres $p_i$ d'une loi de probabilité vérifient :
\[ \text{pour tout } i,\quad 0 \leq p_i \leq 1 \qquad \text{et} \qquad \sum_{i=1}^{n} p_i = p_1 + p_2 + \dots + p_n = 1. \]
\end{propriete}

\textbf{Justification.} Les événements $(X = x_1), (X = x_2), \dots, (X = x_n)$ sont deux à deux incompatibles (une issue $\omega$ ne peut donner qu'une seule valeur de $X$) et leur réunion est l'univers $\Omega$ tout entier (toute issue donne \emph{une} valeur de $X$). Ils forment donc une partition de $\Omega$, et par conséquent la somme de leurs probabilités vaut $P(\Omega) = 1$. $\blacksquare$

Cette propriété est un outil de \emph{vérification} précieux : après avoir calculé une loi, on doit toujours contrôler que la somme des $p_i$ vaut $1$. Elle sert aussi à retrouver une probabilité manquante : si l'on connaît toutes les $p_i$ sauf une, la dernière s'obtient par différence avec $1$.

\begin{methode}[Déterminer la loi de probabilité d'une variable aléatoire]
\begin{enumerate}
\item \textbf{Lister les valeurs} prises par $X$ : pour chaque issue possible, calculer $X(\omega)$, puis écrire l'ensemble $X(\Omega) = \{x_1\,;\,\dots\,;\,x_n\}$ \emph{sans en oublier} (penser en particulier à la valeur $0$ et aux valeurs négatives pour un gain algébrique).
\item \textbf{Calculer chaque probabilité} $p_i = P(X = x_i)$, en dénombrant les issues favorables (équiprobabilité) ou en utilisant un arbre pondéré.
\item \textbf{Présenter le tableau} de la loi : valeurs $x_i$ rangées dans l'ordre croissant en première ligne, probabilités $p_i$ en seconde ligne.
\item \textbf{Vérifier} que $\sum p_i = 1$. Si ce n'est pas le cas, il y a une erreur : valeur oubliée ou calcul faux.
\end{enumerate}
\end{methode}

\begin{exoresolu}[Loi de la somme de deux dés]
On lance deux dés équilibrés et on note $S$ la somme des faces obtenues. Déterminer la loi de probabilité de $S$.
\tcblower
\textbf{Étape 1 : valeurs de $S$.} La plus petite somme est $1+1 = 2$, la plus grande est $6+6 = 12$, et toutes les valeurs intermédiaires sont atteintes : $S(\Omega) = \{2\,;\,3\,;\,4\,;\,5\,;\,6\,;\,7\,;\,8\,;\,9\,;\,10\,;\,11\,;\,12\}$.

\textbf{Étape 2 : calcul des probabilités.} Il y a $36$ couples équiprobables. On dénombre les couples donnant chaque somme : par exemple $(S=4)$ est réalisé par $(1\,;\,3)$, $(2\,;\,2)$, $(3\,;\,1)$, soit $3$ issues, donc $P(S=4) = \dfrac{3}{36}$. En procédant de même pour chaque valeur :

\textbf{Étape 3 : tableau de la loi.}
\begin{center}
\begin{tabularx}{\linewidth}{|l|*{11}{>{\centering\arraybackslash}X|}}
\hline
\rowcolor{popblueL} $s_i$ & 2 & 3 & 4 & 5 & 6 & 7 & 8 & 9 & 10 & 11 & 12 \\
\hline
$P(S = s_i)$ & $\frac{1}{36}$ & $\frac{2}{36}$ & $\frac{3}{36}$ & $\frac{4}{36}$ & $\frac{5}{36}$ & $\frac{6}{36}$ & $\frac{5}{36}$ & $\frac{4}{36}$ & $\frac{3}{36}$ & $\frac{2}{36}$ & $\frac{1}{36}$ \\
\hline
\end{tabularx}
\end{center}

\textbf{Étape 4 : vérification.}
\[ \frac{1+2+3+4+5+6+5+4+3+2+1}{36} = \frac{36}{36} = 1. \]
La loi est correcte. On remarque la symétrie de la loi autour de la valeur $7$, la plus probable.
\end{exoresolu}

\courssub{Représentation graphique : le diagramme en barres}

Une loi de probabilité se visualise par un \cle{diagramme en barres} : sur l'axe des abscisses, on place les valeurs $x_i$ ; au-dessus de chacune, on trace une barre dont la hauteur est la probabilité $p_i$. Ce diagramme donne d'un coup d'œil les valeurs les plus probables et la forme générale de la loi (symétrie, étalement).

\begin{popfigure}
\begin{center}
\begin{tikzpicture}
\begin{axis}[
    width=0.85\linewidth, height=6cm,
    ybar, bar width=8pt,
    xlabel={Valeur $s_i$ de la somme},
    ylabel={$P(S = s_i)$},
    xmin=1, xmax=13,
    ymin=0, ymax=0.18,
    xtick={2,3,4,5,6,7,8,9,10,11,12},
    ytick={0,0.05,0.10,0.15},
    yticklabels={$0$, $0{,}05$, $0{,}10$, $0{,}15$},
    axis lines=left,
    tick label style={font=\small},
    label style={font=\small},
]
\addplot[fill=popblue, draw=popdark] coordinates {
    (2,0.0278) (3,0.0556) (4,0.0833) (5,0.1111) (6,0.1389) (7,0.1667)
    (8,0.1389) (9,0.1111) (10,0.0833) (11,0.0556) (12,0.0278)
};
\end{axis}
\end{tikzpicture}
\end{center}
\legende{Diagramme en barres de la loi de la somme $S$ de deux dés : la loi est symétrique autour de $7$, valeur la plus probable.}
\end{popfigure}

\courssub{Exemple de la vie courante : loi d'un gain algébrique}

Les variables aléatoires sont l'outil naturel pour analyser les jeux d'argent, très présents dans la vie quotidienne (tombola de quartier, jeux de paris). L'exemple suivant est un modèle à retenir : il sera repris dans la section suivante pour le calcul de l'espérance.

\begin{exemplebox}[Un jeu de dé à 100 FCFA]
À la fête de la jeunesse, un forain propose le jeu suivant : le joueur mise $100$ FCFA, puis lance un dé équilibré.
\begin{itemize}
\item Si le dé donne $6$, le forain rend $500$ FCFA au joueur ;
\item si le dé donne $1$, le forain rend $200$ FCFA ;
\item dans les autres cas ($2$, $3$, $4$ ou $5$), le forain ne rend rien.
\end{itemize}
On note $X$ le \cle{gain algébrique} du joueur, c'est-à-dire ce qu'il reçoit \emph{moins} sa mise de $100$ FCFA. Déterminons la loi de $X$.

\textbf{Valeurs de $X$.} Si le dé donne $6$ : $X = 500 - 100 = 400$. Si le dé donne $1$ : $X = 200 - 100 = 100$. Sinon : $X = 0 - 100 = -100$. Donc $X(\Omega) = \{-100\,;\,100\,;\,400\}$.

\textbf{Probabilités.} Le dé étant équilibré : $P(X = 400) = \dfrac{1}{6}$, $P(X = 100) = \dfrac{1}{6}$, et $P(X = -100) = \dfrac{4}{6} = \dfrac{2}{3}$.

\textbf{Tableau de la loi} (valeurs rangées dans l'ordre croissant) :
\begin{center}
\begin{tabularx}{0.7\linewidth}{|l|*{3}{>{\centering\arraybackslash}X|}}
\hline
\rowcolor{popblueL} $x_i$ & $-100$ & $100$ & $400$ \\
\hline
$P(X = x_i)$ & $\dfrac{2}{3}$ & $\dfrac{1}{6}$ & $\dfrac{1}{6}$ \\
\hline
\end{tabularx}
\end{center}

\textbf{Vérification :} $\dfrac{2}{3} + \dfrac{1}{6} + \dfrac{1}{6} = \dfrac{4+1+1}{6} = 1$. La loi est bien définie. On observe que le joueur perd sa mise dans deux tiers des cas : nous verrons à la section suivante, grâce à l'espérance, si ce jeu est équitable.
\end{exemplebox}

\begin{attention}[Les deux erreurs classiques]
\begin{itemize}
\item \textbf{Oublier une valeur possible de $X$.} C'est l'erreur la plus fréquente : on oublie souvent la valeur $0$ (par exemple « nombre de succès » quand il n'y a aucun succès) ou une valeur négative (la perte dans un gain algébrique). Avant de calculer, demande-toi toujours : \emph{quelle est la plus petite valeur possible ? la plus grande ? toutes les valeurs intermédiaires sont-elles atteintes ?}
\item \textbf{Ne pas vérifier que $\sum p_i = 1$.} Cette vérification, qui prend dix secondes, détecte la quasi-totalité des erreurs. Au Baccalauréat, écris explicitement la somme et sa valeur $1$ : c'est un réflexe qui rassure le correcteur et qui te protège.
\end{itemize}
\end{attention}

\courssub{Fonction de répartition}

\textbf{Intuition.} La loi de probabilité répond à la question « quelle est la probabilité que $X$ prenne \emph{exactement} la valeur $x_i$ ? ». Mais dans la pratique, on pose souvent une autre question : « quelle est la probabilité que $X$ ne dépasse pas une certaine valeur ? ». Un fournisseur d'accès veut connaître la probabilité que le nombre de pannes mensuelles soit \emph{au plus} égal à $3$ ; un assureur, la probabilité que le nombre de sinistres soit \emph{au plus} égal à un seuil. C'est la fonction de répartition qui répond à ces questions, en \emph{accumulant} les probabilités.

\begin{definition}[Fonction de répartition]
Soit $X$ une variable aléatoire. On appelle \cle{fonction de répartition} de $X$ la fonction $F$ définie sur $\mathbb{R}$ par :
\[ F(x) = P(X \leq x). \]
\end{definition}

Concrètement, si $X$ prend les valeurs $x_1 < x_2 < \dots < x_n$, alors $F(x)$ est la somme des probabilités des valeurs inférieures ou égales à $x$ :
\begin{itemize}
\item si $x < x_1$ : aucune valeur de $X$ n'est inférieure ou égale à $x$, donc $F(x) = 0$ ;
\item si $x_1 \leq x < x_2$ : $F(x) = p_1$ ;
\item si $x_2 \leq x < x_3$ : $F(x) = p_1 + p_2$ ;
\item \dots
\item si $x \geq x_n$ : toutes les valeurs sont inférieures ou égales à $x$, donc $F(x) = 1$.
\end{itemize}

\begin{propriete}[Propriétés de la fonction de répartition]
La fonction de répartition $F$ d'une variable aléatoire $X$ (prenant un nombre fini de valeurs) est :
\begin{enumerate}
\item une fonction \cle{croissante} sur $\mathbb{R}$ ;
\item une fonction \cle{en escalier} : constante sur chaque intervalle $[x_i\,;\,x_{i+1}[$, avec un « saut » en chaque valeur $x_i$ ;
\item telle que $\displaystyle\lim_{x \to -\infty} F(x) = 0$ et $\displaystyle\lim_{x \to +\infty} F(x) = 1$ (elle vaut $0$ avant la plus petite valeur de $X$ et $1$ à partir de la plus grande).
\end{enumerate}
\end{propriete}

\textbf{Justification de la croissance.} Si $a \leq b$, alors l'événement $(X \leq a)$ est inclus dans l'événement $(X \leq b)$ : toute issue réalisant le premier réalise le second. Or, si $A \subset B$, alors $P(A) \leq P(B)$. Donc $F(a) = P(X \leq a) \leq P(X \leq b) = F(b)$. $\blacksquare$

\begin{propriete}[Lien entre fonction de répartition et loi de probabilité]
Si les valeurs de $X$ sont rangées dans l'ordre strictement croissant $x_1 < x_2 < \dots < x_n$, alors le saut de $F$ en $x_i$ est exactement la probabilité de $x_i$ :
\[ P(X = x_i) = F(x_i) - F(x_{i-1}) \qquad (\text{avec } F(x_0) = 0). \]
Plus généralement, pour tous réels $a < b$ : $P(a < X \leq b) = F(b) - F(a)$.
\end{propriete}

\textbf{Justification.} L'événement $(X \leq x_i)$ est la réunion des deux événements incompatibles $(X \leq x_{i-1})$ et $(X = x_i)$, puisqu'il n'y a aucune valeur de $X$ strictement entre $x_{i-1}$ et $x_i$. Donc $F(x_i) = F(x_{i-1}) + P(X = x_i)$, d'où le résultat. $\blacksquare$

Ce lien est précieux dans les deux sens : on construit $F$ en additionnant les $p_i$, et on retrouve la loi en calculant les sauts de $F$.

\begin{exoresolu}[Fonction de répartition du jeu de dé]
On reprend le jeu de la fête de la jeunesse : $X$ prend les valeurs $-100$, $100$, $400$ avec les probabilités $\dfrac{2}{3}$, $\dfrac{1}{6}$, $\dfrac{1}{6}$. Déterminer la fonction de répartition $F$ de $X$, puis calculer $P(X \leq 250)$.
\tcblower
On accumule les probabilités dans l'ordre croissant des valeurs :
\[ F(x) = \begin{cases}
0 & \text{si } x < -100 \\[6pt]
\dfrac{2}{3} & \text{si } -100 \leq x < 100 \\[6pt]
\dfrac{2}{3} + \dfrac{1}{6} = \dfrac{5}{6} & \text{si } 100 \leq x < 400 \\[6pt]
\dfrac{5}{6} + \dfrac{1}{6} = 1 & \text{si } x \geq 400
\end{cases} \]
Pour $P(X \leq 250)$ : comme $100 \leq 250 < 400$, on lit directement $F(250) = \dfrac{5}{6}$. Autrement dit, la probabilité que le joueur gagne \emph{au plus} $250$ FCFA vaut $\dfrac{5}{6}$. On vérifie le lien avec la loi : $P(X = 400) = F(400) - F(100) = 1 - \dfrac{5}{6} = \dfrac{1}{6}$, ce qui redonne bien la probabilité du tableau.
\end{exoresolu}

\begin{popfigure}
\begin{center}
\begin{tikzpicture}[scale=1.0, >=Stealth]
% Axes
\draw[->, popdark] (-5.2,0) -- (5.6,0) node[below left] {$x$};
\draw[->, popdark] (0,-0.4) -- (0,3.4) node[below left] {$F(x)$};
% Graduations y
\draw[popdark] (-0.08,2) -- (0.08,2) node[left=2pt, xshift=-6pt] {$\frac{2}{3}$};
\draw[popdark] (-0.08,2.5) -- (0.08,2.5) node[left=2pt, xshift=-6pt] {$\frac{5}{6}$};
\draw[popdark] (-0.08,3) -- (0.08,3) node[left=2pt, xshift=-6pt] {$1$};
% Graduations x
\draw[popdark] (-3,0.08) -- (-3,-0.08) node[below] {$-100$};
\draw[popdark] (1,0.08) -- (1,-0.08) node[below] {$100$};
\draw[popdark] (4,0.08) -- (4,-0.08) node[below] {$400$};
% Pointillés de rappel
\draw[dashed, popmuted] (-3,0) -- (-3,2);
\draw[dashed, popmuted] (1,0) -- (1,2.5);
\draw[dashed, popmuted] (4,0) -- (4,3);
\draw[dashed, popmuted] (0,2) -- (-3,2);
\draw[dashed, popmuted] (0,2.5) -- (1,2.5);
\draw[dashed, popmuted] (0,3) -- (4,3);
% Paliers
\draw[very thick, popblue] (-5,0) -- (-3,0);
\draw[very thick, popblue] (-3,2) -- (1,2);
\draw[very thick, popblue] (1,2.5) -- (4,2.5);
\draw[very thick, popblue] (4,3) -- (5.4,3);
% Points fermés et ouverts
\fill[popblue] (-3,2) circle (2pt);
\draw[popblue, fill=white, thick] (-3,0) circle (2pt);
\fill[popblue] (1,2.5) circle (2pt);
\draw[popblue, fill=white, thick] (1,2) circle (2pt);
\fill[popblue] (4,3) circle (2pt);
\draw[popblue, fill=white, thick] (4,2.5) circle (2pt);
\end{tikzpicture}
\end{center}
\legende{Fonction de répartition du gain $X$ : fonction en escalier, croissante, nulle avant $-100$, égale à $1$ à partir de $400$. Les sauts ont pour hauteurs $\frac{2}{3}$, $\frac{1}{6}$ et $\frac{1}{6}$, c'est-à-dire exactement les probabilités $P(X = x_i)$.}
\end{popfigure}

\begin{methode}[Construire et utiliser une fonction de répartition]
\begin{enumerate}
\item Ranger les valeurs de $X$ dans l'ordre croissant : $x_1 < x_2 < \dots < x_n$.
\item Calculer les \emph{probabilités cumulées} : $F(x_1) = p_1$, $F(x_2) = p_1 + p_2$, \dots, $F(x_n) = 1$.
\item Écrire $F$ par morceaux : $0$ avant $x_1$, puis chaque cumul sur l'intervalle $[x_i\,;\,x_{i+1}[$, et $1$ à partir de $x_n$.
\item Pour calculer $P(a < X \leq b)$, utiliser $F(b) - F(a)$ ; pour retrouver $P(X = x_i)$, calculer le saut $F(x_i) - F(x_{i-1})$.
\end{enumerate}
\end{methode}

\begin{attention}[Inégalités strictes ou larges : sois vigilant]
Par définition, $F(x) = P(X \leq x)$ avec une inégalité \emph{large} : la valeur $x$ est \textbf{incluse}. Ainsi $F(100) = \dfrac{5}{6}$ dans l'exemple du jeu (la valeur $100$ est comptée), alors que pour tout $x$ strictement inférieur à $100$, $F(x) = \dfrac{2}{3}$. Sur le graphique, c'est le point \emph{fermé} en haut du saut qui indique la valeur de $F$ en $x_i$. De même, $P(X < x_i) = F(x_{i-1})$ et non $F(x_i)$ : ne confonds pas $(X < x_i)$ et $(X \leq x_i)$.
\end{attention}

\begin{aretenir}[L'essentiel de la section]
\begin{itemize}
\item Une \cle{variable aléatoire} $X$ est une \emph{fonction} de $\Omega$ dans $\mathbb{R}$ ; $(X = x_i)$ est un événement dont on calcule la probabilité.
\item La \cle{loi de probabilité} de $X$ est le tableau des couples $\big(x_i\,;\, P(X = x_i)\big)$ ; on a toujours $\sum p_i = 1$ : cette vérification est obligatoire.
\item La loi se représente par un \cle{diagramme en barres}.
\item La \cle{fonction de répartition} $F(x) = P(X \leq x)$ est croissante, en escalier, nulle avant la plus petite valeur et égale à $1$ à partir de la plus grande.
\item On passe de $F$ à la loi par les sauts : $P(X = x_i) = F(x_i) - F(x_{i-1})$.
\end{itemize}
\end{aretenir}

\begin{savaistu}[Des probabilités aux assurances]
La fonction de répartition est l'outil de base des actuaires, ces mathématiciens qui calculent les primes d'assurance. Pour fixer le prix d'une assurance automobile à Douala, l'assureur étudie la variable aléatoire $X$ = « montant annuel des indemnisations d'un client » et utilise $F(x) = P(X \leq x)$ pour estimer, par exemple, la probabilité que les indemnisations dépassent les primes encaissées. La théorie que tu découvres ici, formalisée par le mathématicien soviétique Andreï Kolmogorov en 1933, fait ainsi tourner l'assurance, la banque et les télécommunications — y compris celles des réseaux MTN et Orange au Cameroun, où les lois de probabilités servent à dimensionner le nombre d'antennes relais.
\end{savaistu}

\begin{exercice}[À toi de jouer]
Une urne contient $3$ boules rouges et $2$ boules vertes, indiscernables au toucher. On tire simultanément $2$ boules et on note $X$ le nombre de boules rouges obtenues.
\begin{enumerate}
\item Déterminer la loi de probabilité de $X$ (on vérifiera que la somme des probabilités vaut $1$).
\item Représenter cette loi par un diagramme en barres.
\item Déterminer la fonction de répartition $F$ de $X$ et en déduire $P(X \leq 1)$.
\end{enumerate}
\end{exercice}

\begin{corrige}[Exercice]
\begin{enumerate}
\item Il y a $\binom{5}{2} = 10$ tirages équiprobables. $X$ prend les valeurs $0$, $1$, $2$ (on ne peut tirer que $0$, $1$ ou $2$ boules rouges).
\[ P(X = 0) = \frac{\binom{2}{2}}{10} = \frac{1}{10}, \quad P(X = 1) = \frac{\binom{3}{1}\binom{2}{1}}{10} = \frac{6}{10}, \quad P(X = 2) = \frac{\binom{3}{2}}{10} = \frac{3}{10}. \]
Vérification : $\dfrac{1}{10} + \dfrac{6}{10} + \dfrac{3}{10} = 1$.
\item Diagramme en barres : barres de hauteurs $0{,}1$ ; $0{,}6$ ; $0{,}3$ au-dessus des valeurs $0$, $1$, $2$.
\item $F(x) = 0$ si $x < 0$ ; $F(x) = \dfrac{1}{10}$ si $0 \leq x < 1$ ; $F(x) = \dfrac{7}{10}$ si $1 \leq x < 2$ ; $F(x) = 1$ si $x \geq 2$. Comme $1 \leq 1 < 2$, on a $P(X \leq 1) = F(1) = \dfrac{7}{10}$.
\end{enumerate}
\end{corrige}

Nous savons désormais décrire complètement une variable aléatoire par sa loi et sa fonction de répartition. Mais ces tableaux sont parfois longs à manipuler : ne pourrait-on résumer une loi par quelques nombres significatifs, comme une « valeur moyenne » ou une mesure de sa dispersion ? C'est précisément l'objet de la section suivante : l'\emph{espérance}, la \emph{variance} et l'\emph{écart-type}.

\coursec{Espérance, variance et écart-type}

Dans la section précédente, nous avons appris à décrire complètement une variable aléatoire par sa \cle{loi de probabilité} : un tableau donnant, pour chaque valeur $x_i$, la probabilité $p_i = P(X = x_i)$. Cette description est exhaustive, mais elle est parfois trop riche : face à un jeu de hasard, un commerçant ou un assureur souhaite souvent résumer la situation par un ou deux nombres. \emph{Combien puis-je espérer gagner en moyenne ? Les gains sont-ils stables ou très dispersés ?} Ce sont précisément les questions auxquelles répondent l'\cle{espérance} et la \cle{variance} (ou son équivalent, l'\cle{écart-type}), que nous étudions dans cette section.

\courssub{Espérance mathématique}

\textbf{Une intuition pour commencer.}\par
Imaginons une tombola organisée par une coopérative scolaire à Yaoundé. Chaque billet coûte $500$ FCFA et permet de gagner, selon le tirage, $0$, $1\,000$ ou $5\,000$ FCFA. Si tu joues une seule fois, ton gain est imprévisible. Mais si tu joues un très grand nombre de fois, ton gain \emph{moyen par partie} va se stabiliser autour d'une valeur précise : c'est cette valeur que l'on appelle \cle{espérance mathématique}. Elle se calcule en faisant la moyenne des gains possibles, pondérée par leurs probabilités.

\begin{definition}[Espérance d'une variable aléatoire]
Soit $X$ une variable aléatoire prenant les valeurs $x_1, x_2, \dots, x_n$ avec les probabilités $p_1, p_2, \dots, p_n$ (où $p_i = P(X = x_i)$). L'\cle{espérance mathématique} de $X$ est le nombre réel noté $E(X)$ défini par :
\[ E(X) = \sum_{i=1}^{n} p_i\, x_i = p_1 x_1 + p_2 x_2 + \cdots + p_n x_n. \]
\end{definition}

\begin{aretenir}[Interprétation concrète de l'espérance]
L'espérance $E(X)$ s'interprète comme la \textbf{valeur moyenne} de $X$ sur un \textbf{grand nombre} de répétitions de l'expérience. C'est l'idée intuitive de la \emph{loi des grands nombres} : si l'on répète l'expérience $N$ fois et que l'on note les valeurs obtenues, la moyenne arithmétique de ces valeurs se rapproche de $E(X)$ lorsque $N$ devient très grand. L'espérance n'est donc \emph{pas} une valeur que $X$ prend nécessairement : c'est une moyenne théorique à long terme.
\end{aretenir}

\begin{attention}[L'espérance n'est pas une valeur possible]
L'espérance peut très bien ne figurer \textbf{pas} parmi les valeurs $x_i$. Par exemple, le nombre moyen d'enfants par ménage peut valoir $3,4$, alors qu'aucun ménage n'a $3,4$ enfants ! Ne dis jamais « $X$ vaut en moyenne exactement $E(X)$ à chaque partie », mais « sur un grand nombre de parties, le gain moyen par partie se rapproche de $E(X)$ ».
\end{attention}

\textbf{Jeu équitable, favorable, défavorable.}\par
Lorsque $X$ représente le \emph{gain algébrique} d'un jeu (gain net, c'est-à-dire en tenant compte de la mise), l'espérance permet de classer le jeu :

\begin{definition}[Jeu équitable, favorable, défavorable]
Soit $X$ le gain algébrique d'un jeu de hasard.
\begin{itemize}
\item Si $E(X) = 0$, le jeu est dit \cle{équitable} : sur un grand nombre de parties, ni le joueur ni l'organisateur ne gagnent en moyenne.
\item Si $E(X) > 0$, le jeu est \cle{favorable} au joueur.
\item Si $E(X) < 0$, le jeu est \cle{défavorable} au joueur (c'est le cas de tous les jeux d'argent organisés commercialement : l'organisateur se réserve une marge).
\end{itemize}
\end{definition}

\begin{exemplebox}[La tombola de la coopérative]
Une tombola propose $200$ billets à $500$ FCFA chacun. Les lots sont : un lot de $20\,000$ FCFA, deux lots de $5\,000$ FCFA, dix lots de $1\,000$ FCFA. On note $X$ le gain \emph{algébrique} (lot reçu moins la mise) d'un joueur achetant un billet. La loi de $X$ est :
\begin{center}
\begin{tabularx}{\linewidth}{|l|X|X|X|X|}
\hline
\rowcolor{popblueL} $x_i$ & $19\,500$ & $4\,500$ & $500$ & $-500$ \\
\hline
$p_i$ & $\dfrac{1}{200}$ & $\dfrac{2}{200}$ & $\dfrac{10}{200}$ & $\dfrac{187}{200}$ \\
\hline
\end{tabularx}
\end{center}
Calculons l'espérance :
\begin{align*}
E(X) &= \frac{1}{200}\times 19\,500 + \frac{2}{200}\times 4\,500 + \frac{10}{200}\times 500 + \frac{187}{200}\times(-500)\\
&= \frac{19\,500 + 9\,000 + 5\,000 - 93\,500}{200} = \frac{-60\,000}{200} = -300.
\end{align*}
Ainsi $E(X) = -300$ FCFA : le jeu est \textbf{défavorable} au joueur. En moyenne, chaque joueur « perd » $300$ FCFA par billet, ce qui correspond à la marge de la coopérative (c'est le but d'une tombola : financer les activités !).
\end{exemplebox}

\courssub{Linéarité de l'espérance}

Très souvent, le gain réel s'obtient à partir d'une variable $X$ par une transformation affine : par exemple, si la mise est de $b$ francs et que le lot est proportionnel à $X$, le gain net est $Y = aX + b$. La propriété suivante, admise, évite de recalculer toute la loi.

\begin{propriete}[Linéarité de l'espérance (admise)]
Soit $X$ une variable aléatoire et $a$, $b$ deux nombres réels. Alors :
\[ E(aX + b) = a\,E(X) + b. \]
En particulier, $E(X + b) = E(X) + b$ et $E(aX) = a\,E(X)$.
\end{propriete}

\begin{exemplebox}[Vérification sur un exemple]
Soit $X$ de loi : $P(X=1) = 0{,}5$, $P(X=2) = 0{,}3$, $P(X=3) = 0{,}2$. Alors $E(X) = 0{,}5\times 1 + 0{,}3\times 2 + 0{,}2\times 3 = 1{,}7$. Posons $Y = 2X + 3$. La loi de $Y$ porte sur les valeurs $5$, $7$, $9$ avec les mêmes probabilités, donc :
\[ E(Y) = 0{,}5\times 5 + 0{,}3\times 7 + 0{,}2\times 9 = 2{,}5 + 2{,}1 + 1{,}8 = 6{,}4. \]
Et l'on vérifie : $2E(X) + 3 = 2\times 1{,}7 + 3 = 6{,}4$. La formule fonctionne.
\end{exemplebox}

\courssub{Variance et écart-type : mesurer la dispersion}

\textbf{Pourquoi l'espérance ne suffit pas.}\par
Considérons deux jeux. Jeu A : tu gagnes $990$ FCFA ou $1\,010$ FCFA avec probabilité $\tfrac12$ chacun. Jeu B : tu gagnes $0$ FCFA ou $2\,000$ FCFA avec probabilité $\tfrac12$ chacun. Les deux jeux ont la même espérance $E = 1\,000$ FCFA. Pourtant, ils ne se ressemblent pas : le jeu A est très stable, le jeu B est très risqué. Il nous faut un outil pour mesurer cet \emph{écart} par rapport à la moyenne : c'est la \cle{variance}.

\begin{definition}[Variance et écart-type]
Soit $X$ une variable aléatoire d'espérance $m = E(X)$, prenant les valeurs $x_1, \dots, x_n$ avec les probabilités $p_1, \dots, p_n$.
La \cle{variance} de $X$ est :
\[ V(X) = \sum_{i=1}^{n} p_i\,\bigl(x_i - m\bigr)^2. \]
L'\cle{écart-type} de $X$ est : $\sigma(X) = \sqrt{V(X)}$.
\end{definition}

\begin{aretenir}[Sens de la variance et de l'écart-type]
La variance est la \textbf{moyenne des carrés des écarts à la moyenne}. Plus les valeurs de $X$ sont dispersées autour de $E(X)$, plus $V(X)$ et $\sigma(X)$ sont grands. On a toujours $V(X) \geq 0$, et $V(X) = 0$ si et seulement si $X$ est constante. L'écart-type a l'avantage de s'exprimer dans la \emph{même unité} que $X$ (des FCFA si $X$ est en FCFA), ce qui le rend plus parlant que la variance.
\end{aretenir}

La définition de la variance est rarement la méthode de calcul la plus rapide. Le théorème suivant, exigible au Baccalauréat, fournit une formule beaucoup plus pratique.

\begin{propriete}[Théorème de König-Huygens]
Pour toute variable aléatoire $X$ :
\[ V(X) = E(X^2) - \bigl[E(X)\bigr]^2, \]
c'est-à-dire $V(X) = \displaystyle\sum_{i=1}^{n} p_i x_i^2 - \left(\sum_{i=1}^{n} p_i x_i\right)^2$.
\end{propriete}

\begin{experience}[Démonstration guidée du théorème de König-Huygens]
Posons $m = E(X)$ et développons la définition de la variance.
\begin{enumerate}
\item \textbf{Développons le carré.} Pour chaque indice $i$ : $(x_i - m)^2 = x_i^2 - 2m\,x_i + m^2$.
\item \textbf{Injectons dans la somme.}
\[ V(X) = \sum_{i=1}^{n} p_i (x_i - m)^2 = \sum_{i=1}^{n} p_i x_i^2 - 2m\sum_{i=1}^{n} p_i x_i + m^2 \sum_{i=1}^{n} p_i. \]
\item \textbf{Reconnaissons les sommes.} Par définition, $\sum p_i x_i^2 = E(X^2)$ et $\sum p_i x_i = m$. De plus, la somme des probabilités vaut $1$ : $\sum p_i = 1$.
\item \textbf{Concluons.}
\[ V(X) = E(X^2) - 2m \times m + m^2 \times 1 = E(X^2) - 2m^2 + m^2 = E(X^2) - m^2. \]
D'où $V(X) = E(X^2) - [E(X)]^2$. $\square$
\end{enumerate}
\end{experience}

\begin{attention}[Les deux erreurs classiques]
\begin{itemize}
\item \textbf{Ne confonds pas} $E(X^2)$ et $[E(X)]^2$ ! Le premier est la moyenne des carrés ($\sum p_i x_i^2$), le second est le carré de la moyenne. Ils sont presque toujours différents, et leur différence est précisément la variance, qui est \emph{positive ou nulle}. Si tu trouves une variance négative, tu as fait une erreur de calcul.
\item \textbf{L'écart-type est la racine carrée de la variance}, pas la variance elle-même. Quand l'énoncé demande $\sigma(X)$, n'oublie pas la dernière étape : $\sigma(X) = \sqrt{V(X)}$.
\end{itemize}
\end{attention}

\begin{methode}[Calculer $E(X)$, $V(X)$, $\sigma(X)$ avec un tableau]
\begin{enumerate}
\item Construis un tableau à quatre lignes : $x_i$ ; $p_i$ ; $p_i x_i$ ; $p_i x_i^2$.
\item Vérifie que $\sum p_i = 1$ (contrôle rapide de la loi).
\item Additionne la troisième ligne : $E(X) = \sum p_i x_i$.
\item Additionne la quatrième ligne : $E(X^2) = \sum p_i x_i^2$.
\item Applique König-Huygens : $V(X) = E(X^2) - [E(X)]^2$.
\item Termine par $\sigma(X) = \sqrt{V(X)}$, en arrondissant proprement.
\end{enumerate}
\end{methode}

\begin{exoresolu}[Espérance et écart-type du gain d'une tombola]
\textbf{Énoncé.} Une association de quartier à Douala organise une tombola de $100$ billets vendus $1\,000$ FCFA chacun. Les lots sont : un vélo d'une valeur de $40\,000$ FCFA, quatre sacs de riz de $5\,000$ FCFA chacun. On note $X$ le gain algébrique (lot moins la mise) d'un acheteur d'un billet. Déterminer la loi de $X$, calculer $E(X)$, $V(X)$ et $\sigma(X)$. Le jeu est-il équitable ?
\tcblower
\textbf{Solution.} Les valeurs possibles de $X$ sont : $40\,000 - 1\,000 = 39\,000$ ; $5\,000 - 1\,000 = 4\,000$ ; et $-1\,000$ (billet perdant, $95$ cas sur $100$).
\begin{center}
\begin{tabularx}{\linewidth}{|l|X|X|X|X|}
\hline
\rowcolor{popblueL} $x_i$ & $39\,000$ & $4\,000$ & $-1\,000$ & \textbf{Totaux} \\
\hline
$p_i$ & $0{,}01$ & $0{,}04$ & $0{,}95$ & $1$ \\
\hline
$p_i x_i$ & $390$ & $160$ & $-950$ & $E(X) = -400$ \\
\hline
$p_i x_i^2$ & $15\,210\,000$ & $640\,000$ & $950\,000$ & $E(X^2) = 16\,800\,000$ \\
\hline
\end{tabularx}
\end{center}
\textbf{Espérance :} $E(X) = 390 + 160 - 950 = -400$ FCFA. Le jeu est \textbf{défavorable} au joueur : il n'est pas équitable (pour être équitable, il faudrait $E(X) = 0$). Chaque billet rapporte en moyenne $400$ FCFA à l'association, soit $40\,000$ FCFA au total, ce qui correspond exactement à la recette ($100\,000$ FCFA) moins la valeur des lots ($60\,000$ FCFA).

\textbf{Variance} (König-Huygens) :
\[ V(X) = E(X^2) - [E(X)]^2 = 16\,800\,000 - (-400)^2 = 16\,800\,000 - 160\,000 = 16\,640\,000. \]
\textbf{Écart-type :} $\sigma(X) = \sqrt{16\,640\,000} \approx 4\,079$ FCFA. L'écart-type est très grand devant l'espérance : le gain est extrêmement dispersé (la plupart perdent $1\,000$ FCFA, un seul gagne beaucoup), ce qui est typique d'une tombola.
\end{exoresolu}

\textbf{Deux lois de même espérance, de variances différentes.}\par
La figure ci-dessous compare les jeux A et B évoqués plus haut : même espérance $1\,000$ FCFA, mais des dispersions très différentes. Pour le jeu A, $V(A) = 100$ et $\sigma(A) = 10$ FCFA ; pour le jeu B, $V(B) = 1\,000\,000$ et $\sigma(B) = 1\,000$ FCFA.

\begin{popfigure}
\begin{center}
\begin{minipage}{0.45\linewidth}
\centering
\begin{tikzpicture}
\begin{axis}[
    ybar, bar width=12pt,
    width=\linewidth, height=5cm,
    ymin=0, ymax=0.6,
    ylabel={Probabilité},
    xlabel={Gain (FCFA)},
    xtick={990, 1010},
    xticklabel style={font=\small, rotate=45},
    ymajorgrids=true, grid style={popline!40},
    title={Jeu A : $\sigma = 10$ FCFA},
    title style={font=\small},
]
\addplot+[ybar, fill=popblue, draw=popdark] coordinates {(990,0.5) (1010,0.5)};
\end{axis}
\end{tikzpicture}
\end{minipage}
\hfill
\begin{minipage}{0.45\linewidth}
\centering
\begin{tikzpicture}
\begin{axis}[
    ybar, bar width=12pt,
    width=\linewidth, height=5cm,
    ymin=0, ymax=0.6,
    ylabel={Probabilité},
    xlabel={Gain (FCFA)},
    xtick={0, 2000},
    xticklabel style={font=\small, rotate=45},
    ymajorgrids=true, grid style={popline!40},
    title={Jeu B : $\sigma = 1\,000$ FCFA},
    title style={font=\small},
]
\addplot+[ybar, fill=poporange, draw=popdark] coordinates {(0,0.5) (2000,0.5)};
\end{axis}
\end{tikzpicture}
\end{minipage}
\end{center}
\legende{Deux lois de même espérance $E = 1\,000$ FCFA mais de variances très différentes : le jeu A est concentré autour de la moyenne, le jeu B est très dispersé. L'échelle horizontale est respectée dans chaque graphique.}
\end{popfigure}

\begin{propriete}[Effet d'une transformation affine sur la variance (admise)]
Soit $X$ une variable aléatoire et $a$, $b$ deux réels. Alors :
\[ V(aX + b) = a^2\,V(X) \qquad \text{et} \qquad \sigma(aX + b) = |a|\,\sigma(X). \]
Ajouter une constante $b$ translate toutes les valeurs sans changer la dispersion ; multiplier par $a$ dilate les écarts, d'où le facteur $a^2$ sur la variance.
\end{propriete}

\begin{exemplebox}[Application directe]
Si $V(X) = 4$, alors $V(3X + 5) = 3^2 \times 4 = 36$ et $\sigma(3X+5) = 3 \times 2 = 6$. Remarque que la constante $5$ n'intervient pas.
\end{exemplebox}

\begin{savaistu}[L'espérance, outil des assureurs]
Les compagnies d'assurance fixent leurs primes grâce à l'espérance mathématique des sinistres. Si un assureur estime qu'un conducteur a une probabilité $0{,}02$ d'avoir un accident coûtant en moyenne $2\,000\,000$ FCFA dans l'année, le coût moyen du risque est $E = 0{,}02 \times 2\,000\,000 = 40\,000$ FCFA. La prime annuelle sera donc \emph{supérieure} à $40\,000$ FCFA (on ajoute les frais de gestion et une marge) : l'assurance est, comme le casino, un « jeu » à espérance négative pour le client... mais qui le protège d'une perte catastrophique. C'est le même raisonnement qui permet à la LONACAM ou aux loteries de garantir leurs lots tout en dégageant des recettes.
\end{savaistu}

\begin{exercice}[Application : Jeu de cartes]
Un opérateur de jeux propose le tirage suivant : on tire une carte au hasard dans un jeu de $32$ cartes. Tirer un as rapporte $1\,000$ FCFA, tirer un roi rapporte $500$ FCFA, toute autre carte ne rapporte rien. La mise est de $200$ FCFA. On note $X$ le gain algébrique du joueur.
\begin{enumerate}
\item Déterminer la loi de probabilité de $X$.
\item Calculer $E(X)$. Le jeu est-il équitable, favorable ou défavorable ?
\item Calculer $V(X)$ puis $\sigma(X)$ (arrondir au franc près).
\item Quelle mise rendrait le jeu équitable ?
\end{enumerate}
\end{exercice}

\begin{corrige}[Exercice : Jeu de cartes]
\begin{enumerate}
\item Il y a $4$ as et $4$ rois dans $32$ cartes. Les valeurs de $X$ sont $800$ (as), $300$ (roi), $-200$ (autre) :
$P(X = 800) = \dfrac{4}{32} = \dfrac18$ ; $P(X = 300) = \dfrac18$ ; $P(X = -200) = \dfrac{24}{32} = \dfrac34$.
\item $E(X) = \dfrac18 \times 800 + \dfrac18 \times 300 + \dfrac34 \times (-200) = 100 + 37{,}5 - 150 = -12{,}5$ FCFA. Le jeu est \textbf{défavorable} au joueur.
\item $E(X^2) = \dfrac18 \times 800^2 + \dfrac18 \times 300^2 + \dfrac34 \times (-200)^2 = 80\,000 + 11\,250 + 30\,000 = 121\,250$. D'où $V(X) = 121\,250 - (-12{,}5)^2 = 121\,250 - 156{,}25 = 121\,093{,}75$, et $\sigma(X) = \sqrt{121\,093{,}75} \approx 348$ FCFA.
\item Notons $m$ la mise. Le gain algébrique devient $X' = X + (200 - m)$, donc $E(X') = E(X) + 200 - m = 187{,}5 - m$. Le jeu est équitable si $E(X') = 0$, c'est-à-dire $m = 187{,}5$ FCFA.
\end{enumerate}
\end{corrige}

\begin{aretenir}[L'essentiel de la section]
\begin{itemize}
\item $E(X) = \sum p_i x_i$ : gain moyen sur un grand nombre de parties. Jeu équitable si $E(X) = 0$.
\item Linéarité : $E(aX + b) = aE(X) + b$.
\item Variance : $V(X) = \sum p_i (x_i - E(X))^2 = E(X^2) - [E(X)]^2$ (König-Huygens, formule de calcul à privilégier).
\item Écart-type : $\sigma(X) = \sqrt{V(X)}$ ; il mesure la dispersion dans l'unité de $X$.
\item $V(aX+b) = a^2 V(X)$ : la translation ne change pas la dispersion.
\end{itemize}
\end{aretenir}

\coursec{Épreuve de Bernoulli et schéma de Bernoulli}

Dans la section précédente, tu as appris à associer à une expérience aléatoire une \cle{variable aléatoire} $X$, puis à résumer son comportement par deux nombres : l'\cle{espérance} $E(X)$, qui donne la valeur moyenne attendue, et la \cle{variance} $V(X)$ (ou l'écart-type $\sigma(X)$), qui mesure la dispersion autour de cette moyenne. Nous allons maintenant mettre ces outils au service d'une situation extrêmement fréquente, tant dans la vie courante qu'au Baccalauréat : \emph{répéter plusieurs fois la même expérience qui n'a que deux issues possibles}. Réussite ou échec à un examen blanc, pièce conforme ou défectueuse à la sortie d'une usine de Brasseries du Cameroun, réponse juste ou fausse à une question de QCM : partout, le même modèle mathématique se cache. Ce modèle porte le nom du mathématicien suisse Jacques Bernoulli.

\courssub{Épreuve de Bernoulli et variable de Bernoulli}

\textbf{Intuition.} Imagine que tu lances une pièce de monnaie en te demandant simplement : « Pile ou Face ? ». Ou qu'un technicien de CAMTEL teste une ligne téléphonique : « fonctionnelle ou défectueuse ? ». Dans les deux cas, l'expérience n'a que \emph{deux issues}, et l'une d'elles nous intéresse particulièrement : on l'appelle le \cle{succès}. Toute la richesse du modèle tient dans un seul nombre : la probabilité $p$ du succès.

\begin{definition}[Épreuve de Bernoulli]
On appelle \cle{épreuve de Bernoulli} toute expérience aléatoire qui ne comporte que \textbf{deux issues} :
\begin{itemize}
\item le \cle{succès}, noté $S$, de probabilité $p$ ;
\item l'\cle{échec}, noté $\bar{S}$ (ou $E$), de probabilité $q = 1 - p$.
\end{itemize}
Le nombre $p$ (avec $0 \leq p \leq 1$) est appelé le \cle{paramètre} de l'épreuve de Bernoulli.
\end{definition}

\begin{attention}[Deux issues, pas forcément « naturelles »]
Une expérience peut avoir plusieurs issues et être \emph{ramenée} à une épreuve de Bernoulli en regroupant les issues. Par exemple, on lance un dé à six faces et on s'intéresse uniquement à l'événement « obtenir un 6 » : alors $S = \{\text{obtenir un 6}\}$ avec $p = \dfrac{1}{6}$, et $\bar{S} = \{\text{obtenir 1, 2, 3, 4 ou 5}\}$ avec $q = \dfrac{5}{6}$. C'est toi qui choisis ce que tu appelles « succès », selon la question posée.
\end{attention}

\textbf{La variable de Bernoulli.} À une épreuve de Bernoulli, on associe naturellement une variable aléatoire qui « code » le résultat : $1$ si succès, $0$ si échec.

\begin{definition}[Variable de Bernoulli]
La \cle{variable de Bernoulli} $X$ de paramètre $p$ est la variable aléatoire définie par :
\[ X = \begin{cases} 1 & \text{en cas de succès} \\ 0 & \text{en cas d'échec} \end{cases} \qquad \text{avec } P(X=1) = p \text{ et } P(X=0) = 1-p. \]
On dit que $X$ suit la \cle{loi de Bernoulli} de paramètre $p$, notée $\mathcal{B}(1\,;\,p)$.
\end{definition}

\begin{propriete}[Espérance et variance d'une variable de Bernoulli]
Si $X$ suit la loi de Bernoulli de paramètre $p$, alors :
\[ E(X) = p \qquad \text{et} \qquad V(X) = p(1-p). \]
\end{propriete}

\begin{experience}[Démonstration guidée (exigible)]
On calcule directement à partir des définitions de l'espérance et de la variance.
\begin{enumerate}
\item \textbf{Espérance.} La loi de $X$ ne comporte que deux valeurs :
\[ E(X) = 1 \times P(X=1) + 0 \times P(X=0) = 1 \times p + 0 \times (1-p) = p. \]
\item \textbf{Calcul de $E(X^2)$.} Comme $X$ ne prend que les valeurs $0$ et $1$, on a $X^2 = X$ (car $0^2=0$ et $1^2=1$). Donc :
\[ E(X^2) = 1^2 \times p + 0^2 \times (1-p) = p. \]
\item \textbf{Variance.} On utilise la formule de König-Huygens $V(X) = E(X^2) - \big[E(X)\big]^2$ :
\[ V(X) = p - p^2 = p(1-p). \]
\end{enumerate}
La démonstration est terminée. Retiens la démarche : elle est courte et peut être demandée telle quelle au Baccalauréat.
\end{experience}

\begin{exemplebox}[Lancer d'un dé]
On lance un dé équilibré et on note $X$ la variable qui vaut $1$ si on obtient un 6, $0$ sinon. Alors $X$ suit la loi de Bernoulli de paramètre $p = \dfrac{1}{6}$, et :
\[ E(X) = \frac{1}{6} \approx 0{,}17 \qquad ; \qquad V(X) = \frac{1}{6} \times \frac{5}{6} = \frac{5}{36} \approx 0{,}14. \]
\end{exemplebox}

\courssub{Schéma de Bernoulli : répéter des épreuves identiques et indépendantes}

\textbf{Intuition.} Une seule épreuve, c'est peu d'information. Mais si un opérateur de téléphonie teste $20$ lignes, si un élève répond à $10$ questions d'un QCM, si un contrôleur qualité prélève $50$ sachets de lait en poudre à la sortie d'une chaîne de production, alors on répète \emph{la même} épreuve de Bernoulli, et la question devient : « combien de succès va-t-on obtenir au total ? ». Pour que le modèle soit valable, deux conditions sont indispensables : les épreuves doivent être \textbf{identiques} (même probabilité $p$ à chaque fois) et \textbf{indépendantes} (le résultat d'une épreuve ne modifie pas les probabilités des suivantes).

\begin{definition}[Schéma de Bernoulli]
On appelle \cle{schéma de Bernoulli} la répétition de $n$ épreuves de Bernoulli \textbf{identiques} et \textbf{indépendantes}, où $n$ est un entier naturel non nul. À chaque épreuve, la probabilité du succès est la même, égale à $p$.
\end{definition}

\begin{methode}[Justifier qu'une situation est un schéma de Bernoulli]
Pour pouvoir utiliser ce modèle (et, dans la section suivante, la loi binomiale), tu dois rédiger une justification en trois points :
\begin{enumerate}
\item \textbf{Deux issues.} Chaque épreuve élémentaire n'a que deux issues : succès $S$ de probabilité $p$, échec de probabilité $1-p$. Précise clairement ce que tu appelles succès.
\item \textbf{Même probabilité.} Les épreuves sont identiques : la probabilité $p$ du succès est la même à chaque répétition.
\item \textbf{Indépendance.} Les épreuves sont indépendantes : le résultat de l'une n'influence pas les autres. C'est souvent ici que se joue la validité du modèle (tirages \emph{avec remise}, grande population, etc.).
\end{enumerate}
Conclusion type à rédiger : « On répète $n$ fois, de façon identique et indépendante, une épreuve de Bernoulli de paramètre $p$ : il s'agit donc d'un schéma de Bernoulli de paramètres $n$ et $p$. »
\end{methode}

\textbf{Modélisation par un arbre.} Un schéma de Bernoulli se représente par un arbre à $n$ niveaux : à chaque niveau, deux branches ($S$ de poids $p$, $\bar{S}$ de poids $q = 1-p$). Grâce à l'indépendance, la probabilité d'un chemin est le \emph{produit} des probabilités rencontrées le long de ce chemin. Un chemin qui contient exactement $k$ succès et $n-k$ échecs a donc pour probabilité $p^k (1-p)^{n-k}$, quel que soit l'ordre des succès. Il reste à compter \emph{combien} de chemins contiennent exactement $k$ succès : c'est le nombre de façons de choisir les positions des $k$ succès parmi les $n$ épreuves, c'est-à-dire le coefficient binomial $\dbinom{n}{k}$.

\begin{popfigure}
\begin{center}
\begin{tikzpicture}[
  grow=right, level distance=2.6cm,
  level 1/.style={sibling distance=4.4cm},
  level 2/.style={sibling distance=2.2cm},
  level 3/.style={sibling distance=1.1cm},
  every node/.style={font=\small},
  branche/.style={draw=popmuted},
  brancheok/.style={draw=poporange, very thick},
  feuille/.style={font=\footnotesize}
]
\node {$\bullet$}
  child { node[feuille] {$\bar{S}$}
    child { node[feuille] {$\bar{S}$}
      child { node[feuille] {$\bar{S}$ \; $0$ succès} edge from parent[branche] node[below, font=\scriptsize] {$q$} }
      child { node[feuille] {$S$ \; $1$ succès} edge from parent[branche] node[above, font=\scriptsize] {$p$} }
      edge from parent[branche] node[below, font=\scriptsize] {$q$} }
    child { node[feuille] {$S$}
      child { node[feuille] {$\bar{S}$ \; $1$ succès} edge from parent[branche] node[below, font=\scriptsize] {$q$} }
      child { node[feuille, fill=poporangeL, rounded corners=1pt] {$S$ \; $\mathbf{2}$ \textbf{succès}} edge from parent[brancheok] node[above, font=\scriptsize] {$p$} }
      edge from parent[brancheok] node[above, font=\scriptsize] {$p$} }
    edge from parent[branche] node[below, font=\scriptsize] {$q$} }
  child { node[feuille] {$S$}
    child { node[feuille] {$\bar{S}$}
      child { node[feuille] {$\bar{S}$ \; $1$ succès} edge from parent[branche] node[below, font=\scriptsize] {$q$} }
      child { node[feuille, fill=poporangeL, rounded corners=1pt] {$S$ \; $\mathbf{2}$ \textbf{succès}} edge from parent[brancheok] node[above, font=\scriptsize] {$p$} }
      edge from parent[brancheok] node[below, font=\scriptsize] {$q$} }
    child { node[feuille, fill=poporangeL, rounded corners=1pt] {$S$}
      child { node[feuille, fill=poporangeL, rounded corners=1pt] {$\bar{S}$ \; $\mathbf{2}$ \textbf{succès}} edge from parent[brancheok] node[below, font=\scriptsize] {$q$} }
      child { node[feuille] {$S$ \; $3$ succès} edge from parent[branche] node[above, font=\scriptsize] {$p$} }
      edge from parent[brancheok] node[above, font=\scriptsize] {$p$} }
    edge from parent[brancheok] node[above, font=\scriptsize] {$p$} };
\end{tikzpicture}
\end{center}
\legende{Arbre d'un schéma de Bernoulli à $n = 3$ épreuves. Les trois chemins épais et orangés contiennent exactement $k = 2$ succès : il y en a $\dbinom{3}{2} = 3$, chacun de probabilité $p^2 q$.}
\end{popfigure}

\begin{aretenir}[Compter les chemins]
Dans un schéma de Bernoulli à $n$ épreuves, le nombre de chemins de l'arbre contenant exactement $k$ succès est :
\[ \binom{n}{k} = \frac{n!}{k!\,(n-k)!}. \]
Chacun de ces chemins a pour probabilité $p^k(1-p)^{n-k}$. Cette observation est la clé de la loi binomiale, que nous étudierons dans la section suivante.
\end{aretenir}

\begin{exemplebox}[Vérification sur l'arbre à 3 épreuves]
Pour $n = 3$, les chemins à $2$ succès sont $(S,S,\bar{S})$, $(S,\bar{S},S)$ et $(\bar{S},S,S)$ : on retrouve bien $\dbinom{3}{2} = 3$ chemins, chacun de probabilité $p \times p \times q = p^2 q$. La probabilité d'obtenir exactement $2$ succès est donc $3p^2q$. Par exemple, avec $p = 0{,}5$ (pièce équilibrée), elle vaut $3 \times 0{,}25 \times 0{,}5 = 0{,}375$.
\end{exemplebox}

\courssub{Reconnaître (ou rejeter) un schéma de Bernoulli}

\textbf{Exemples où le modèle s'applique.}
\begin{itemize}
\item \textbf{Tirages avec remise.} On tire une boule dans une urne, on note le résultat, \emph{on la remet}, et on recommence. L'urne est dans le même état à chaque tirage : épreuves identiques et indépendantes.
\item \textbf{Contrôle qualité en grande série.} Une usine produit des milliers de pièces par jour ; on en prélève $n$ au hasard. Même si le prélèvement se fait techniquement sans remise, le prélèvement de quelques pièces ne modifie pratiquement pas la proportion de défectueuses : on assimile à des tirages avec remise (voir la discussion ci-dessous).
\item \textbf{Réponses au hasard à un QCM.} Un candidat répond au hasard et indépendamment à chaque question : chaque question est une épreuve de Bernoulli dont le paramètre dépend du nombre d'options proposées.
\end{itemize}

\begin{exoresolu}[Un élève répond au hasard à un QCM]
Un élève répond complètement au hasard à un QCM de $10$ questions. Chaque question propose $2$ options, dont une seule est correcte. Justifier que le nombre $X$ de bonnes réponses relève d'un schéma de Bernoulli, et préciser ses paramètres. Quelle est la probabilité qu'il obtienne exactement $5$ bonnes réponses ?
\tcblower
\textbf{Solution.}
\begin{enumerate}
\item \textbf{Deux issues.} Pour chaque question, deux issues : succès $S$ « la réponse est correcte » de probabilité $p = \dfrac{1}{2}$ (une option correcte parmi 2, choix au hasard), ou échec de probabilité $1-p = \dfrac{1}{2}$.
\item \textbf{Épreuves identiques.} Chaque question a la même structure (2 options, réponse au hasard) : $p = \dfrac{1}{2}$ à chaque épreuve.
\item \textbf{Indépendance.} L'élève répond au hasard à chaque question sans tenir compte des réponses précédentes : les épreuves sont indépendantes.
\end{enumerate}
On répète donc $n = 10$ fois, de façon identique et indépendante, une épreuve de Bernoulli de paramètre $p = \dfrac{1}{2}$ : c'est un schéma de Bernoulli de paramètres $n = 10$ et $p = \dfrac{1}{2}$.

Pour obtenir exactement $5$ bonnes réponses, il faut choisir les $5$ questions réussies parmi $10$ : il y a $\dbinom{10}{5} = 252$ chemins, chacun de probabilité $\left(\dfrac{1}{2}\right)^{5}\left(\dfrac{1}{2}\right)^{5} = \left(\dfrac{1}{2}\right)^{10}$. Donc :
\[ P(X = 5) = \binom{10}{5}\left(\frac{1}{2}\right)^{10} = \frac{252}{1024} \approx 0{,}246. \]
L'élève a environ $24,6\,\%$ de chances d'obtenir exactement la moyenne en répondant au hasard.
\end{exoresolu}

\textbf{Contre-exemple : les tirages sans remise.} Une urne contient $4$ boules rouges et $6$ boules noires. On tire une première boule : la probabilité qu'elle soit rouge est $\dfrac{4}{10} = 0{,}4$. Si on ne la remet pas et qu'elle était rouge, la probabilité de tirer une rouge au deuxième tirage devient $\dfrac{3}{9} = \dfrac{1}{3}$ ; si elle était noire, elle devient $\dfrac{4}{9}$. La probabilité du succès \emph{dépend du résultat précédent} : les épreuves ne sont ni identiques ni indépendantes. Ce n'est \textbf{pas} un schéma de Bernoulli.

\begin{attention}[Erreur fréquente : appliquer la loi binomiale à des tirages sans remise]
C'est l'un des pièges classiques du Baccalauréat. Avant d'écrire « $X$ suit une loi binomiale », vérifie impérativement l'\textbf{indépendance} des épreuves. Des tirages \emph{sans remise} dans une \emph{petite} population ne constituent pas un schéma de Bernoulli : dans ce cas, il faut revenir aux probabilités conditionnelles et à l'arbre pondéré, ou au dénombrement direct. Appliquer la formule binomiale sans justification fait perdre des points, même si le calcul numérique est correct.
\end{attention}

\textbf{Discussion nuancée : quand peut-on assimiler à des tirages avec remise ?} Supposons maintenant qu'un contrôleur prélève $10$ sachets dans un stock de $10\,000$ sachets dont $5\,\%$ sont défectueux, \emph{sans remise}. Après le premier tirage d'un sachet défectueux, la proportion passe de $\dfrac{500}{10\,000} = 0{,}05$ à $\dfrac{499}{9\,999} \approx 0{,}0499$ : la variation est négligeable. Lorsque la taille de la population est \emph{très grande} devant le nombre de tirages, les épreuves sont \emph{quasi indépendantes} et on assimile les tirages sans remise à des tirages avec remise : le schéma de Bernoulli (et donc la loi binomiale) devient une excellente approximation. C'est exactement ce qui justifie le contrôle qualité en grande série évoqué plus haut, et c'est une formulation souvent attendue dans les sujets du Bac : « la population étant très grande, on assimile ce prélèvement à un tirage avec remise ».

\begin{savaistu}[Jacques Bernoulli et la loi des grands nombres]
Jacques Bernoulli (1654--1705), membre d'une illustre famille de mathématiciens suisses, a consacré vingt ans à son ouvrage \emph{Ars Conjectandi}, publié après sa mort en 1713. Il y démontre la première version de la \emph{loi des grands nombres} : quand on répète un schéma de Bernoulli un très grand nombre de fois, la fréquence observée du succès se rapproche de la probabilité $p$. C'est le fondement théorique des sondages : interroger $1\,000$ électeurs choisis au hasard permet d'estimer une intention de vote à l'échelle de tout le Cameroun, avec une marge d'erreur que les mathématiques savent quantifier.
\end{savaistu}

\begin{exercice}[À toi de jouer]
Une entreprise de Douala produit des câbles électriques. On estime que $2\,\%$ des câbles présentent un défaut. On prélève au hasard $20$ câbles dans la production journalière, très importante.
\begin{enumerate}
\item Justifier que le nombre $X$ de câbles défectueux dans le prélèvement relève d'un schéma de Bernoulli dont on précisera les paramètres.
\item Combien de chemins de l'arbre (théorique) conduisent à exactement $2$ câbles défectueux ?
\item Écrire, sans la calculer, la probabilité d'obtenir exactement $2$ câbles défectueux.
\end{enumerate}
\end{exercice}

\begin{corrige}[Exercice : câbles électriques]
\begin{enumerate}
\item Chaque câble prélevé donne lieu à une épreuve à deux issues : succès $S$ « le câble est défectueux » de probabilité $p = 0{,}02$, ou échec de probabilité $0{,}98$. La production étant très importante, on assimile le prélèvement à $20$ tirages avec remise : les épreuves sont identiques et indépendantes. C'est donc un schéma de Bernoulli de paramètres $n = 20$ et $p = 0{,}02$.
\item Le nombre de chemins contenant exactement $2$ succès est $\dbinom{20}{2} = \dfrac{20 \times 19}{2} = 190$.
\item Chaque chemin à $2$ succès a pour probabilité $(0{,}02)^2 \times (0{,}98)^{18}$, donc :
\[ P(X = 2) = \binom{20}{2} (0{,}02)^2 (0{,}98)^{18} = 190 \times (0{,}02)^2 \times (0{,}98)^{18}. \]
\end{enumerate}
\end{corrige}

Tu sais désormais reconnaître, justifier et représenter un schéma de Bernoulli. Dans la section suivante, nous exploiterons le comptage des chemins par les coefficients $\dbinom{n}{k}$ pour définir la \cle{loi binomiale} et en donner les formules d'espérance et de variance, qui te permettront de résoudre efficacement les problèmes du Baccalauréat.

\coursec{Loi binomiale et applications}

Dans la section précédente, nous avons appris à reconnaître un \cle{schéma de Bernoulli} : la répétition de $n$ épreuves de Bernoulli identiques et indépendantes, chacune ayant une probabilité $p$ de succès. Nous avons compté les chemins d'un arbre pour calculer, par exemple, la probabilité d'obtenir exactement deux succès en trois épreuves. Mais que faire si l'on répète l'épreuve $10$, $50$ ou $1000$ fois ? Dessiner l'arbre devient impossible. C'est précisément là qu'intervient la \cle{loi binomiale} : elle fournit une \emph{formule directe} qui remplace l'arbre. Cette loi est l'un des outils les plus utilisés de la vie courante : contrôle qualité dans une usine des Brasseries du Cameroun, sondage d'opinion avant une élection, efficacité d'un vaccin dans un centre de santé de Yaoundé. Cette section te donne la maîtrise complète de cet outil, très fréquent au Baccalauréat.

\courssub{Définition de la loi binomiale}

Reprenons la situation générale. On répète $n$ fois, de façon identique et indépendante, une épreuve de Bernoulli de probabilité de succès $p$. À l'issue des $n$ épreuves, on s'intéresse au \emph{nombre total de succès} obtenus. Ce nombre est aléatoire : il peut valoir $0$ (aucun succès), $1$, $2$, \dots, jusqu'à $n$ (que des succès). C'est donc une variable aléatoire, et sa loi porte un nom.

\begin{definition}[Loi binomiale]
Soit $n$ un entier naturel non nul et $p$ un réel de $[0\,;1]$. On considère un schéma de Bernoulli de $n$ épreuves, de probabilité de succès $p$. La variable aléatoire $X$ qui compte le \cle{nombre de succès} obtenus au cours des $n$ épreuves suit la \cle{loi binomiale} de paramètres $n$ et $p$. On note :
\[ X \hookrightarrow \mathcal{B}(n\,;p). \]
Les valeurs prises par $X$ sont les entiers $0, 1, 2, \dots, n$.
\end{definition}

\begin{exemplebox}[Reconnaître une loi binomiale]
Un opérateur téléphonique constate qu'un appel passé en heure de pointe aboutit avec une probabilité $0{,}8$, indépendamment des autres appels. Un client passe $10$ appels. La variable $X$ qui compte le nombre d'appels aboutis suit la loi $\mathcal{B}(10\,;0{,}8)$ : $n = 10$ répétitions identiques et indépendantes, succès « l'appel aboutit » de probabilité $p = 0{,}8$.
\end{exemplebox}

\begin{attention}[Vérifier les conditions avant d'écrire $\mathcal{B}(n\,;p)$]
On ne peut invoquer la loi binomiale que si les trois conditions sont réunies : (1) même épreuve répétée $n$ fois ; (2) deux issues seulement (succès/échec) ; (3) \textbf{indépendance} des épreuves. En particulier, un tirage \emph{sans remise} dans une petite population n'est pas un schéma de Bernoulli. En revanche, un tirage avec remise, ou un prélèvement dans une très grande population (que l'on assimile alors à un tirage avec remise), convient. Au Baccalauréat, cette justification doit être rédigée explicitement : elle vaut des points.
\end{attention}

\courssub{Formule de la loi binomiale}

\begin{propriete}[Formule de la loi binomiale]
Si $X \hookrightarrow \mathcal{B}(n\,;p)$, alors pour tout entier $k$ tel que $0 \leq k \leq n$ :
\[ P(X = k) = \binom{n}{k} p^k (1-p)^{\,n-k}, \]
où $\binom{n}{k} = \dfrac{n!}{k!(n-k)!}$ est le nombre de façons de choisir $k$ positions parmi $n$.
\end{propriete}

Cette formule n'est pas tombée du ciel : elle se comprend entièrement grâce à l'arbre de probabilités. Suivons le raisonnement, dont la justification est exigible au Baccalauréat.

\begin{experience}[Pourquoi la formule est-elle vraie ?]
Plaçons-nous dans un schéma de Bernoulli à $n$ épreuves, avec $p$ la probabilité de succès et $q = 1-p$ celle d'échec.
\begin{enumerate}
\item \textbf{Probabilité d'un chemin.} Réaliser l'événement « $X = k$ », c'est suivre un chemin de l'arbre comportant exactement $k$ succès et $n-k$ échecs. Par indépendance des épreuves, la probabilité d'un tel chemin est le produit des probabilités rencontrées le long des branches : $k$ facteurs $p$ et $n-k$ facteurs $q$, soit $p^k q^{\,n-k}$. \emph{Tous} les chemins à $k$ succès ont donc la \emph{même} probabilité.
\item \textbf{Combien de tels chemins ?} Un chemin à $k$ succès est entièrement déterminé par le \emph{choix des positions} des $k$ succès parmi les $n$ épreuves. Il y a exactement $\binom{n}{k}$ tels choix, donc $\binom{n}{k}$ chemins.
\item \textbf{Conclusion.} L'événement « $X = k$ » est la réunion de $\binom{n}{k}$ événements (chemins) deux à deux incompatibles, chacun de probabilité $p^k q^{\,n-k}$. D'où :
\[ P(X = k) = \binom{n}{k} p^k (1-p)^{\,n-k}. \]
\end{enumerate}
\end{experience}

\begin{exemplebox}[Application directe]
Reprenons les appels téléphoniques : $X \hookrightarrow \mathcal{B}(10\,;0{,}8)$. La probabilité qu'exactement $7$ appels sur $10$ aboutissent est :
\[ P(X = 7) = \binom{10}{7} (0{,}8)^7 (0{,}2)^3 = 120 \times 0{,}2097 \times 0{,}008 \approx 0{,}201. \]
Il y a environ $20\,\%$ de chances qu'exactement $7$ appels sur $10$ aboutissent.
\end{exemplebox}

\begin{attention}[Les deux erreurs classiques au Bac]
\begin{itemize}
\item \textbf{Oublier le coefficient $\binom{n}{k}$.} Écrire $P(X = k) = p^k(1-p)^{n-k}$ ne donne que la probabilité d'\emph{un seul} chemin (par exemple « les $k$ premières épreuves sont des succès »). C'est l'erreur la plus fréquente : le coefficient binomial compte \emph{tous} les chemins favorables.
\item \textbf{Intervertir les exposants.} L'exposant de $p$ est $k$ (le nombre de succès), celui de $1-p$ est $n-k$ (le nombre d'échecs). Vérifie toujours que la somme des exposants vaut $n$.
\end{itemize}
\end{attention}

Puisque $X$ est une variable aléatoire, la somme des probabilités de toutes ses valeurs doit valoir $1$. Vérifions-le : c'est une magnifique application de la \emph{formule du binôme de Newton}, vue dans le chapitre sur le dénombrement.

\begin{experience}[Vérification que $\sum_{k=0}^{n} P(X=k) = 1$]
Rappelons la formule du binôme de Newton : pour tous réels $a$ et $b$,
\[ (a+b)^n = \sum_{k=0}^{n} \binom{n}{k} a^k b^{\,n-k}. \]
Appliquons-la avec $a = p$ et $b = 1-p$ :
\[ \sum_{k=0}^{n} P(X = k) = \sum_{k=0}^{n} \binom{n}{k} p^k (1-p)^{\,n-k} = \bigl(p + (1-p)\bigr)^n = 1^n = 1. \]
La loi binomiale est donc bien une loi de probabilité. Retiens le lien profond : le nom « binomiale » vient précisément du binôme de Newton, et les coefficients $\binom{n}{k}$ sont ceux du triangle de Pascal.
\end{experience}

\courssub{Espérance et variance de la loi binomiale}

\begin{propriete}[Espérance et variance — résultats admis]
Si $X \hookrightarrow \mathcal{B}(n\,;p)$, alors :
\[ E(X) = np \qquad \text{et} \qquad V(X) = np(1-p), \quad \text{donc} \quad \sigma(X) = \sqrt{np(1-p)}. \]
Ces formules sont \textbf{admises} au programme de Terminale ; leur démonstration n'est pas exigible.
\end{propriete}

L'intuition derrière $E(X) = np$ est limpide : si chaque épreuve réussit en moyenne une proportion $p$ du temps, alors sur $n$ épreuves on attend en moyenne $n \times p$ succès. Sur $10$ appels qui aboutissent $8$ fois sur $10$, on « attend » $10 \times 0{,}8 = 8$ appels aboutis : c'est exactement le bon sens.

\begin{exemplebox}[Vérification numérique sur un petit cas]
Prenons $X \hookrightarrow \mathcal{B}(2\,;0{,}5)$ (deux lancers d'une pièce équilibrée, $X$ = nombre de « Pile »). La loi de $X$, calculée avec la formule, est :
\begin{center}
\begin{tabularx}{\linewidth}{|l|X|X|X|}
\hline
\rowcolor{popblueL} $k$ & $0$ & $1$ & $2$ \\
\hline
$P(X=k)$ & $0{,}25$ & $0{,}5$ & $0{,}25$ \\
\hline
\end{tabularx}
\end{center}
Calculons l'espérance « à la main » : $E(X) = 0 \times 0{,}25 + 1 \times 0{,}5 + 2 \times 0{,}25 = 1$. La formule donne $np = 2 \times 0{,}5 = 1$ : accord parfait. Pour la variance : $V(X) = E(X^2) - [E(X)]^2 = (0 \times 0{,}25 + 1 \times 0{,}5 + 4 \times 0{,}25) - 1^2 = 1{,}5 - 1 = 0{,}5$, et la formule donne $np(1-p) = 2 \times 0{,}5 \times 0{,}5 = 0{,}5$. Les formules admises sont bien cohérentes avec les définitions.
\end{exemplebox}

\courssub{Calculs types avec la loi binomiale}

Au Baccalauréat, trois types de questions reviennent sans cesse. Voici la stratégie pour chacun.

\begin{methode}[Les trois calculs types]
Soit $X \hookrightarrow \mathcal{B}(n\,;p)$.
\begin{enumerate}
\item \textbf{Probabilité d'une valeur exacte} : appliquer directement la formule $P(X = k) = \binom{n}{k} p^k (1-p)^{n-k}$.
\item \textbf{Probabilité cumulée} $P(X \leq k)$ : additionner les probabilités $P(X=0) + P(X=1) + \dots + P(X=k)$. Il n'y a pas de raccourci : on calcule chaque terme puis on somme. De même, $P(X \geq k) = P(X=k) + \dots + P(X=n)$, ou encore $1 - P(X \leq k-1)$ si cela demande moins de calculs.
\item \textbf{Probabilité « au moins un succès »} : passer par l'événement contraire. L'événement « $X \geq 1$ » a pour contraire « $X = 0$ » (aucun succès, c'est-à-dire $n$ échecs), d'où :
\[ P(X \geq 1) = 1 - P(X = 0) = 1 - (1-p)^n. \]
Cette astuce transforme une somme de $n$ termes en un seul calcul : utilise-la systématiquement dès que tu lis « au moins un ».
\end{enumerate}
\end{methode}

\begin{exoresolu}[Contrôle qualité dans une usine]
Une machine d'une usine de conditionnement à Douala produit $2\,\%$ de pièces défectueuses. On prélève au hasard un lot de $50$ pièces dans la production (la production est suffisamment importante pour assimiler le prélèvement à un tirage avec remise). On note $X$ le nombre de pièces défectueuses dans le lot.
\begin{enumerate}
\item Justifier que $X$ suit une loi binomiale et préciser ses paramètres.
\item Calculer la probabilité que le lot contienne \emph{au plus une} pièce défectueuse (arrondir à $10^{-3}$).
\item Quel est le nombre moyen de pièces défectueuses attendu par lot ?
\end{enumerate}
\tcblower
\textbf{1.} Chaque pièce prélevée est une épreuve de Bernoulli : succès « la pièce est défectueuse » de probabilité $p = 0{,}02$. On répète $n = 50$ épreuves, identiques et indépendantes (prélèvement assimilé à un tirage avec remise). Donc $X \hookrightarrow \mathcal{B}(50\,;0{,}02)$.

\textbf{2.} « Au plus une pièce défectueuse » signifie $X \leq 1$, c'est-à-dire $X = 0$ ou $X = 1$ :
\begin{align*}
P(X = 0) &= \binom{50}{0}(0{,}02)^0(0{,}98)^{50} = (0{,}98)^{50} \approx 0{,}3642, \\
P(X = 1) &= \binom{50}{1}(0{,}02)^1(0{,}98)^{49} = 50 \times 0{,}02 \times (0{,}98)^{49} \approx 0{,}3716.
\end{align*}
Donc $P(X \leq 1) = P(X=0) + P(X=1) \approx 0{,}3642 + 0{,}3716 \approx 0{,}736$. Il y a environ $73{,}6\,\%$ de chances qu'un lot de $50$ pièces contienne au plus une pièce défectueuse.

\textbf{3.} Le nombre moyen attendu est l'espérance : $E(X) = np = 50 \times 0{,}02 = 1$. En moyenne, on attend \textbf{une pièce défectueuse par lot} de $50$ pièces.
\end{exoresolu}

\courssub{Représentation graphique et forme de la loi}

Représentons la loi binomiale par un diagramme en barres : en abscisse les valeurs $k$, en ordonnée les probabilités $P(X=k)$. Voici la loi $\mathcal{B}(10\,;0{,}3)$.

\begin{popfigure}
\begin{center}
\begin{tikzpicture}
\begin{axis}[
    width=0.85\linewidth, height=6.5cm,
    ybar, bar width=8pt,
    xlabel={$k$}, ylabel={$P(X=k)$},
    xmin=-0.7, xmax=10.7, ymin=0, ymax=0.30,
    xtick={0,1,2,3,4,5,6,7,8,9,10},
    ytick={0,0.05,0.10,0.15,0.20,0.25},
    axis lines=left,
    tick label style={font=\small},
    label style={font=\small}
]
\addplot[fill=popblue, draw=popdark] coordinates {
(0,0.0282) (1,0.1211) (2,0.2335) (3,0.2668) (4,0.2001)
(5,0.1029) (6,0.0368) (7,0.0090) (8,0.0014) (9,0.0001) (10,0.000006)
};
\end{axis}
\end{tikzpicture}
\end{center}
\legende{Diagramme en barres de la loi $\mathcal{B}(10\,;0{,}3)$. Le maximum est atteint en $k = 3 = E(X)$ : les valeurs les plus probables se concentrent autour de l'espérance.}
\end{popfigure}

On observe que la barre la plus haute correspond à $k = 3$, précisément l'espérance $np = 10 \times 0{,}3 = 3$. Ce n'est pas un hasard : la loi binomiale se concentre toujours autour de son espérance. Observons maintenant l'effet du paramètre $p$ en comparant deux lois à $n = 20$ fixé.

\begin{popfigure}
\begin{center}
\begin{tikzpicture}
\begin{axis}[
    width=0.9\linewidth, height=7cm,
    ybar, bar width=3.2pt,
    xlabel={$k$}, ylabel={$P(X=k)$},
    xmin=-0.8, xmax=20.8, ymin=0, ymax=0.30,
    xtick={0,2,4,6,8,10,12,14,16,18,20},
    axis lines=left,
    tick label style={font=\small},
    label style={font=\small},
    legend style={font=\small, at={(0.97,0.97)}, anchor=north east}
]
\addplot[fill=popblue, draw=popdark] coordinates {
(0,0.000001) (1,0.000019) (2,0.000181) (3,0.001087) (4,0.004621)
(5,0.014786) (6,0.036964) (7,0.073929) (8,0.120134) (9,0.160179)
(10,0.176197) (11,0.160179) (12,0.120134) (13,0.073929) (14,0.036964)
(15,0.014786) (16,0.004621) (17,0.001087) (18,0.000181) (19,0.000019) (20,0.000001)
};
\addlegendentry{$\mathcal{B}(20\,;0{,}5)$}
\addplot[fill=poporange, draw=popdark] coordinates {
(0,0.1216) (1,0.2702) (2,0.2852) (3,0.1901) (4,0.0898)
(5,0.0319) (6,0.0089) (7,0.0020) (8,0.0004) (9,0.0001)
(10,0.00001) (11,0) (12,0) (13,0) (14,0) (15,0) (16,0) (17,0) (18,0) (19,0) (20,0)
};
\addlegendentry{$\mathcal{B}(20\,;0{,}1)$}
\end{axis}
\end{tikzpicture}
\end{center}
\legende{Comparaison des lois $\mathcal{B}(20\,;0{,}5)$ et $\mathcal{B}(20\,;0{,}1)$. Pour $p = 0{,}5$, la loi est symétrique en forme de cloche centrée sur $np = 10$ ; pour $p = 0{,}1$, elle est décalée vers la gauche, centrée sur $np = 2$.}
\end{popfigure}

\begin{aretenir}[Lecture graphique de la loi binomiale]
\begin{itemize}
\item Le diagramme en barres de $\mathcal{B}(n\,;p)$ est centré autour de l'espérance $np$.
\item Pour $p = 0{,}5$, le diagramme est \textbf{symétrique} et présente une forme en \emph{cloche}.
\item Pour $p$ proche de $0$ (respectivement de $1$), le diagramme est décalé vers les petites (respectivement grandes) valeurs de $k$ : il est \emph{asymétrique}.
\item Quand $n$ grandit, la forme en cloche devient de plus en plus régulière : c'est l'embryon de la fameuse \emph{courbe en cloche} (loi normale) que tu rencontreras dans les études supérieures.
\end{itemize}
\end{aretenir}

\courssub{Applications de la loi binomiale}

La loi binomiale intervient dès qu'on répète une même observation à deux issues. Voici les grands domaines d'application, avec des situations que tu peux rencontrer au Baccalauréat comme dans la vie.

\textbf{Contrôle qualité.} C'est l'exemple traité plus haut : un fabricant teste un échantillon de produits et décide d'accepter ou de rejeter un lot selon le nombre de pièces défectueuses. La règle de décision (« rejeter le lot si $X \geq 2$ ») repose entièrement sur des calculs binomiaux.

\textbf{Sondages.} Avant une élection, un institut interroge $n = 1000$ personnes. Si la proportion réelle d'électeurs favorables à un candidat est $p$, le nombre $X$ de réponses favorables suit approximativement $\mathcal{B}(1000\,;p)$. On attend $E(X) = 1000\,p$ réponses favorables, avec un écart-type $\sqrt{1000\,p(1-p)}$ qui mesure la \emph{marge d'erreur} du sondage.

\textbf{Jeux de hasard.} Un joueur mise $n$ fois de suite à un jeu où sa probabilité de gain est $p$. Son nombre de gains suit $\mathcal{B}(n\,;p)$. La probabilité de gagner \emph{au moins une fois} est $1-(1-p)^n$ : c'est elle qui explique pourquoi « jouer souvent augmente les chances de gagner au moins une fois », tout en rappelant que l'espérance de gain reste généralement défavorable.

\textbf{Médecine.} Un traitement guérit une maladie avec une probabilité $p = 0{,}9$. Administré à $n = 20$ patients (guérisons supposées indépendantes), le nombre de guérisons suit $\mathcal{B}(20\,;0{,}9)$. On attend $E(X) = 18$ guérisons, et la probabilité que \emph{tous} les patients guérissent est $P(X = 20) = 0{,}9^{20} \approx 0{,}12$ : seulement $12\,\%$, alors que chaque patient a $90\,\%$ de chances de guérir ! Ce genre de calcul aide les médecins et les épidémiologistes à évaluer l'efficacité réelle d'un traitement ou d'une campagne de vaccination.

\begin{savaistu}[La loi binomiale, pilier des sondages modernes]
C'est Jacques Bernoulli (1654--1705) qui, dans son ouvrage \emph{Ars Conjectandi}, a posé les fondations mathématiques de ces situations répétées — d'où le nom « épreuve de Bernoulli ». Trois siècles plus tard, la même formule $\binom{n}{k}p^k(1-p)^{n-k}$ est utilisée chaque jour : par les instituts de sondage pour annoncer « ce candidat est crédité de $52\,\%$ des intentions de vote, à plus ou moins $3$ points près », par les opérateurs comme MTN ou Orange Cameroun pour dimensionner leurs réseaux (combien d'abonnés appelleront simultanément ?), et par les laboratoires pharmaceutiques pour décider si un nouveau médicament est réellement efficace. Maîtriser la loi binomiale, c'est comprendre la logique cachée derrière tous ces chiffres.
\end{savaistu}

\begin{exercice}[Pour t'entraîner]
Un vaccin est efficace dans $95\,\%$ des cas. On vaccine $30$ personnes d'un quartier de Bafoussam (efficacités supposées indépendantes). On note $X$ le nombre de personnes pour lesquelles le vaccin est efficace.
\begin{enumerate}
\item Quelle loi suit $X$ ? Préciser les paramètres.
\item Calculer la probabilité que le vaccin soit efficace pour exactement $28$ personnes (arrondir à $10^{-3}$).
\item Calculer la probabilité que le vaccin soit inefficace pour \emph{au moins une} personne.
\item Combien de personnes peut-on espérer voir protégées ?
\end{enumerate}
\end{exercice}

\begin{corrige}[Exercice]
\textbf{1.} On répète $n = 30$ épreuves de Bernoulli identiques et indépendantes, de probabilité de succès $p = 0{,}95$ (« le vaccin est efficace »). Donc $X \hookrightarrow \mathcal{B}(30\,;0{,}95)$.

\textbf{2.} $P(X = 28) = \binom{30}{28}(0{,}95)^{28}(0{,}05)^{2} = 435 \times (0{,}95)^{28} \times 0{,}0025$. Or $(0{,}95)^{28} \approx 0{,}2377$, donc $P(X=28) \approx 435 \times 0{,}2377 \times 0{,}0025 \approx 0{,}259$.

\textbf{3.} « Inefficace pour au moins une personne » est le contraire de « efficace pour les $30$ personnes » :
\[ P(X \leq 29) = 1 - P(X = 30) = 1 - (0{,}95)^{30} \approx 1 - 0{,}2146 \approx 0{,}785. \]
Il y a donc environ $78{,}5\,\%$ de chances qu'au moins une personne ne soit pas protégée, malgré une efficacité individuelle de $95\,\%$.

\textbf{4.} $E(X) = np = 30 \times 0{,}95 = 28{,}5$ : on peut espérer qu'environ $28$ à $29$ personnes soient protégées.
\end{corrige}

\begin{aretenir}[L'essentiel de la section]
\begin{itemize}
\item $X$ compte le nombre de succès dans un schéma de Bernoulli à $n$ épreuves de paramètre $p$ : $X \hookrightarrow \mathcal{B}(n\,;p)$.
\item Formule fondamentale : $P(X = k) = \binom{n}{k} p^k (1-p)^{\,n-k}$, pour $k = 0, 1, \dots, n$. Le coefficient $\binom{n}{k}$ compte les chemins de l'arbre à $k$ succès.
\item Espérance et variance (admises) : $E(X) = np$ et $V(X) = np(1-p)$, d'où $\sigma(X) = \sqrt{np(1-p)}$.
\item Réflexe « au moins un » : $P(X \geq 1) = 1 - (1-p)^n$.
\item Le diagramme en barres est centré sur $np$, en forme de cloche symétrique quand $p = 0{,}5$.
\end{itemize}
\end{aretenir}

\newpage

\coursec{Probabilités conditionnelles et probabilités totales}

Dans la leçon commune, tu as appris à calculer des probabilités par équiprobabilité, à utiliser les arbres et à reconnaître un schéma de Bernoulli. Mais dans la vie réelle, l'information arrive souvent \emph{progressivement} : un test de dépistage est positif, un client a déjà attendu dix minutes, la première carte tirée est un as. Comment cette information modifie-t-elle la probabilité d'un événement ? C'est toute la question des \cle{probabilités conditionnelles}.

\courssub{Probabilité conditionnelle : définition et intuition}

\begin{experience}[Une intuition pour commencer]
Dans une classe de Terminale D de 40 élèves, 18 sont des filles et 12 élèves portent des lunettes, dont 5 filles. On choisit un élève au hasard. La probabilité qu'il porte des lunettes est $\dfrac{12}{40}=0{,}3$. Mais si l'on \emph{sait} que l'élève choisi est une fille, l'univers se réduit aux 18 filles : la probabilité devient $\dfrac{5}{18}\approx 0{,}28$. L'information « c'est une fille » a changé la probabilité. On note $P_F(L)=\dfrac{5}{18}$ : probabilité de $L$ \emph{sachant} $F$.
\end{experience}

\begin{definition}[Probabilité conditionnelle]
Soit $A$ et $B$ deux événements d'un univers $\Omega$, avec $P(A)\neq 0$. On appelle \cle{probabilité de $B$ sachant $A$} le nombre
\[ P_A(B)=\dfrac{P(A\cap B)}{P(A)}. \]
On le note aussi $P(B\mid A)$. Il se lit : « probabilité que $B$ se réalise, sachant que $A$ est réalisé ».
\end{definition}

\begin{attention}[Ne pas confondre $P_A(B)$ et $P_B(A)$]
En général $P_A(B)\neq P_B(A)$. Dans l'exemple de la classe : $P_F(L)=\dfrac{5}{18}$ (lunettes sachant fille), mais $P_L(F)=\dfrac{5}{12}$ (fille sachant lunettes). Vérifie toujours \emph{quel événement est la condition} : c'est celui qui est déjà réalisé, celui qui réduit l'univers.
\end{attention}

\begin{propriete}[Probabilités composées]
Pour tous événements $A$ et $B$ avec $P(A)\neq 0$ :
\[ P(A\cap B)=P(A)\times P_A(B). \]
C'est une simple réécriture de la définition, mais c'est la règle qui justifie la multiplication des probabilités le long des branches d'un arbre. Elle se généralise à $n$ événements :
\[ P(A_1\cap\cdots\cap A_n)=P(A_1)\times P_{A_1}(A_2)\times\cdots\times P_{A_1\cap\cdots\cap A_{n-1}}(A_n). \]
\end{propriete}

\courssub{L'arbre pondéré : l'outil de calcul}

\begin{methode}[Construire et exploiter un arbre pondéré]
\begin{enumerate}
\item Identifier l'ordre chronologique ou logique des événements : la première ramification porte sur la condition, la seconde sur l'événement conditionné.
\item Inscrire sur chaque branche la probabilité correspondante (probabilité conditionnelle à partir du deuxième niveau).
\item \textbf{Règle 1} : la probabilité d'un chemin est le \emph{produit} des probabilités de ses branches.
\item \textbf{Règle 2} : la probabilité d'un événement réalisé par plusieurs chemins est la \emph{somme} des probabilités de ces chemins.
\item Vérifier que la somme des probabilités issues d'un même nœud vaut $1$.
\end{enumerate}
\end{methode}

\begin{popfigure}
\begin{tikzpicture}[grow=right, level distance=3.5cm,
level 1/.style={sibling distance=3cm},
level 2/.style={sibling distance=1.5cm},
every node/.style={font=\small, align=center},
edge from parent/.style={draw=popmuted,-{Stealth}}]
\node {}
child { node {$\bar{A}$}
  child { node[align=center] {$\bar{B}$ \\ $P(\bar{A}\cap\bar{B})$} edge from parent node[below,pos=0.6]{$P_{\bar{A}}(\bar{B})$} }
  child { node[align=center] {$B$ \\ $P(\bar{A}\cap B)$} edge from parent node[above,pos=0.6]{$P_{\bar{A}}(B)$} }
  edge from parent node[below,pos=0.6]{$P(\bar{A})$} }
child { node {$A$}
  child { node[align=center] {$\bar{B}$ \\ $P(A\cap\bar{B})$} edge from parent node[below,pos=0.6]{$P_A(\bar{B})$} }
  child { node[align=center] {$B$ \\ $P(A\cap B)$} edge from parent node[above,pos=0.6]{$P_A(B)$} }
  edge from parent node[above,pos=0.6]{$P(A)$} };
\end{tikzpicture}
\legende{Arbre pondéré à deux niveaux : chaque chemin donne une intersection (probabilité indiquée en bout de branche), par exemple $P(A\cap B)=P(A)\times P_A(B)$.}
\end{popfigure}

\courssub{Système complet et formule des probabilités totales}

\begin{definition}[Système complet d'événements]
Des événements $A_1, A_2, \ldots, A_n$ forment un \cle{système complet} (ou une partition de $\Omega$) s'ils sont deux à deux incompatibles ($A_i\cap A_j=\varnothing$ pour $i\neq j$) et si leur réunion est $\Omega$ ($\bigcup_{i=1}^n A_i=\Omega$). Le cas le plus utilisé est le couple $(A,\bar{A})$.
\end{definition}

\begin{propriete}[Formule des probabilités totales]
Si $(A_1,\ldots,A_n)$ est un système complet avec $P(A_i)\neq 0$ pour tout $i$, alors pour tout événement $B$ :
\[ P(B)=\sum_{i=1}^{n} P(A_i)\times P_{A_i}(B). \]
Dans le cas particulier $(A,\bar{A})$ :
\[ P(B)=P(A)\times P_A(B)+P(\bar{A})\times P_{\bar{A}}(B). \]
Sur l'arbre, c'est exactement la règle 2 : on additionne les probabilités de tous les chemins menant à $B$.
\end{propriete}

\begin{propriete}[Formule de Bayes (ou formule des causes)]
Dans les mêmes conditions, pour tout $k\in\{1,\ldots,n\}$ tel que $P(B)\neq 0$ :
\[ P_{B}(A_k)=\dfrac{P(A_k)\times P_{A_k}(B)}{\sum_{i=1}^{n} P(A_i)\times P_{A_i}(B)}. \]
Elle permet de « remonter l'arbre » : connaître la probabilité d'une cause $A_k$ sachant l'effet $B$ observé.
\end{propriete}

\begin{exoresolu}[Test de dépistage d'une maladie]
Dans une région, $2\,\%$ de la population est atteinte d'une maladie $M$. Un test de dépistage est positif chez $95\,\%$ des malades et chez $3\,\%$ des non-malades. On note $T$ : « le test est positif ». Calculer la probabilité qu'une personne choisie au hasard ait un test positif. Si le test est positif, quelle est la probabilité que la personne soit réellement malade ?
\tcblower
Le couple $(M,\bar{M})$ forme un système complet. Les données sont :
$P(M)=0{,}02$, $P_M(T)=0{,}95$, $P_{\bar{M}}(T)=0{,}03$.
D'après la formule des probabilités totales :
\begin{align*}
P(T) &= P(M)\times P_M(T)+P(\bar{M})\times P_{\bar{M}}(T)\\
&= 0{,}02\times 0{,}95 + 0{,}98\times 0{,}03\\
&= 0{,}019 + 0{,}0294 = 0{,}0484.
\end{align*}
La probabilité d'un test positif est donc $0{,}0484$, soit environ $4{,}84\,\%$.
\medskip
Pour la seconde question, on utilise la formule de Bayes :
\[ P_T(M)=\dfrac{P(M\cap T)}{P(T)}=\dfrac{P(M)\times P_M(T)}{P(T)}=\dfrac{0{,}019}{0{,}0484}\approx 0{,}3926. \]
\emph{Remarque} : malgré un test performant, la probabilité d'être malade sachant un test positif n'est que d'environ $39\,\%$ ! C'est le paradoxe bien connu des maladies rares (faible prévalence).
\end{exoresolu}

\begin{exemplebox}[Fiabilité d'une ligne de transmission]
Un message envoyé par un centre CAMTEL transite par deux relais. Le premier relais fonctionne avec probabilité $0{,}9$ ; s'il fonctionne, le message est transmis correctement avec probabilité $0{,}98$ ; s'il tombe en panne, le message est perdu. La probabilité de transmission correcte est
\[ P(T)=0{,}9\times 0{,}98 + 0{,}1\times 0 = 0{,}882. \]
\end{exemplebox}

\begin{savaistu}[Le paradoxe du dépistage et le filtrage anti-spam]
Le contre-exemple du test médical (maladie rare $\Rightarrow$ beaucoup de faux positifs) est le même mécanisme que les filtres anti-spam d'Orange Cameroun ou MTN. Si $1\,\%$ des mails sont des spams et que le filtre a $99\,\%$ de détection mais $5\,\%$ de faux positifs, un mail classé « spam » n'a que $\approx 17\,\%$ de chances d'être un vrai spam. C'est pourquoi les dossiers « indésirables » contiennent souvent des mails légitimes : le filtre privilégie la sécurité (ne pas rater un spam) au détriment de la précision.
\end{savaistu}

\begin{attention}[Pièges fréquents]
\begin{itemize}
\item La formule des probabilités totales exige un \emph{système complet} : vérifie que les événements de la partition sont incompatibles et recouvrent tout $\Omega$.
\item Sur un arbre, n'additionne jamais les probabilités le long d'un même chemin : on \emph{multiplie} le long d'un chemin, on \emph{additionne} entre chemins distincts.
\item $P_A(B)$ n'est définie que si $P(A)\neq 0$ : mentionne cette condition dans ta rédaction au Bac.
\item Pour Bayes, ne confonds pas $P_A(B)$ (donnée directe, souvent fiabilité d'un test) et $P_B(A)$ (cherchée, probabilité a posteriori).
\end{itemize}
\end{attention}

\begin{exercice}[Contrôle qualité à la SABC]
Une usine remplit des bouteilles avec deux machines : $M_1$ produit $60\,\%$ des bouteilles dont $1\,\%$ sont défectueuses ; $M_2$ produit le reste dont $2\,\%$ sont défectueuses. On prélève une bouteille au hasard. 
\begin{enumerate}
\item Construire l'arbre pondéré.
\item Calculer la probabilité qu'elle soit défectueuse.
\item Sachant qu'elle est défectueuse, quelle est la probabilité qu'elle provienne de $M_1$ ?
\end{enumerate}
\end{exercice}

\begin{corrige}[Exercice : contrôle qualité]
\begin{enumerate}
\item Premier niveau (machine) : $P(M_1)=0{,}6$, $P(M_2)=0{,}4$. Second niveau (défaut $D$) : $P_{M_1}(D)=0{,}01$, $P_{M_1}(\bar{D})=0{,}99$ ; $P_{M_2}(D)=0{,}02$, $P_{M_2}(\bar{D})=0{,}98$. Les feuilles de l'arbre portent $P(M_1\cap D)=0{,}006$, $P(M_1\cap\bar{D})=0{,}594$, $P(M_2\cap D)=0{,}008$, $P(M_2\cap\bar{D})=0{,}392$.
\item Probabilités totales : $P(D)=0{,}6\times 0{,}01+0{,}4\times 0{,}02=0{,}006+0{,}008=0{,}014$.
\item Formule de Bayes : $P_D(M_1)=\dfrac{P(M_1\cap D)}{P(D)}=\dfrac{0{,}006}{0{,}014}=\dfrac{3}{7}\approx 0{,}4286$.
\end{enumerate}
\end{corrige}

\begin{aretenir}[L'essentiel du complément]
\begin{itemize}
\item $P_A(B)=\dfrac{P(A\cap B)}{P(A)}$ pour $P(A)\neq 0$ ; donc $P(A\cap B)=P(A)\times P_A(B)$.
\item Sur un arbre pondéré : produit le long d'un chemin, somme entre les chemins.
\item Formule des probabilités totales : $P(B)=\sum P(A_i)P_{A_i}(B)$, $(A_i)$ système complet.
\item Formule de Bayes : $P_B(A_k)=\dfrac{P(A_k)P_{A_k}(B)}{\sum P(A_i)P_{A_i}(B)}$ (remonter l'arbre).
\end{itemize}
\end{aretenir}

\newpage

\coursec{Activité d'intégration}

\coursec{Leçon M11 : Probabilités, variables aléatoires et loi binomiale}

\courssub{Activité d'intégration : La tombola du quartier Deïdo}

\courssub{Objectifs de l'activité}

À l'issue de cette activité d'intégration, tu seras capable de :

\begin{itemize}
\item modéliser une situation concrète de la vie courante (une tombola de quartier) à l'aide d'une \cle{variable aléatoire} et de sa \cle{loi de probabilité} ;
\item calculer une \cle{espérance mathématique} et l'interpréter pour juger si un jeu est équitable, favorable ou défavorable ;
\item reconnaître et justifier un \cle{schéma de Bernoulli}, puis utiliser la \cle{loi binomiale} pour calculer une probabilité et un nombre moyen de succès ;
\item discuter de manière critique une hypothèse de modélisation (tirages avec ou sans remise) ;
\item rédiger une conclusion argumentée permettant de prendre une décision : conseiller ou déconseiller l'achat de billets de tombola.
\end{itemize}

\courssub{Situation}

\textbf{La tombola du quartier Deïdo, à Douala.}

Pour financer la réfection du terrain de football du quartier, l'association des jeunes de Deïdo organise une grande tombola lors de la fête de fin d'année. Voici les conditions du jeu, affichées sur la banderole devant le foyer culturel :

\begin{itemize}
\item \num{500} billets numérotés de 1 à 500 sont mis en vente, au prix de \num{500} FCFA l'unité ;
\item le tirage aura lieu en public, dans un sac opaque contenant les \num{500} numéros ; chaque numéro a la même chance d'être tiré ;
\item \textbf{10 billets gagnants} donnent droit à un lot d'une valeur de \num{5000} FCFA chacun (sacs de riz, bons d'achat chez le commerçant du coin) ;
\item \textbf{1 billet gagnant} donne droit au gros lot : \num{50000} FCFA en espèces ;
\item les 11 numéros gagnants sont tirés successivement et sont tous distincts : un même billet ne peut pas gagner deux lots.
\end{itemize}

Mama Ngo Bell, qui tient une boutique de provisions face au marché, hésite : son neveu, élève en Terminale C au lycée de Deïdo, lui propose d'acheter 20 billets « pour augmenter ses chances ». Elle se pose trois questions :

\begin{enumerate}
\item Si j'achète un seul billet, combien puis-je espérer gagner \emph{en moyenne} ? Le jeu est-il équitable ?
\item L'association est-elle sûre de gagner de l'argent, et combien ?
\item Si j'achète 20 billets, quelle est la probabilité que je gagne au moins un lot ? Combien de lots puis-je espérer remporter en moyenne ?
\end{enumerate}

Par groupes de quatre, vous allez aider Mama Ngo Bell à prendre sa décision, en mobilisant tous les outils de la leçon : équiprobabilité, loi d'une variable aléatoire, espérance, schéma de Bernoulli et loi binomiale.

\courssub{Consignes}

\textbf{Partie A — Le point de vue de l'acheteur d'un billet.}

\begin{enumerate}
\item On choisit au hasard un billet parmi les 500. Justifier que l'on est en situation d'\cle{équiprobabilité}.
\item On note $X$ la variable aléatoire égale au \emph{gain algébrique} de l'acheteur d'un billet, c'est-à-dire la valeur du lot reçu \emph{moins} le prix du billet. Déterminer les valeurs prises par $X$.
\item Déterminer la loi de probabilité de $X$ (on présentera les résultats dans un tableau). Vérifier que la somme des probabilités vaut 1.
\item Calculer l'espérance $E(X)$. Interpréter le résultat : que perd (ou gagne) en moyenne un joueur qui achèterait un billet à un très grand nombre de tombolas identiques ? Le jeu est-il équitable ?
\end{enumerate}

\textbf{Partie B — Le point de vue de l'association.}

\begin{enumerate}
\setcounter{enumi}{4}
\item Calculer la recette totale de la vente des billets et le coût total des lots distribués. En déduire le bénéfice de l'association. Ce bénéfice est-il certain ou aléatoire ? Justifier.
\item Retrouver le lien entre le bénéfice par billet de l'association et l'espérance $E(X)$ calculée à la question 4.
\end{enumerate}

\textbf{Partie C — Le commerçant qui achète 20 billets.}

Mama Ngo Bell achète finalement 20 billets. On s'intéresse au nombre $Y$ de lots qu'elle remporte.

\begin{enumerate}
\setcounter{enumi}{6}
\item Expliquer pourquoi, rigoureusement, les 20 billets ne constituent pas des épreuves indépendantes (les numéros gagnants sont tirés sans remise).
\item Pour simplifier, on assimile la situation à 20 tirages \emph{avec remise}, chaque billet ayant une probabilité $p = \dfrac{11}{500}$ d'être gagnant. Justifier alors que $Y$ suit une loi binomiale dont on précisera les paramètres. Pourquoi cette approximation est-elle raisonnable ici ?
\item Calculer la probabilité que Mama Ngo Bell ne gagne \emph{aucun} lot, puis la probabilité qu'elle gagne \emph{au moins un} lot. Arrondir à $10^{-3}$.
\item Calculer $E(Y)$ et interpréter : combien de lots peut-elle espérer remporter en moyenne ?
\item Sachant qu'elle a dépensé \num{10000} FCFA en billets, et en utilisant l'espérance de la Partie A, calculer son gain algébrique moyen total. L'achat de 20 billets améliore-t-il son espérance de gain ?
\end{enumerate}

\textbf{Partie D — Conclusion argumentée.}

\begin{enumerate}
\setcounter{enumi}{11}
\item Rédiger en quelques lignes un conseil à Mama Ngo Bell : faut-il acheter des billets de tombola ? Votre argumentaire devra distinguer le point de vue financier (espérance de gain) et le point de vue solidaire (financement du terrain de football).
\end{enumerate}

\courssub{Ressources à mobiliser}

\begin{aretenir}[Boîte à outils de la leçon]
\begin{itemize}
\item \textbf{Équiprobabilité} : si l'univers $\Omega$ contient $n$ issues équiprobables, alors pour tout événement $A$ : $P(A) = \dfrac{\text{card}(A)}{\text{card}(\Omega)} = \dfrac{\text{nombre de cas favorables}}{\text{nombre de cas possibles}}$.
\item \textbf{Loi d'une variable aléatoire} $X$ prenant les valeurs $x_1, x_2, \dots, x_k$ : c'est la donnée des probabilités $p_i = P(X = x_i)$, avec $p_1 + p_2 + \dots + p_k = 1$.
\item \textbf{Fonction de répartition} : $F(x) = P(X \le x)$. C'est une fonction croissante, continue à droite, avec des sauts aux valeurs prises par $X$.
\item \textbf{Espérance} : $E(X) = \displaystyle\sum_{i=1}^{k} p_i\, x_i$. Le jeu est \emph{équitable} si $E(X) = 0$, favorable au joueur si $E(X) > 0$, défavorable si $E(X) < 0$.
\item \textbf{Variance} : $V(X) = E\bigl[(X - E(X))^2\bigr] = E(X^2) - [E(X)]^2 = \displaystyle\sum_{i=1}^{k} p_i\, (x_i - E(X))^2$.
\item \textbf{Écart-type} : $\sigma(X) = \sqrt{V(X)}$. Il mesure la dispersion des valeurs autour de l'espérance (même unité que $X$).
\item \textbf{Schéma de Bernoulli} : $n$ répétitions \emph{identiques et indépendantes} d'une épreuve de Bernoulli de probabilité de succès $p$. La variable $Y$ qui compte les succès suit la loi binomiale $\mathcal{B}(n, p)$.
\item \textbf{Loi binomiale} : $P(Y = k) = \dbinom{n}{k} p^k (1-p)^{n-k}$, avec $E(Y) = np$ et $V(Y) = np(1-p)$.
\item \textbf{Événement contraire} : $P(Y \geq 1) = 1 - P(Y = 0)$, souvent bien plus rapide à calculer.
\end{itemize}
\end{aretenir}

\courssub{Démarche et corrigé commenté}

\begin{exoresolu}[La tombola du quartier Deïdo — Partie A : Point de vue de l'acheteur]
\textbf{1. Équiprobabilité.} \textbf{2. Valeurs de $X$.} \textbf{3. Loi de $X$.} \textbf{4. Espérance et interprétation.}
\tcblower
\textbf{1. Équiprobabilité.} Le numéro gagnant est tiré au hasard dans un sac opaque contenant les 500 numéros : chaque numéro a donc la même probabilité $\dfrac{1}{500}$ d'être tiré. On est bien en situation d'équiprobabilité, avec $\text{card}(\Omega) = 500$.

\textbf{2. Valeurs de $X$.} Trois cas se présentent pour l'acheteur d'un billet à 500 FCFA :
\begin{itemize}
\item il gagne le gros lot : $X = \num{50000} - 500 = \num{49500}$ ;
\item il gagne un lot de \num{5000} FCFA : $X = \num{5000} - 500 = \num{4500}$ ;
\item il ne gagne rien : $X = 0 - 500 = -500$.
\end{itemize}
Donc $X$ prend les valeurs $\num{49500}$, $\num{4500}$ et $-500$.

\textbf{3. Loi de probabilité de $X$.} Par équiprobabilité :
\[
P(X = \num{49500}) = \frac{1}{500}, \qquad P(X = \num{4500}) = \frac{10}{500}, \qquad P(X = -500) = \frac{489}{500}.
\]
\begin{center}
\begin{tabularx}{\linewidth}{|l|X|X|X|}
\hline
\rowcolor{popblueL} Valeur $x_i$ & \centering $\num{49500}$ & \centering $\num{4500}$ & \multicolumn{1}{c|}{$-500$} \\
\hline
$P(X = x_i)$ & \centering $\dfrac{1}{500}$ & \centering $\dfrac{10}{500}$ & \multicolumn{1}{c|}{$\dfrac{489}{500}$} \\
\hline
\end{tabularx}
\end{center}
Vérification : $\dfrac{1}{500} + \dfrac{10}{500} + \dfrac{489}{500} = \dfrac{500}{500} = 1$. La loi est bien définie.

\textbf{4. Espérance de $X$.}
\begin{align*}
E(X) &= \frac{1}{500}\times \num{49500} + \frac{10}{500}\times \num{4500} + \frac{489}{500}\times(-500)\\
&= \frac{\num{49500} + \num{45000} - \num{244500}}{500} = \frac{-\num{150000}}{500} = -300.
\end{align*}
\emph{Interprétation} : en moyenne, l'acheteur d'un billet \textbf{perd 300 FCFA} par billet. Comme $E(X) < 0$, le jeu est \textbf{défavorable au joueur} : il n'est pas équitable. Cela ne signifie pas que tout joueur perd : un joueur peut gagner le gros lot ; mais sur un très grand nombre de participations, la perte moyenne se rapproche de 300 FCFA par billet (loi des grands nombres).
\end{exoresolu}

\begin{exoresolu}[Partie B — Le point de vue de l'association]
\textbf{5. Recettes, coûts et bénéfice.} \textbf{6. Lien avec $E(X)$.}
\tcblower
\textbf{5. Bénéfice de l'association.}
\begin{itemize}
\item Recette de la vente : $500 \times 500 = \num{250000}$ FCFA.
\item Coût total des lots : $10 \times \num{5000} + 1 \times \num{50000} = \num{50000} + \num{50000} = \num{100000}$ FCFA.
\end{itemize}
\[
\text{Bénéfice} = \num{250000} - \num{100000} = \num{150000} \text{ FCFA}.
\]
Ce bénéfice est \textbf{certain} (et non aléatoire) : dès lors que les 500 billets sont vendus, la recette est fixe, et les 11 lots seront distribués quoiqu'il arrive, pour un coût connu d'avance. Le hasard ne joue que sur \emph{qui} gagne, pas sur les montants. L'association financera donc le terrain à hauteur de \num{150000} FCFA.

\textbf{6. Lien avec l'espérance.} Le bénéfice par billet de l'association est $\dfrac{\num{150000}}{500} = 300$ FCFA : c'est exactement l'\emph{opposé} de $E(X) = -300$. C'est logique : le gain de l'organisateur est le contraire du gain algébrique du joueur, donc $E(\text{gain association}) = -E(X) = 300$ FCFA par billet. Ce que les joueurs perdent en moyenne, l'association le gagne.
\end{exoresolu}

\begin{exoresolu}[Partie C — Les 20 billets de Mama Ngo Bell]
\textbf{7. Non-indépendance stricte.} \textbf{8. Modèle binomial approché.} \textbf{9. $P(Y \geq 1)$.} \textbf{10. $E(Y)$.} \textbf{11. Gain moyen total.}
\tcblower
\textbf{7. Non-indépendance stricte.} Les 11 numéros gagnants sont tirés \emph{sans remise} et sont distincts : si l'un des billets de Mama Ngo Bell est déjà gagnant, il reste moins de lots disponibles pour ses autres billets. La probabilité qu'un deuxième billet gagne dépend du résultat du premier : les épreuves ne sont donc \emph{pas rigoureusement indépendantes}. Le modèle exact serait plus compliqué (loi hypergéométrique, hors programme).

\textbf{8. Modèle binomial approché.} On assimile la situation à 20 tirages \emph{avec remise} : chaque billet est alors une épreuve de Bernoulli de succès « le billet gagne un lot », de probabilité
\[
p = \frac{11}{500} = \num{0,022},
\]
et les 20 épreuves sont identiques et indépendantes. C'est donc un \cle{schéma de Bernoulli} : $Y$ suit la loi binomiale $\mathcal{B}\left(20\,;\,\num{0,022}\right)$. L'approximation est raisonnable car 20 billets représentent une petite fraction (4 \%) des 500 billets : le fait qu'un billet soit gagnant modifie très peu les chances des autres.

\textbf{9. Probabilité de gagner au moins un lot.} On passe par l'événement contraire :
\[
P(Y = 0) = \binom{20}{0}\,(\num{0,022})^0\,(\num{0,978})^{20} = (\num{0,978})^{20} \approx \num{0,641}.
\]
\[
P(Y \geq 1) = 1 - P(Y = 0) \approx 1 - \num{0,641} = \num{0,359}.
\]
Mama Ngo Bell a donc environ \textbf{36 \% de chances} de remporter au moins un lot : c'est loin d'être sûr, malgré ses 20 billets.

\textbf{10. Nombre moyen de lots.}
\[
E(Y) = np = 20 \times \frac{11}{500} = \frac{220}{500} = \num{0,44}.
\]
En moyenne, elle remporte \num{0,44} lot : autrement dit, moins d'un lot sur deux participations de ce type.

\textbf{11. Gain algébrique moyen total.} Chaque billet a une espérance de gain de $-300$ FCFA ; par additivité de l'espérance, pour 20 billets :
\[
E(\text{gain total}) = 20 \times (-300) = -\num{6000} \text{ FCFA}.
\]
Sur ses \num{10000} FCFA investis, elle peut espérer récupérer en moyenne environ \num{4000} FCFA seulement. Acheter plus de billets augmente la probabilité de gagner \emph{au moins une fois}, mais \textbf{aggrave la perte moyenne} : l'espérance de gain par franc dépensé reste la même.
\end{exoresolu}

\begin{exoresolu}[Partie D — Conclusion argumentée]
\textbf{12) Rédiger un conseil à Mama Ngo Bell.}
\tcblower
\textbf{Exemple de conclusion rédigée.}

« Du strict point de vue financier, la tombola est défavorable à l'acheteur : chaque billet de 500 FCFA rapporte en moyenne 200 FCFA de lots, soit une perte moyenne de 300 FCFA par billet, et acheter 20 billets porte la perte moyenne à \num{6000} FCFA, pour une probabilité de gagner au moins un lot d'environ 36 \% seulement. Si l'objectif de Mama Ngo Bell est de gagner de l'argent, il faut donc lui \emph{déconseiller} cet achat : aucune stratégie d'achat ne peut rendre l'espérance positive. En revanche, la tombola n'est pas une arnaque : la règle du jeu est transparente, le bénéfice certain de \num{150000} FCFA finance le terrain de football du quartier, et l'achat de billets peut être vu comme un \emph{don partiellement remboursé en moyenne}, avec une petite chance de gros lot. Notre conseil : acheter un ou deux billets par solidarité avec l'association, en sachant que l'on « donne » en moyenne 300 FCFA par billet — mais ne pas considérer la tombola comme un placement. »
\end{exoresolu}

\begin{attention}[Pièges fréquents dans ce type de problème]
\begin{itemize}
\item \textbf{Oublier le prix du billet} : le gain algébrique est le lot \emph{moins} 500 FCFA. Calculer l'espérance avec les lots bruts donnerait $+200$ FCFA et ferait croire à tort que le jeu est favorable !
\item \textbf{Confondre « probable » et « rentable »} : avoir 36 \% de chances de gagner au moins un lot ne signifie pas que l'opération est rentable ; un lot de \num{5000} FCFA ne rembourse que la moitié des \num{10000} FCFA dépensés.
\item \textbf{Utiliser la loi binomiale sans vérifier les conditions} : il faut des épreuves \emph{identiques et indépendantes}. Ici, l'indépendance n'est qu'approchée (tirages sans remise assimilés à des tirages avec remise) : le dire explicitement fait partie de la qualité de la modélisation.
\item \textbf{Oublier de vérifier} que la somme des probabilités d'une loi vaut 1.
\end{itemize}
\end{attention}

\courssub{Bilan de l'activité}

Cette activité a permis de mettre en évidence plusieurs idées essentielles de la leçon :

\begin{itemize}
\item L'\textbf{équiprobabilité} fournit naturellement la loi de probabilité d'une variable aléatoire dès que l'on sait dénombrer les cas favorables (1 gros lot, 10 lots, 489 billets perdants sur 500).
\item L'\textbf{espérance mathématique} est l'outil de décision central : elle résume en un seul nombre le comportement moyen d'un jeu et permet de trancher objectivement la question « le jeu est-il équitable ? ». Ici, $E(X) = -300$ FCFA tranche sans appel.
\item Le gain de l'organisateur est l'\textbf{opposé} du gain des joueurs : une tombola associative est conçue pour être défavorable aux joueurs, et c'est précisément ce qui garantit le bénéfice certain de l'association.
\item Le \textbf{schéma de Bernoulli} et la \textbf{loi binomiale} permettent de traiter les achats multiples : probabilité de gagner au moins une fois (via l'événement contraire) et nombre moyen de succès $E(Y) = np$.
\item Une modélisation exige de \textbf{discuter ses hypothèses} : l'assimilation de tirages sans remise à des tirages avec remise est une approximation, légitime ici car $20$ est petit devant $500$.
\item Enfin, les mathématiques éclairent la décision sans la dicter : le calcul montre que la tombola est un mauvais \emph{investissement} mais peut être un bon \emph{geste solidaire}. Savoir distinguer ces deux points de vue, c'est exactement l'esprit de l'approche par les compétences : mobiliser des savoirs pour résoudre des situations de la vie courante et argumenter un choix.
\end{itemize}

\begin{savaistu}[Pour aller plus loin]
Les loteries nationales et les jeux de hasard fonctionnent tous sur ce principe : l'espérance de gain du joueur est toujours strictement négative, la différence finançant l'organisateur et les lots. Au Cameroun, les jeux de paris sportifs très populaires chez les jeunes obéissent à la même logique mathématique : sur le long terme, c'est l'opérateur qui gagne. Comprendre l'espérance mathématique, c'est se donner les moyens de ne pas se laisser piéger par le mirage du « gros lot ».
\end{savaistu}

\newpage

\coursec{Méthodes et exercices résolus}

\coursec{Probabilités, variables aléatoires et loi binomiale}

\courssub{Vocabulaire des probabilités et équiprobabilité}

\begin{definition}[Expérience aléatoire, Univers et Événements]
Une \cle{expérience aléatoire} est une expérience dont le résultat ne peut être prédit avec certitude avant sa réalisation, bien que l'ensemble des résultats possibles soit connu.
L'\cle{univers} $\Omega$ est l'ensemble de tous les résultats possibles (issues élémentaires).
Un \cle{événement} est une partie de $\Omega$ (un sous-ensemble des issues).
Un \cle{événement élémentaire} est un singleton $\{ \omega \}$ ne contenant qu'une seule issue.
L'\cle{événement contraire} de $A$ (noté $\bar{A}$ ou $A^c$) est l'ensemble des issues de $\Omega$ qui ne sont pas dans $A$ : $\bar{A} = \Omega \setminus A$.
Deux événements $A$ et $B$ sont \cle{incompatibles} (ou mutuellement exclusifs) si $A \cap B = \emptyset$ : ils ne peuvent pas se réaliser simultanément.
\end{definition}

\begin{definition}[Équiprobabilité]
On dit qu'il y a \cle{équiprobabilité} lorsque toutes les issues élémentaires de $\Omega$ ont la même chance de se réaliser.
C'est le cas typiquement pour : un dé parfait (6 faces), une pièce de monnaie non truquée (2 faces), un tirage au hasard dans une urne où les objets sont indiscernables au toucher, un tirage au sort équitable.
Si $\Omega$ a $N$ issues équiprobables, alors $P(\{\omega\}) = \frac{1}{N}$ pour tout $\omega \in \Omega$.
\end{definition}

\begin{propriete}[Probabilité dans le cas d'équiprobabilité]
Si $\Omega$ est fini et muni d'une équiprobabilité, pour tout événement $A \subset \Omega$ :
\[ P(A) = \frac{\text{Card}(A)}{\text{Card}(\Omega)} = \frac{\text{nombre de cas favorables}}{\text{nombre de cas possibles}}. \]
On a toujours $0 \leq P(A) \leq 1$, $P(\emptyset)=0$, $P(\Omega)=1$ et $P(\bar{A}) = 1 - P(A)$.
\end{propriete}

\begin{methode}[Calculer une probabilité dans un cas d'équiprobabilité]
Pour calculer la probabilité d'un événement $A$ lorsque toutes les issues sont équiprobables :
\begin{enumerate}
\item \textbf{Justifier l'équiprobabilité} : décrire l'expérience et préciser pourquoi les issues ont toutes la même chance de se réaliser (pièce non truquée, dé équilibré, tirage au hasard dans une urne...).
\item \textbf{Déterminer l'univers} $\Omega$ et calculer son cardinal $\text{Card}(\Omega)$, souvent à l'aide d'un dénombrement (arrangements, combinaisons, $p$-listes).
\item \textbf{Dénombrer les issues favorables} à $A$, c'est-à-dire calculer $\text{Card}(A)$.
\item \textbf{Conclure} avec la formule :
\[ P(A) = \frac{\text{Card}(A)}{\text{Card}(\Omega)} = \frac{\text{nombre de cas favorables}}{\text{nombre de cas possibles}}. \]
\item Simplifier la fraction et, si demandé, donner une valeur décimale ou un pourcentage.
\end{enumerate}
\textbf{Réflexes} : un tirage \emph{simultané} (sans ordre) se dénombre avec les combinaisons $\binom{n}{p}$ ; un tirage \emph{successif sans remise} avec les arrangements $A_n^p$ ; un tirage \emph{successif avec remise} avec les $p$-listes $n^p$.
\end{methode}

\begin{exoresolu}[Tirage simultané dans une urne]
Une urne contient 5 boules rouges et 4 boules vertes, indiscernables au toucher. On tire simultanément 3 boules de l'urne.
\begin{enumerate}
\item Déterminer le nombre de tirages possibles.
\item Calculer la probabilité de l'événement $A$ : « obtenir exactement 2 boules rouges ».
\item Calculer la probabilité de l'événement $B$ : « obtenir au moins une boule verte ».
\end{enumerate}
\tcblower
\textbf{1. Nombre de tirages possibles.}

Le tirage est simultané : l'ordre n'intervient pas et il n'y a pas de remise. On choisit donc 3 boules parmi 9 :
\[ \text{Card}(\Omega) = \binom{9}{3} = \frac{9 \times 8 \times 7}{3 \times 2 \times 1} = 84. \]
Tous les tirages sont équiprobables car les boules sont indiscernables au toucher.

\textbf{2. Probabilité de $A$.}

Obtenir exactement 2 boules rouges signifie : choisir 2 boules parmi les 5 rouges \emph{et} 1 boule parmi les 4 vertes. D'après le principe multiplicatif :
\[ \text{Card}(A) = \binom{5}{2} \times \binom{4}{1} = 10 \times 4 = 40. \]
Par équiprobabilité :
\[ P(A) = \frac{\text{Card}(A)}{\text{Card}(\Omega)} = \frac{40}{84} = \frac{10}{21}. \]

\textbf{3. Probabilité de $B$.}

L'événement contraire $\bar{B}$ est « n'obtenir aucune boule verte », c'est-à-dire obtenir 3 boules rouges :
\[ \text{Card}(\bar{B}) = \binom{5}{3} = 10 \quad \text{donc} \quad P(\bar{B}) = \frac{10}{84} = \frac{5}{42}. \]
Alors :
\[ P(B) = 1 - P(\bar{B}) = 1 - \frac{5}{42} = \frac{37}{42}. \]

\textbf{Conclusion} : $P(A) = \dfrac{10}{21}$ et $P(B) = \dfrac{37}{42}$. Le passage par l'événement contraire a évité un dénombrement pénible en trois cas.
\end{exoresolu}

\courssub{Opérations sur les événements et arbres de probabilités}

\begin{definition}[Réunion, Intersection, Contraire]
Soient $A$ et $B$ deux événements de $\Omega$.
\begin{itemize}
\item La \cle{réunion} $A \cup B$ (se lit « $A$ \emph{ou} $B$ ») est l'événement « $A$ se réalise \textbf{ou} $B$ se réalise (ou les deux) ».
\item L'\cle{intersection} $A \cap B$ (se lit « $A$ \emph{et} $B$ ») est l'événement « $A$ se réalise \textbf{et} $B$ se réalise ».
\item L'\cle{événement contraire} $\bar{A}$ (ou $A^c$) est l'événement « $A$ ne se réalise pas ».
\end{itemize}
\end{definition}

\begin{propriete}[Formules fondamentales sur les probabilités]
Pour tous événements $A, B$ de $\Omega$ :
\begin{itemize}
\item $P(\bar{A}) = 1 - P(A)$.
\item $P(A \cup B) = P(A) + P(B) - P(A \cap B)$.
\item Si $A$ et $B$ sont incompatibles ($A \cap B = \emptyset$), alors $P(A \cup B) = P(A) + P(B)$.
\item Probabilité conditionnelle : si $P(B) > 0$, $P_A(B) = \dfrac{P(A \cap B)}{P(B)}$.
\item $A$ et $B$ sont \cle{indépendants} si et seulement si $P(A \cap B) = P(A) \times P(B)$.
\end{itemize}
\end{propriete}

\begin{methode}[Utiliser la formule de la réunion et l'événement contraire]
\begin{enumerate}
\item \textbf{Traduire l'énoncé} : identifier les événements $A$ et $B$, et exprimer ce qui est demandé ($A \cup B$ se lit « $A$ \emph{ou} $B$ » ; $A \cap B$ se lit « $A$ \emph{et} $B$ » ; $\bar{A}$ se lit « non $A$ »).
\item \textbf{Événement contraire} : utiliser
\[ P(\bar{A}) = 1 - P(A). \]
C'est souvent le réflexe gagnant dès qu'apparaît « au moins un ».
\item \textbf{Réunion} : appliquer la formule
\[ P(A \cup B) = P(A) + P(B) - P(A \cap B). \]
\item \textbf{Cas particulier} : si $A$ et $B$ sont \cle{incompatibles} ($A \cap B = \emptyset$), alors $P(A \cap B) = 0$ et la formule devient $P(A \cup B) = P(A) + P(B)$.
\item \textbf{Vérifier} que le résultat est compris entre 0 et 1.
\end{enumerate}
\end{methode}

\begin{exoresolu}[Enquête dans un lycée de Yaoundé]
Dans un lycée, $60\,\%$ des élèves pratiquent le football, $35\,\%$ pratiquent la musique, et $20\,\%$ pratiquent les deux. On interroge un élève au hasard. On note $F$ : « l'élève pratique le football » et $M$ : « l'élève pratique la musique ».
\begin{enumerate}
\item Calculer la probabilité que l'élève pratique le football ou la musique.
\item Calculer la probabilité qu'il ne pratique ni l'un ni l'autre.
\item Calculer la probabilité qu'il pratique le football mais pas la musique.
\end{enumerate}
\tcblower
On a $P(F) = 0{,}6$, $P(M) = 0{,}35$ et $P(F \cap M) = 0{,}2$.

\textbf{1.} D'après la formule de la réunion :
\[ P(F \cup M) = P(F) + P(M) - P(F \cap M) = 0{,}6 + 0{,}35 - 0{,}2 = 0{,}75. \]

\textbf{2.} « Ni l'un ni l'autre » est l'événement contraire de $F \cup M$, c'est-à-dire $\overline{F \cup M}$ :
\[ P(\overline{F \cup M}) = 1 - P(F \cup M) = 1 - 0{,}75 = 0{,}25. \]

\textbf{3.} « Football mais pas musique » est l'événement $F \cap \bar{M}$. Comme $F = (F \cap M) \cup (F \cap \bar{M})$ et que ces deux événements sont incompatibles :
\[ P(F) = P(F \cap M) + P(F \cap \bar{M}) \implies P(F \cap \bar{M}) = 0{,}6 - 0{,}2 = 0{,}4. \]

\textbf{Conclusion} : la probabilité de pratiquer au moins une des deux activités est $0{,}75$ ; celle de n'en pratiquer aucune est $0{,}25$ ; celle de pratiquer uniquement le football est $0{,}4$.
\end{exoresolu}

\begin{methode}[Construire et exploiter un arbre de probabilités]
Un arbre de probabilités modélise une expérience à plusieurs étapes successives.
\begin{enumerate}
\item \textbf{Identifier les étapes} de l'expérience (ex: premier tirage, second tirage).
\item \textbf{Dessiner les branches} pour chaque issue possible à chaque étape.
\item \textbf{Étiqueter chaque branche} par la probabilité \emph{conditionnelle} correspondante (probabilité de l'issue sachant le chemin parcouru).
\item \textbf{Calculer la probabilité d'un chemin} (issue composée) en multipliant les probabilités sur les branches de ce chemin (formule des probabilités composées : $P(A \cap B) = P(A) \times P_A(B)$).
\item \textbf{Calculer la probabilité d'un événement} en sommant les probabilités des chemins qui le réalisent.
\end{enumerate}
\end{methode}

\begin{exoresolu}[Arbre de probabilités : contrôle qualité sans remise]
Une boîte contient 10 composants dont 3 sont défectueux. On prélève 2 composants successivement \textbf{sans remise}. On note $D_i$ : « le $i$-ème composant est défectueux ».
\begin{enumerate}
\item Construire l'arbre de probabilités.
\item Calculer la probabilité que les deux composants soient défectueux.
\item Calculer la probabilité qu'exactement un des deux soit défectueux.
\end{enumerate}
\tcblower
\textbf{1. Arbre de probabilités.}

Première étape : $P(D_1) = \frac{3}{10} = 0{,}3$, $P(\bar{D_1}) = \frac{7}{10} = 0{,}7$.
Seconde étape (sans remise, les probabilités changent) :
\begin{itemize}
\item Si $D_1$ réalisé : il reste 9 composants dont 2 défectueux. $P_{D_1}(D_2) = \frac{2}{9}$, $P_{D_1}(\bar{D_2}) = \frac{7}{9}$.
\item Si $\bar{D_1}$ réalisé : il reste 9 composants dont 3 défectueux. $P_{\bar{D_1}}(D_2) = \frac{3}{9} = \frac{1}{3}$, $P_{\bar{D_1}}(\bar{D_2}) = \frac{6}{9} = \frac{2}{3}$.
\end{itemize}

\begin{popfigure}
\begin{tikzpicture}[grow=right, level distance=3cm,
level 1/.style={sibling distance=3cm},
level 2/.style={sibling distance=1.5cm},
every node/.style={font=\small}]
\node {}
child { node {$\bar{D_1}$}
  child { node {$\bar{D_2}$} edge from parent node[below, popmuted] {\scriptsize $\frac{2}{3}$} }
  child { node {$D_2$} edge from parent node[above, popmuted] {\scriptsize $\frac{1}{3}$} }
  edge from parent node[below, popmuted] {\scriptsize $0{,}7$} }
child { node {$D_1$}
  child { node {$\bar{D_2}$} edge from parent node[below, popmuted] {\scriptsize $\frac{7}{9}$} }
  child { node {$D_2$} edge from parent node[above, popmuted] {\scriptsize $\frac{2}{9}$} }
  edge from parent node[above, popmuted] {\scriptsize $0{,}3$} };
\end{tikzpicture}
\legende{Arbre de probabilités pour un tirage sans remise de 2 composants.}
\end{popfigure}

\textbf{2. Deux défectueux :} chemin $D_1 \cap D_2$.
\[ P(D_1 \cap D_2) = P(D_1) \times P_{D_1}(D_2) = 0{,}3 \times \frac{2}{9} = \frac{0{,}6}{9} = \frac{1}{15} \approx 0{,}0667. \]

\textbf{3. Exactement un défectueux :} chemins $D_1 \cap \bar{D_2}$ et $\bar{D_1} \cap D_2$.
\begin{align*}
P(D_1 \cap \bar{D_2}) &= 0{,}3 \times \frac{7}{9} = \frac{2{,}1}{9} = \frac{7}{30} ; \\
P(\bar{D_1} \cap D_2) &= 0{,}7 \times \frac{1}{3} = \frac{7}{30}.
\end{align*}
\[ P(\text{exactement 1}) = \frac{7}{30} + \frac{7}{30} = \frac{7}{15} \approx 0{,}4667. \]

\textbf{Conclusion} : sans remise, les tirages ne sont pas indépendants ($P(D_2)$ dépend de $D_1$), l'arbre est l'outil indispensable.
\end{exoresolu}

\courssub{Variables aléatoires et loi de probabilité}

\begin{definition}[Variable aléatoire]
Une \cle{variable aléatoire} (v.a.) est une application $X : \Omega \to \mathbb{R}$ qui associe un nombre réel à chaque issue de l'univers.
On note $X(\Omega) = \{x_1, x_2, \ldots, x_n\}$ l'ensemble des valeurs prises par $X$.
Pour chaque valeur $x_i$, l'événement « $X = x_i \}$ » est $\{ \omega \in \Omega \mid X(\omega) = x_i \}$.
\end{definition}

\begin{definition}[Loi de probabilité d'une variable aléatoire]
La \cle{loi de probabilité} d'une v.a. $X$ à valeurs dans $\{x_1, \ldots, x_n\}$ est le tableau qui associe à chaque valeur $x_i$ sa probabilité $p_i = P(X = x_i)$.
\begin{center}
\begin{tabularx}{\linewidth}{|l|*{4}{>{\centering\arraybackslash}X|}}
\hline
\rowcolor{popblueL} $x_i$ & $x_1$ & $x_2$ & $\ldots$ & $x_n$ \\
\hline
$P(X = x_i)$ & $p_1$ & $p_2$ & $\ldots$ & $p_n$ \\
\hline
\end{tabularx}
\end{center}
Les $p_i$ satisfont $p_i \geq 0$ et $\sum_{i=1}^n p_i = 1$.
\end{definition}

\begin{definition}[Fonction de répartition]
La \cle{fonction de répartition} de $X$ est la fonction $F : \mathbb{R} \to [0,1]$ définie par $F(x) = P(X \leq x)$.
C'est une fonction croissante, continue à droite, à paliers, telle que $\lim_{x \to -\infty} F(x) = 0$ et $\lim_{x \to +\infty} F(x) = 1$.
Pour $x \in [x_k, x_{k+1})$, $F(x) = \sum_{i=1}^k p_i$.
\end{definition}

\begin{propriete}[Propriétés de la fonction de répartition]
Pour tout réel $a \leq b$ :
\[ P(a < X \leq b) = F(b) - F(a). \]
En particulier, pour une valeur $x_i$ prise par $X$ :
\[ P(X = x_i) = F(x_i) - \lim_{x \to x_i^-} F(x) \quad \text{(hauteur du palier)}. \]
\end{propriete}

\begin{methode}[Établir la loi de probabilité d'une variable aléatoire $X$]
\begin{enumerate}
\item \textbf{Déterminer l'ensemble des valeurs} prises par $X$ : $X(\Omega) = \{x_1, x_2, \ldots, x_n\}$. C'est l'étape la plus importante : lister \emph{toutes} les valeurs possibles.
\item Pour \textbf{chaque valeur} $x_i$, identifier l'événement $(X = x_i)$ et calculer $p_i = P(X = x_i)$, souvent par équiprobabilité ou à l'aide d'un arbre.
\item \textbf{Présenter la loi} dans un tableau.
\item \textbf{Contrôle obligatoire} : vérifier que $p_1 + p_2 + \cdots + p_n = 1$.
\item Si on demande la \cle{fonction de répartition} $F$, calculer les sommes cumulées : $F(x) = P(X \leq x)$, définie par paliers sur $\mathbb{R}$.
\end{enumerate}
\end{methode}

\begin{exoresolu}[Gains à une tombola de quartier]
Lors d'une tombola organisée à Douala, un sac contient 10 tickets : 1 ticket rapporte \SI{5000}{F}, 3 tickets rapportent \SI{1000}{F} et 6 tickets ne rapportent rien. On tire un ticket au hasard et on note $X$ le gain (en francs).
\begin{enumerate}
\item Déterminer la loi de probabilité de $X$.
\item Déterminer la fonction de répartition de $X$.
\end{enumerate}
\tcblower
\textbf{1. Loi de probabilité de $X$.}

Les valeurs possibles de $X$ sont : $0$, $1000$ et $5000$. Les tickets étant indiscernables, il y a équiprobabilité sur les 10 tickets :
\begin{align*}
P(X = 0) &= \frac{6}{10} = 0{,}6 ; \\
P(X = 1000) &= \frac{3}{10} = 0{,}3 ; \\
P(X = 5000) &= \frac{1}{10} = 0{,}1.
\end{align*}
D'où la loi de $X$ :
\begin{center}
\begin{tabularx}{\linewidth}{|l|*{3}{>{\centering\arraybackslash}X|}}
\hline
\rowcolor{popblueL} $x_i$ & $0$ & $1000$ & $5000$ \\
\hline
$P(X = x_i)$ & $0{,}6$ & $0{,}3$ & $0{,}1$ \\
\hline
\end{tabularx}
\end{center}
\textbf{Contrôle} : $0{,}6 + 0{,}3 + 0{,}1 = 1$. La loi est bien définie.

\textbf{2. Fonction de répartition.}

$F(x) = P(X \leq x)$ est définie sur $\mathbb{R}$ par paliers :
\[ F(x) = \begin{cases} 0 & \text{si } x < 0 \\ 0{,}6 & \text{si } 0 \leq x < 1000 \\ 0{,}9 & \text{si } 1000 \leq x < 5000 \\ 1 & \text{si } x \geq 5000 \end{cases} \]
Par exemple, pour $1000 \leq x < 5000$ : $F(x) = P(X = 0) + P(X = 1000) = 0{,}6 + 0{,}3 = 0{,}9$.

\textbf{Conclusion} : $X$ prend les valeurs $0$, $1000$, $5000$ avec les probabilités $0{,}6$ ; $0{,}3$ ; $0{,}1$, et sa fonction de répartition est une fonction en escalier, croissante, de 0 à 1.
\end{exoresolu}

\courssub{Espérance, variance et écart-type}

\begin{definition}[Espérance mathématique]
L'\cle{espérance} (ou moyenne) d'une variable aléatoire $X$ de loi $(x_i, p_i)$ est la moyenne pondérée de ses valeurs par leurs probabilités :
\[ E(X) = \sum_{i=1}^{n} x_i \, p_i = x_1 p_1 + x_2 p_2 + \cdots + x_n p_n. \]
C'est la valeur moyenne que l'on attendrait en répétant l'expérience un très grand nombre de fois (Loi des grands nombres).
\end{definition}

\begin{definition}[Variance et écart-type]
La \cle{variance} mesure la dispersion des valeurs de $X$ autour de son espérance :
\[ V(X) = E\left( [X - E(X)]^2 \right) = \sum_{i=1}^{n} (x_i - E(X))^2 p_i. \]
L'\cle{écart-type} est la racine carrée de la variance : $\sigma(X) = \sqrt{V(X)}$. Il s'exprime dans la même unité que $X$.
\end{definition}

\begin{propriete}[Formule de König-Huygens et linéarité]
\begin{itemize}
\item \textbf{König-Huygens} (à connaître par cœur, utilisable sans démonstration au Bac) :
\[ V(X) = E(X^2) - [E(X)]^2 = \left(\sum_{i=1}^{n} x_i^2 \, p_i\right) - [E(X)]^2. \]
\item \textbf{Linéarité de l'espérance} : pour tous réels $a, b$, $E(aX + b) = a E(X) + b$.
\item \textbf{Variance affine} : $V(aX + b) = a^2 V(X)$.
\end{itemize}
\end{propriete}

\begin{methode}[Calculer $E(X)$, $V(X)$ et $\sigma(X)$]
Soit $X$ une variable aléatoire de loi $(x_i, p_i)$.
\begin{enumerate}
\item \textbf{Espérance} : multiplier chaque valeur par sa probabilité puis additionner :
\[ E(X) = \sum_{i=1}^{n} x_i \, p_i. \]
\item \textbf{Variance} : le plus rapide est la formule de König-Huygens :
\[ V(X) = E(X^2) - [E(X)]^2 = \left(\sum_{i=1}^{n} x_i^2 \, p_i\right) - [E(X)]^2. \]
On calcule donc d'abord $E(X^2)$ en ajoutant une ligne $x_i^2$ au tableau.
\item \textbf{Écart-type} : prendre la racine carrée :
\[ \sigma(X) = \sqrt{V(X)}. \]
\item \textbf{Interpréter} : $E(X)$ est la valeur moyenne attendue sur un grand nombre de répétitions ; $\sigma(X)$ mesure la dispersion autour de cette moyenne.
\end{enumerate}
\textbf{Réflexe} : toujours vérifier que $V(X) \geq 0$ ; une variance négative signale obligatoirement une erreur de calcul.
\end{methode}

\begin{exoresolu}[Rentabilité d'un jeu de hasard]
On reprend la tombola de l'exercice précédent : $X$ prend les valeurs $0$, $1000$, $5000$ avec les probabilités $0{,}6$ ; $0{,}3$ ; $0{,}1$.
\begin{enumerate}
\item Calculer $E(X)$, $V(X)$ et $\sigma(X)$.
\item Le ticket coûte \SI{500}{F}. Le jeu est-il favorable au joueur ?
\end{enumerate}
\tcblower
\textbf{1.} Complétons le tableau de la loi avec les lignes utiles :
\begin{center}
\begin{tabularx}{\linewidth}{|l|*{3}{>{\centering\arraybackslash}X|}>{\centering\arraybackslash}p{2.2cm}|}
\hline
\rowcolor{popblueL} $x_i$ & $0$ & $1000$ & $5000$ & Total \\
\hline
$p_i$ & $0{,}6$ & $0{,}3$ & $0{,}1$ & $1$ \\
\hline
$x_i p_i$ & $0$ & $300$ & $500$ & $800$ \\
\hline
$x_i^2 p_i$ & $0$ & \num{300000} & \num{2500000} & \num{2800000} \\
\hline
\end{tabularx}
\end{center}

\textbf{Espérance} :
\[ E(X) = 0 \times 0{,}6 + 1000 \times 0{,}3 + 5000 \times 0{,}1 = 0 + 300 + 500 = 800 \text{ F}. \]

\textbf{Variance} (formule de König-Huygens) :
\begin{align*}
E(X^2) &= 0^2 \times 0{,}6 + 1000^2 \times 0{,}3 + 5000^2 \times 0{,}1 = \num{300000} + \num{2500000} = \num{2800000} ; \\
V(X) &= E(X^2) - [E(X)]^2 = \num{2800000} - 800^2 = \num{2800000} - \num{640000} = \num{2160000}.
\end{align*}

\textbf{Écart-type} :
\[ \sigma(X) = \sqrt{\num{2160000}} \approx \num{1469,7} \text{ F}. \]

\textbf{2.} Le gain moyen par partie est $E(X) = \SI{800}{F}$, alors que le ticket coûte \SI{500}{F}. Le gain moyen net est $800 - 500 = \SI{300}{F} > 0$ : le jeu est donc \textbf{favorable au joueur} en moyenne. (En pratique, l'organisateur aurait intérêt à vendre le ticket plus cher !)

\textbf{Conclusion} : $E(X) = \SI{800}{F}$, $V(X) = \num{2160000}$ et $\sigma(X) \approx \SI{1470}{F}$. La forte valeur de $\sigma(X)$ traduit une grande dispersion des gains autour de la moyenne.
\end{exoresolu}

\courssub{Épreuve de Bernoulli et schéma de Bernoulli}

\begin{definition}[Épreuve de Bernoulli]
Une \cle{épreuve de Bernoulli} est une expérience aléatoire qui ne présente que \emph{deux issues} possibles :
\begin{itemize}
\item \textbf{Succès} (noté $S$), de probabilité $p$ ($0 < p < 1$).
\item \textbf{Échec} (noté $\bar{S}$ ou $E$), de probabilité $q = 1 - p$.
\end{itemize}
Exemples : lancer une pièce (Pile/Face), tester un composant (Conforme/Défectueux), un appel qui aboutit ou non.
\end{definition}

\begin{definition}[Schéma de Bernoulli]
On répète $n$ fois ($n \in \mathbb{N}^*$) une même épreuve de Bernoulli de paramètre $p$, de manière \emph{identique et indépendante}.
L'indépendance est garantie en pratique par un \emph{tirage avec remise}, ou par un tirage dans une population très grande (assimilé à un tirage avec remise).
La variable aléatoire $X$ qui compte le nombre de succès sur les $n$ épreuves suit la \cle{loi binomiale} de paramètres $n$ et $p$, notée $X \hookrightarrow \mathcal{B}(n, p)$.
\end{definition}

\begin{propriete}[Loi binomiale]
Si $X \hookrightarrow \mathcal{B}(n, p)$, alors pour tout $k \in \{0, 1, \ldots, n\}$ :
\[ P(X = k) = \binom{n}{k} p^k (1-p)^{n-k}. \]
Espérance et variance :
\[ E(X) = np, \qquad V(X) = np(1-p), \qquad \sigma(X) = \sqrt{np(1-p)}. \]
\end{propriete}

\begin{experience}[Découverte de la formule de la loi binomiale]
On considère un schéma de Bernoulli à $n=3$ épreuves de paramètre $p$.
\begin{enumerate}
\item Dessiner l'arbre des 3 épreuves (succès $S$ prob. $p$, échec $E$ prob. $q=1-p$).
\item Combien de chemins mènent à exactement $k=1$ succès ? Écrire leur probabilité commune. En déduire $P(X=1)$.
\item Faire de même pour $k=2$ et $k=3$.
\item Généraliser pour $n$ épreuves : pourquoi le coefficient $\binom{n}{k}$ apparaît-il ?
\end{enumerate}
\tcblower
\textbf{1.} L'arbre a $2^3 = 8$ feuilles.
\textbf{2.} $k=1$ succès : 3 chemins ($SEE, ESE, EES$). Chacun a probabilité $p q^2$. $P(X=1) = 3 p q^2 = \binom{3}{1} p^1 q^2$.
\textbf{3.} $k=2$ : 3 chemins ($SSE, SES, ESS$), prob. $p^2 q$. $P(X=2) = 3 p^2 q = \binom{3}{2} p^2 q^1$.
$k=3$ : 1 chemin ($SSS$), prob. $p^3$. $P(X=3) = \binom{3}{3} p^3 q^0$.
\textbf{4.} Pour $n$ épreuves, obtenir $k$ succès revient à choisir les $k$ positions des succès parmi les $n$ épreuves : $\binom{n}{k}$ choix. Chaque chemin a probabilité $p^k q^{n-k}$. D'où la formule générale.
\end{experience}

\begin{methode}[Identifier une épreuve et un schéma de Bernoulli]
\begin{enumerate}
\item \textbf{Épreuve de Bernoulli} : vérifier que l'expérience n'a que \emph{deux issues} : succès $S$ (probabilité $p$) et échec $\bar{S}$ (probabilité $q = 1 - p$).
\item \textbf{Schéma de Bernoulli} : vérifier que l'on répète $n$ fois cette épreuve de façon \emph{identique et indépendante}. L'indépendance est garantie en pratique par un \emph{tirage avec remise}, ou par un tirage dans une population très grande (assimilé à un tirage avec remise).
\item \textbf{Conclure} : la variable aléatoire $X$ qui compte le nombre de succès suit la \cle{loi binomiale} $\mathcal{B}(n, p)$.
\item \textbf{Tracer un arbre} si $n$ est petit (2 ou 3) : chaque chemin a une probabilité obtenue en multipliant les probabilités des branches.
\end{enumerate}
\begin{attention}[Condition d'indépendance]
Un tirage \emph{sans remise} dans une petite urne ne donne \textbf{pas} un schéma de Bernoulli : les tirages ne sont pas indépendants. Ne pas appliquer la loi binomiale dans ce cas !
\end{attention}
\end{methode}

\begin{exoresolu}[Contrôle qualité chez un embouteilleur]
Une usine d'embouteillage de Bafoussam produit des bouteilles dont $2\,\%$ présentent un défaut de bouchage. On prélève au hasard 3 bouteilles dans la production, avec remise. On note $X$ le nombre de bouteilles défectueuses.
\begin{enumerate}
\item Justifier que $X$ suit une loi binomiale et préciser ses paramètres.
\item Construire l'arbre de probabilités et déterminer la loi de $X$.
\end{enumerate}
\tcblower
\textbf{1.} Chaque prélèvement est une épreuve de Bernoulli : succès $S$ = « la bouteille est défectueuse » de probabilité $p = 0{,}02$, échec $\bar{S}$ de probabilité $q = 1 - 0{,}02 = 0{,}98$.

On répète $n = 3$ fois cette épreuve, de manière identique et indépendante (tirage \emph{avec remise}). Donc $X$, qui compte le nombre de succès, suit la loi binomiale :
\[ X \hookrightarrow \mathcal{B}(3\,;\,0{,}02). \]

\textbf{2.} Arbre de probabilités :
\begin{popfigure}
\begin{tikzpicture}[grow=right, level distance=2.6cm,
level 1/.style={sibling distance=2.6cm},
level 2/.style={sibling distance=1.3cm},
level 3/.style={sibling distance=0.7cm},
every node/.style={font=\small}]
\node {}
child { node {$\bar{S}$}
  child { node {$\bar{S}$}
    child { node {$\bar{S}$} edge from parent node[below, popmuted] {\scriptsize $0{,}98$} }
    child { node {$S$} edge from parent node[above, popmuted] {\scriptsize $0{,}02$} }
    edge from parent node[below, popmuted] {\scriptsize $0{,}98$} }
  child { node {$S$}
    child { node {$\bar{S}$} edge from parent node[below, popmuted] {\scriptsize $0{,}98$} }
    child { node {$S$} edge from parent node[above, popmuted] {\scriptsize $0{,}02$} }
    edge from parent node[above, popmuted] {\scriptsize $0{,}02$} }
  edge from parent node[below, popmuted] {\scriptsize $0{,}98$} }
child { node {$S$}
  child { node {$\bar{S}$}
    child { node {$\bar{S}$} edge from parent node[below, popmuted] {\scriptsize $0{,}98$} }
    child { node {$S$} edge from parent node[above, popmuted] {\scriptsize $0{,}02$} }
    edge from parent node[below, popmuted] {\scriptsize $0{,}98$} }
  child { node {$S$}
    child { node {$\bar{S}$} edge from parent node[below, popmuted] {\scriptsize $0{,}98$} }
    child { node {$S$} edge from parent node[above, popmuted] {\scriptsize $0{,}02$} }
    edge from parent node[above, popmuted] {\scriptsize $0{,}02$} }
  edge from parent node[above, popmuted] {\scriptsize $0{,}02$} };
\end{tikzpicture}
\legende{Schéma de Bernoulli à 3 épreuves : $p = 0{,}02$, $q = 0{,}98$.}
\end{popfigure}

En comptant les succès sur chaque chemin et en multipliant les probabilités le long des branches :
\begin{align*}
P(X = 0) &= (0{,}98)^3 = \num{0,941192} ; \\
P(X = 1) &= 3 \times 0{,}02 \times (0{,}98)^2 = 3 \times \num{0,019208} = \num{0,057624} ; \\
P(X = 2) &= 3 \times (0{,}02)^2 \times 0{,}98 = \num{0,001176} ; \\
P(X = 3) &= (0{,}02)^3 = \num{0,000008}.
\end{align*}
\textbf{Contrôle} : $\num{0,941192} + \num{0,057624} + \num{0,001176} + \num{0,000008} = 1$.

\textbf{Conclusion} : $X \hookrightarrow \mathcal{B}(3 ; 0{,}02)$ et sa loi est donnée par le tableau ci-dessus ; la probabilité d'avoir au moins une bouteille défectueuse sur les 3 est $1 - P(X=0) \approx 0{,}0588$.
\end{exoresolu}

\courssub{Loi binomiale et applications}

\begin{propriete}[Formules usuelles de la loi binomiale $\mathcal{B}(n, p)$]
Si $X \hookrightarrow \mathcal{B}(n, p)$ :
\begin{itemize}
\item $P(X = k) = \binom{n}{k} p^k (1-p)^{n-k}$, pour $k \in \{0, \ldots, n\}$.
\item $P(X \geq 1) = 1 - P(X = 0) = 1 - (1-p)^n$ (très utile pour « au moins un »).
\item $P(X \leq k) = \sum_{i=0}^{k} P(X=i)$ (somme des probabilités).
\item $E(X) = np$, $V(X) = np(1-p)$, $\sigma(X) = \sqrt{np(1-p)}$.
\end{itemize}
\end{propriete}

\begin{methode}[Exploiter la loi binomiale $\mathcal{B}(n, p)$]
\begin{enumerate}
\item \textbf{Identifier} $n$ (nombre d'épreuves) et $p$ (probabilité de succès à chaque épreuve). Vérifier l'indépendance.
\item \textbf{Annoncer clairement} : « $X$ suit la loi binomiale de paramètres $n = \ldots$ et $p = \ldots$ ».
\item \textbf{Appliquer la formule} adaptée à la question :
\begin{itemize}
\item « exactement $k$ » : $P(X=k) = \binom{n}{k} p^k (1-p)^{n-k}$.
\item « au moins $k$ » : $P(X \geq k) = 1 - P(X \leq k-1)$ ou somme directe si $k$ proche de $n$.
\item « au plus $k$ » : $P(X \leq k) = \sum_{i=0}^k P(X=i)$.
\end{itemize}
\item \textbf{Calculer} avec la calculatrice (fonction \texttt{binompdf} et \texttt{binomcdf} si disponible) ou à la main si $n$ petit.
\item \textbf{Interpréter} le résultat en pourcentage ou en décision.
\end{enumerate}
\end{methode}

\begin{exoresolu}[Réseau mobile en zone rurale]
Dans une localité, la probabilité qu'un appel passé sur le réseau aboutisse du premier coup est $0{,}8$. Un technicien effectue 10 appels tests, supposés indépendants. On note $X$ le nombre d'appels qui aboutissent.
\begin{enumerate}
\item Quelle est la loi de $X$ ? Calculer $E(X)$ et $\sigma(X)$.
\item Calculer la probabilité qu'exactement 8 appels aboutissent.
\item Calculer la probabilité qu'au moins 8 appels aboutissent (arrondir à $10^{-4}$).
\end{enumerate}
\tcblower
\textbf{1.} Chaque appel est une épreuve de Bernoulli de succès $p = 0{,}8$. Les 10 appels sont indépendants, donc :
\[ X \hookrightarrow \mathcal{B}(10\,;\,0{,}8). \]
Espérance : $E(X) = np = 10 \times 0{,}8 = 8$. En moyenne, 8 appels sur 10 aboutissent.

Variance : $V(X) = np(1-p) = 10 \times 0{,}8 \times 0{,}2 = 1{,}6$, d'où $\sigma(X) = \sqrt{1{,}6} \approx 1{,}26$.

\textbf{2.} D'après la formule de la loi binomiale avec $k = 8$ :
\begin{align*}
P(X = 8) &= \binom{10}{8} (0{,}8)^8 (0{,}2)^2 \\
&= 45 \times \num{0,16777216} \times 0{,}04 \\
&\approx \num{0,3020}.
\end{align*}

\textbf{3.} « Au moins 8 » signifie $X = 8$, $X = 9$ ou $X = 10$ :
\begin{align*}
P(X = 9) &= \binom{10}{9} (0{,}8)^9 (0{,}2)^1 = 10 \times \num{0,134217728} \times 0{,}2 \approx \num{0,2684} ; \\
P(X = 10) &= (0{,}8)^{10} \approx \num{0,1074}.
\end{align*}
Donc :
\[ P(X \geq 8) = P(X=8) + P(X=9) + P(X=10) \approx 0{,}3020 + 0{,}2684 + 0{,}1074 = 0{,}6778. \]

\textbf{Conclusion} : la probabilité qu'au moins 8 des 10 appels aboutissent est d'environ $0{,}6778$, soit environ $67{,}8\,\%$. Le réseau est conforme si l'opérateur exige ce seuil avec une probabilité d'au moins $60\,\%$.
\end{exoresolu}

\begin{exoresolu}[Problème de décision : deux fournisseurs]
Une entreprise de Buea achète ses câbles électriques chez deux fournisseurs. Chez le fournisseur A, $5\,\%$ des câbles sont défectueux et le câble coûte \SI{9000}{F}. Chez le fournisseur B, $10\,\%$ sont défectueux et le câble coûte \SI{8000}{F}. Un câble défectueux entraîne un coût de réparation de \SI{20000}{F}. On note $C_A$ et $C_B$ le coût total (achat + réparation éventuelle) d'un câble selon le fournisseur.
\begin{enumerate}
\item Déterminer les lois de $C_A$ et $C_B$.
\item Calculer $E(C_A)$ et $E(C_B)$, puis conseiller l'entreprise.
\end{enumerate}
\tcblower
\textbf{1.} Chez A : le câble coûte \SI{9000}{F} s'il est bon (probabilité $0{,}95$) et $9000 + 20000 = \SI{29000}{F}$ s'il est défectueux (probabilité $0{,}05$).
\begin{center}
\begin{tabularx}{\linewidth}{|l|*{2}{>{\centering\arraybackslash}X|}}
\hline
\rowcolor{popblueL} $c$ & $9000$ & $29000$ \\
\hline
$P(C_A = c)$ & $0{,}95$ & $0{,}05$ \\
\hline
\end{tabularx}
\end{center}
Chez B : le câble coûte \SI{8000}{F} (probabilité $0{,}9$) ou \SI{28000}{F} (probabilité $0{,}1$).
\begin{center}
\begin{tabularx}{\linewidth}{|l|*{2}{>{\centering\arraybackslash}X|}}
\hline
\rowcolor{popblueL} $c$ & $8000$ & $28000$ \\
\hline
$P(C_B = c)$ & $0{,}9$ & $0{,}1$ \\
\hline
\end{tabularx}
\end{center}
Contrôle : $0{,}95 + 0{,}05 = 1$ et $0{,}9 + 0{,}1 = 1$.

\textbf{2.} Calcul des espérances :
\begin{align*}
E(C_A) &= 9000 \times 0{,}95 + 29000 \times 0{,}05 = 8550 + 1450 = \SI{10000}{F} ; \\
E(C_B) &= 8000 \times 0{,}9 + 28000 \times 0{,}1 = 7200 + 2800 = \SI{10000}{F}.
\end{align*}
Les deux coûts moyens sont \textbf{égaux} : \SI{10000}{F}. Pour départager, calculons les variances :
\begin{align*}
V(C_A) &= 9000^2 \times 0{,}95 + 29000^2 \times 0{,}05 - 10000^2 = \num{76950000} + \num{42050000} - \num{100000000} = \num{19000000} ; \\
V(C_B) &= 8000^2 \times 0{,}9 + 28000^2 \times 0{,}1 - 10000^2 = \num{57600000} + \num{78400000} - \num{100000000} = \num{36000000}.
\end{align*}
Donc $\sigma(C_A) \approx \SI{4359}{F}$ et $\sigma(C_B) = \SI{6000}{F}$.

\textbf{Conclusion} : à espérance de coût égale, le fournisseur A présente un écart-type plus faible, donc un risque moindre de mauvaise surprise. L'entreprise prudente choisira le \textbf{fournisseur A}.
\end{exoresolu}

\begin{savaistu}[De Bernoulli à nos téléphones : l'histoire d'une loi universelle]
La loi binomiale doit son nom à \textbf{Jakob Bernoulli} (1655--1705), mathématicien suisse de la célèbre famille Bernoulli. Dans son ouvrage posthume \emph{Ars Conjectandi} (1713), il démontre le \textbf{théorème de Bernoulli} (loi faible des grands nombres) : la fréquence d'un événement converge vers sa probabilité.
Son ami \textbf{Abraham de Moivre} (1667--1754) a ensuite montré que pour $n$ grand, la loi binomiale $\mathcal{B}(n,p)$ est bien approchée par une loi normale (théorème central limite), ce qui évitait des calculs fastidieux à l'époque.
\textbf{Au Cameroun aujourd'hui} : cette loi modélise le taux d'erreur binaire (BER) dans les fibres optiques de CAMTEL, le taux de germination des semences améliorées à l'IRAD, la fiabilité des antennes relais MTN/Orange, ou encore le taux de réussite au Baccalauréat dans un établissement. C'est l'outil mathématique de base de la \emph{fiabilité industrielle} et du \emph{contrôle qualité}.
\end{savaistu}

\begin{attention}[Les 5 erreurs qui coûtent des points au Bac]
\begin{enumerate}
\item \textbf{Appliquer la loi binomiale sans vérifier l'indépendance.} Un tirage sans remise dans une petite population n'est PAS un schéma de Bernoulli. Rédige toujours la justification : « les épreuves sont identiques et indépendantes (tirage avec remise), donc $X \hookrightarrow \mathcal{B}(n, p)$ ».
\item \textbf{Oublier le coefficient binomial.} Écrire $P(X = k) = p^k (1-p)^{n-k}$ sans le $\binom{n}{k}$ est faux dès que $k \neq 0$ et $k \neq n$. Le coefficient compte le nombre de chemins de l'arbre menant à $k$ succès.
\item \textbf{Confondre « ou » et « et ».} $P(A \cup B) \neq P(A) + P(B)$ en général : il faut retrancher $P(A \cap B)$. N'oublie ce terme que si tu as \emph{justifié} que les événements sont incompatibles.
\item \textbf{Ne pas vérifier la loi de probabilité.} Avant de calculer $E(X)$, contrôle que la somme des $p_i$ vaut 1 ; une variance négative ou une probabilité supérieure à 1 doit te faire reprendre tes calculs immédiatement.
\item \textbf{Mal traiter le « au moins un ».} Ne calcule jamais $P(X \geq 1)$ en sommant tous les cas : utilise $P(X \geq 1) = 1 - P(X = 0) = 1 - (1-p)^n$. Et dans un dénombrement, « au moins une boule verte » se traite par l'événement contraire, pas par une somme de cas fastidieuse.
\end{enumerate}
\end{attention}

\newpage

\coursec{Exercices}

\courssub{Je vérifie mes connaissances}

\begin{exercice}[QCM : vocabulaire et calculs de base]
Pour chaque question, une seule réponse est exacte. Entoure-la et justifie brièvement.
\begin{enumerate}
\item On lance un dé équilibré à 6 faces. La probabilité d'obtenir un nombre pair est :
\begin{itemize}
\item[a)] $\dfrac{1}{6}$ \quad b) $\dfrac{1}{3}$ \quad c) $\dfrac{1}{2}$ \quad d) $\dfrac{2}{3}$
\end{itemize}
\item Si $A$ et $B$ sont deux événements incompatibles avec $P(A)=0{,}3$ et $P(B)=0{,}5$, alors $P(A\cup B)$ vaut :
\begin{itemize}
\item[a)] $0{,}15$ \quad b) $0{,}8$ \quad c) $0{,}2$ \quad d) $1$
\end{itemize}
\item Une variable aléatoire $X$ suit la loi binomiale $\mathcal{B}(10\,;\,0{,}2)$. Son espérance est :
\begin{itemize}
\item[a)] $0{,}2$ \quad b) $2$ \quad c) $1{,}6$ \quad d) $8$
\end{itemize}
\item Soit $X$ une variable aléatoire de variance $V(X)=9$. Son écart-type $\sigma(X)$ vaut :
\begin{itemize}
\item[a)] $81$ \quad b) $4{,}5$ \quad c) $3$ \quad d) $9$
\end{itemize}
\end{enumerate}
\end{exercice}

\begin{exercice}[Vrai ou faux justifié]
Réponds par VRAI ou FAUX en justifiant chaque réponse par une propriété du cours ou un contre-exemple.
\begin{enumerate}
\item Pour tout événement $A$, on a $P(\bar{A}) = 1 - P(A)$.
\item Si $P(A)=0{,}4$ et $P(B)=0{,}7$, alors $P(A\cup B)=1{,}1$.
\item L'espérance d'une variable aléatoire est toujours un nombre positif.
\item Si $X$ suit la loi binomiale $\mathcal{B}(n\,;\,p)$, alors $V(X)=np(1-p)$.
\end{enumerate}
\end{exercice}

\begin{exercice}[Définitions à compléter]
Recopie et complète les phrases suivantes avec les mots ou expressions du cours.
\begin{enumerate}
\item Une expérience dont on ne peut pas prévoir le résultat avec certitude, mais dont on connaît l'ensemble des résultats possibles, est appelée \ldots\ ; cet ensemble de résultats est appelé \ldots\ et noté \ldots
\item Dans le cas de l'\ldots, la probabilité d'un événement $A$ est $P(A)=\dfrac{\ldots}{\ldots}$.
\item Une \ldots\ est une application qui, à chaque issue d'une expérience aléatoire, associe un nombre réel.
\item Une épreuve de Bernoulli est une expérience aléatoire qui n'a que \ldots\ issues : le \ldots\ (probabilité $p$) et l'\ldots\ (probabilité $1-p$).
\end{enumerate}
\end{exercice}

\begin{exercice}[Calculs directs]
\begin{enumerate}
\item On lance deux pièces de monnaie équilibrées. Donner l'univers $\Omega$ et calculer la probabilité d'obtenir exactement une fois « Pile ».
\item Soit $A$ et $B$ deux événements tels que $P(A)=0{,}5$, $P(B)=0{,}4$ et $P(A\cap B)=0{,}2$. Calculer $P(A\cup B)$ et $P(\bar{A})$.
\item Une variable aléatoire $X$ prend les valeurs $0$, $1$, $2$ avec les probabilités $0{,}2$, $0{,}5$ et $0{,}3$. Calculer $E(X)$.
\item $X$ suit la loi binomiale $\mathcal{B}(50\,;\,0{,}1)$. Donner $E(X)$, $V(X)$ et $\sigma(X)$.
\end{enumerate}
\end{exercice}

\courssub{Je m'entraîne}

\begin{exercice}[Loto du quartier]
Une urne contient 10 boules indiscernables au toucher : 3 rouges, 5 bleues et 2 vertes. On tire une boule au hasard.
\begin{enumerate}
\item Calculer la probabilité de chacun des événements : $R$ : « la boule est rouge », $B$ : « la boule est bleue », $V$ : « la boule est verte ».
\item Calculer la probabilité que la boule tirée ne soit pas bleue.
\item On tire maintenant deux boules successivement et avec remise. Construire l'arbre de probabilités (on ne distinguera que « rouge » et « non rouge ») et calculer la probabilité d'obtenir deux boules rouges.
\end{enumerate}
\end{exercice}

\begin{exercice}[Gain à la tombola]
À la fête de la jeunesse à Ngaoundéré, une tombola propose 100 tickets à \SI{200}{FCFA} l'unité. Parmi eux, 2 tickets font gagner \SI{5000}{FCFA}, 8 tickets font gagner \SI{1000}{FCFA}, et les autres ne rapportent rien. On note $X$ le gain algébrique (gain brut moins le prix du ticket) d'un joueur qui achète un ticket.
\begin{enumerate}
\item Déterminer la loi de probabilité de $X$ (présenter les résultats dans un tableau).
\item Calculer l'espérance $E(X)$. Interpréter le résultat : le jeu est-il favorable au joueur ?
\item Calculer la variance $V(X)$ puis l'écart-type $\sigma(X)$ (arrondir au franc près).
\end{enumerate}
\end{exercice}

\begin{exercice}[Contrôle qualité chez un opérateur]
Un centre d'appel d'Orange Cameroun constate que la probabilité qu'un appel dure plus de 5 minutes est $0{,}4$. On interroge au hasard 5 appels, supposés indépendants, et on note $X$ le nombre d'appels durant plus de 5 minutes.
\begin{enumerate}
\item Justifier que $X$ suit une loi binomiale dont on précisera les paramètres.
\item Calculer la probabilité qu'exactement 2 appels durent plus de 5 minutes (arrondir à $10^{-3}$).
\item Calculer la probabilité qu'au moins un appel dure plus de 5 minutes.
\item Donner $E(X)$ et interpréter.
\end{enumerate}
\end{exercice}

\begin{exercice}[Fonction de répartition]
Soit $X$ une variable aléatoire dont la fonction de répartition $F$ est donnée par :
\[
F(x)=
\begin{cases}
0 & \text{si } x<1\\
0{,}1 & \text{si } 1\le x<2\\
0{,}4 & \text{si } 2\le x<3\\
0{,}8 & \text{si } 3\le x<4\\
1 & \text{si } x\ge 4
\end{cases}
\]
\begin{enumerate}
\item Déterminer la loi de probabilité de $X$.
\item Calculer $E(X)$ et $V(X)$.
\item Calculer $P(X>2)$ et $P(1\le X\le 3)$.
\end{enumerate}
\end{exercice}

\courssub{Je me prépare au Bac}

\begin{exercice}[Tirage dans une urne — type Bac]
Une urne contient 4 boules blanches et 6 boules noires, indiscernables au toucher. On tire successivement et sans remise 2 boules de l'urne.

\textbf{Partie A : étude du tirage}
\begin{enumerate}
\item Construire l'arbre de probabilités complet de cette expérience, en précisant les probabilités sur chaque branche.
\item Calculer la probabilité des événements suivants :
\begin{enumerate}
\item $A$ : « les deux boules tirées sont blanches » ;
\item $B$ : « les deux boules tirées sont de la même couleur » ;
\item $C$ : « les deux boules tirées sont de couleurs différentes ».
\end{enumerate}
\item Vérifier que $P(C)=1-P(B)$. Expliquer pourquoi ce résultat était prévisible.
\end{enumerate}

\textbf{Partie B : variable aléatoire}

On note $X$ le nombre de boules blanches obtenues à l'issue du tirage.
\begin{enumerate}
\setcounter{enumi}{3}
\item Déterminer la loi de probabilité de $X$ (présenter les résultats dans un tableau, sous forme de fractions irréductibles).
\item Calculer l'espérance $E(X)$ et l'écart-type $\sigma(X)$ (arrondir à $10^{-2}$).
\item On propose le jeu suivant : le joueur mise \SI{500}{FCFA} ; il gagne \SI{1000}{FCFA} par boule blanche obtenue. On note $Y$ le gain algébrique du joueur. Exprimer $Y$ en fonction de $X$, puis calculer $E(Y)$. Le jeu est-il équitable ?
\end{enumerate}
\end{exercice}

\begin{exercice}[Sachets mal scellés à Bonabéri — type Bac]
Dans une usine de conditionnement d'arachides de Bonabéri (Douala), on estime que $3\,\%$ des sachets sont mal scellés. La production journalière est suffisamment importante pour qu'on puisse assimiler tout prélèvement à un tirage avec remise. On prélève au hasard 20 sachets dans la production et on note $X$ le nombre de sachets mal scellés dans ce prélèvement. On donnera les résultats arrondis à $10^{-4}$. On rappelle : $0{,}97^{20}\approx 0{,}5438$.

\textbf{Partie A : loi de $X$}
\begin{enumerate}
\item Justifier soigneusement que $X$ suit une loi binomiale et préciser ses paramètres.
\item Calculer la probabilité qu'aucun sachet du prélèvement ne soit mal scellé.
\item Calculer la probabilité qu'exactement un sachet soit mal scellé.
\item En déduire la probabilité qu'au moins deux sachets soient mal scellés.
\item Calculer le nombre moyen de sachets mal scellés dans un prélèvement de 20 sachets.
\end{enumerate}

\textbf{Partie B : décision commerciale}

Le chef de production affirme : « Sur une journée de \num{10000} sachets, il y a forcément moins de 400 sachets mal scellés. »
\begin{enumerate}
\setcounter{enumi}{5}
\item Calculer l'espérance du nombre de sachets mal scellés sur une journée de \num{10000} sachets et commenter l'affirmation du chef de production.
\item L'usine souhaite que la probabilité qu'un prélèvement de 20 sachets ne contienne aucun sachet mal scellé soit supérieure à $0{,}9$. On note $p$ le taux de sachets défectueux. Montrer que la condition s'écrit $(1-p)^{20}>0{,}9$, puis déterminer la valeur maximale de $p$ (arrondir à $10^{-3}$). On donne : $0{,}9^{1/20}\approx 0{,}9947$.
\end{enumerate}
\end{exercice}

\begin{exercice}[QCM du concours d'entrée — type Bac]
Un candidat répond totalement au hasard à un QCM de 8 questions. Chaque question propose 4 réponses dont une seule est correcte, et les réponses sont indépendantes les unes des autres. On note $X$ le nombre de bonnes réponses du candidat. Les résultats seront arrondis à $10^{-4}$.

\textbf{Partie A : étude probabiliste}
\begin{enumerate}
\item Justifier que $X$ suit une loi binomiale et préciser ses paramètres.
\item Calculer la probabilité que le candidat donne exactement 4 bonnes réponses.
\item Calculer la probabilité que le candidat donne au moins la moitié de bonnes réponses, c'est-à-dire $P(X\ge 4)$. On donne :
\[
P(X=4)\approx 0{,}0865,\quad P(X=5)\approx 0{,}0231,\quad P(X=6)\approx 0{,}0038,\quad P(X=7)\approx 0{,}0004.
\]
\item Calculer l'espérance $E(X)$ et l'écart-type $\sigma(X)$. Interpréter $E(X)$.
\end{enumerate}

\textbf{Partie B : barème et prime d'assurance}

Le concours attribue $+1$ point par bonne réponse et $-0{,}25$ point par mauvaise réponse. On note $N$ la note finale du candidat sur 8.
\begin{enumerate}
\setcounter{enumi}{4}
\item Exprimer $N$ en fonction de $X$, puis calculer $E(N)$. Commenter.
\item Une compagnie d'assurance de Yaoundé assure ses clients contre les dégâts des eaux. Elle estime que la probabilité qu'un client déclare un sinistre dans l'année est $0{,}05$ et que le coût moyen d'un sinistre est de \SI{400000}{FCFA}. En s'inspirant de la notion d'espérance, calculer le coût moyen annuel par client pour la compagnie, puis proposer une prime annuelle minimale (arrondie au millier de francs supérieur) permettant à la compagnie d'être rentable, sachant qu'elle souhaite dégager une marge de $20\,\%$ sur le coût moyen.
\end{enumerate}
\end{exercice}

\begin{attention}[Erreurs fréquentes à éviter au Bac]
\begin{itemize}
\item Confondre tirage \textbf{avec remise} (probabilités constantes, indépendance, loi binomiale possible) et tirage \textbf{sans remise} (les probabilités changent au second tirage : pas de loi binomiale).
\item Écrire $P(A\cup B)=P(A)+P(B)$ sans vérifier que $A$ et $B$ sont incompatibles ; la formule générale est $P(A\cup B)=P(A)+P(B)-P(A\cap B)$.
\item Oublier le coefficient binomial $\dbinom{n}{k}$ dans $P(X=k)=\dbinom{n}{k}p^k(1-p)^{n-k}$.
\item Pour « au moins un », penser à passer par l'événement contraire : $P(X\ge 1)=1-P(X=0)$.
\item Donner la variance à la place de l'écart-type : $\sigma(X)=\sqrt{V(X)}$.
\end{itemize}
\end{attention}

\newpage

\coursec{Corrigés des exercices}

\begin{corrige}[Exercice 1 — QCM : vocabulaire et calculs de base]
\begin{enumerate}
\item Réponse \textbf{c)}. Les issues paires sont $\{2\,;\,4\,;\,6\}$, soit 3 issues sur 6. Par équiprobabilité : $P=\dfrac{3}{6}=\dfrac{1}{2}$.
\item Réponse \textbf{b)}. $A$ et $B$ incompatibles signifie $A\cap B=\varnothing$, donc
\[P(A\cup B)=P(A)+P(B)=0{,}3+0{,}5=0{,}8.\]
\item Réponse \textbf{b)}. Pour $X\sim\mathcal{B}(n\,;\,p)$, $E(X)=np=10\times 0{,}2=2$.
\item Réponse \textbf{c)}. $\sigma(X)=\sqrt{V(X)}=\sqrt{9}=3$.
\end{enumerate}
\end{corrige}

\begin{corrige}[Exercice 2 — Vrai ou faux justifié]
\begin{enumerate}
\item \textbf{VRAI.} $A$ et $\bar{A}$ forment une partition de $\Omega$, donc $P(A)+P(\bar{A})=P(\Omega)=1$, d'où $P(\bar{A})=1-P(A)$.
\item \textbf{FAUX.} Une probabilité est toujours comprise entre $0$ et $1$ : $P(A\cup B)=1{,}1$ est impossible. La formule correcte est $P(A\cup B)=P(A)+P(B)-P(A\cap B)$ ; sans connaître $P(A\cap B)$, on sait seulement que $P(A\cup B)\le 1$.
\item \textbf{FAUX.} Contre-exemple : $X$ prend les valeurs $-2$ et $1$ avec les probabilités $0{,}5$ et $0{,}5$ ; alors $E(X)=-2\times 0{,}5+1\times 0{,}5=-0{,}5<0$. L'espérance d'un gain algébrique peut être négative (jeu défavorable).
\item \textbf{VRAI.} C'est une propriété du cours : si $X\sim\mathcal{B}(n\,;\,p)$, alors $V(X)=np(1-p)$ et $\sigma(X)=\sqrt{np(1-p)}$.
\end{enumerate}
\end{corrige}

\begin{corrige}[Exercice 3 — Définitions à compléter]
\begin{enumerate}
\item \ldots est appelée \cle{expérience aléatoire} ; cet ensemble de résultats est appelé \cle{univers} (ou univers des possibles) et noté \cle{$\Omega$}.
\item Dans le cas de l'\cle{équiprobabilité}, la probabilité d'un événement $A$ est
\[P(A)=\dfrac{\text{nombre d'issues favorables à } A}{\text{nombre total d'issues possibles}}=\dfrac{\text{Card}(A)}{\text{Card}(\Omega)}.\]
\item Une \cle{variable aléatoire} est une application qui, à chaque issue d'une expérience aléatoire, associe un nombre réel.
\item Une épreuve de Bernoulli n'a que \cle{deux} issues : le \cle{succès} (probabilité $p$) et l'\cle{échec} (probabilité $1-p$).
\end{enumerate}
\end{corrige}

\begin{corrige}[Exercice 4 — Calculs directs]
\begin{enumerate}
\item $\Omega=\{PP\,;\,PF\,;\,FP\,;\,FF\}$ : 4 issues équiprobables. Les issues contenant exactement un « Pile » sont $PF$ et $FP$ :
\[P=\dfrac{2}{4}=\dfrac{1}{2}.\]
\item $P(A\cup B)=P(A)+P(B)-P(A\cap B)=0{,}5+0{,}4-0{,}2=0{,}7$ ; \quad $P(\bar{A})=1-P(A)=1-0{,}5=0{,}5$.
\item $E(X)=0\times 0{,}2+1\times 0{,}5+2\times 0{,}3=0+0{,}5+0{,}6=1{,}1$.
\item $E(X)=np=50\times 0{,}1=5$ ; \quad $V(X)=np(1-p)=50\times 0{,}1\times 0{,}9=4{,}5$ ; \quad $\sigma(X)=\sqrt{4{,}5}\approx 2{,}12$.
\end{enumerate}
\end{corrige}

\begin{corrige}[Exercice 5 — Loto du quartier]
\begin{enumerate}
\item Par équiprobabilité (boules indiscernables au toucher) :
\[P(R)=\dfrac{3}{10}=0{,}3 ;\qquad P(B)=\dfrac{5}{10}=0{,}5 ;\qquad P(V)=\dfrac{2}{10}=0{,}2.\]
Vérification : $0{,}3+0{,}5+0{,}2=1$.
\item $P(\bar{B})=1-P(B)=1-0{,}5=0{,}5$. On retrouve bien $P(R)+P(V)=0{,}3+0{,}2=0{,}5$.
\item À chaque tirage, $P(R)=0{,}3$ et $P(\bar{R})=0{,}7$ : la remise rend les tirages indépendants et les probabilités identiques au second tirage.
\begin{popfigure}
\begin{center}
\begin{tikzpicture}[grow=right, level distance=3.2cm, level 1/.style={sibling distance=2.2cm}, level 2/.style={sibling distance=1.2cm}, every node/.style={font=\small}]
\node {}
  child { node {$\bar{R}$}
    child { node {$\bar{R}$ \quad $0{,}7\times 0{,}7=0{,}49$} edge from parent node[below] {$0{,}7$} }
    child { node {$R$ \quad $0{,}7\times 0{,}3=0{,}21$} edge from parent node[above] {$0{,}3$} }
    edge from parent node[below] {$0{,}7$} }
  child { node {$R$}
    child { node {$\bar{R}$ \quad $0{,}3\times 0{,}7=0{,}21$} edge from parent node[below] {$0{,}7$} }
    child { node {$R$ \quad $0{,}3\times 0{,}3=0{,}09$} edge from parent node[above] {$0{,}3$} }
    edge from parent node[above] {$0{,}3$} };
\end{tikzpicture}
\end{center}
\legende{Arbre de probabilités du tirage de deux boules avec remise.}
\end{popfigure}
La probabilité d'obtenir deux boules rouges est le produit le long du chemin $R$--$R$ :
\[P(R\cap R)=0{,}3\times 0{,}3=0{,}09.\]
\end{enumerate}
\begin{attention}[Point de vigilance]
Sur un arbre, on \emph{multiplie} les probabilités le long d'un chemin et on \emph{additionne} les probabilités des chemins menant au même événement.
\end{attention}
\end{corrige}

\begin{corrige}[Exercice 6 — Gain à la tombola]
\begin{enumerate}
\item Le gain algébrique est le gain brut diminué des \SI{200}{FCFA} du ticket : $X$ prend les valeurs $5000-200=4800$, $1000-200=800$ et $0-200=-200$.
\begin{center}
\begin{tabularx}{\linewidth}{|l|X|X|X|}
\hline
\rowcolor{popblueL} Valeur $x_i$ & $4800$ & $800$ & $-200$ \\
\hline
$P(X=x_i)$ & $\dfrac{2}{100}=0{,}02$ & $\dfrac{8}{100}=0{,}08$ & $\dfrac{90}{100}=0{,}9$ \\
\hline
\end{tabularx}
\end{center}
Vérification : $0{,}02+0{,}08+0{,}9=1$.
\item $E(X)=4800\times 0{,}02+800\times 0{,}08+(-200)\times 0{,}9=96+64-180=-20$.

En moyenne, un joueur \textbf{perd \SI{20}{FCFA} par ticket} : le jeu est défavorable au joueur (et rapporte en moyenne \SI{20}{FCFA} par ticket à l'organisateur).
\item On calcule d'abord :
\[E(X^2)=4800^2\times 0{,}02+800^2\times 0{,}08+(-200)^2\times 0{,}9=460800+51200+36000=548000.\]
Par la formule de \cle{Koenig-Huygens} :
\[V(X)=E(X^2)-[E(X)]^2=548000-(-20)^2=548000-400=547600.\]
\[\sigma(X)=\sqrt{547600}=740\ \text{FCFA (valeur exacte)}.\]
L'écart-type, très grand devant l'espérance, traduit la forte dispersion des gains autour de la moyenne : la plupart des joueurs perdent \SI{200}{FCFA}, quelques-uns gagnent gros.
\end{enumerate}
\end{corrige}

\begin{corrige}[Exercice 7 — Contrôle qualité chez un opérateur]
\begin{enumerate}
\item Pour chaque appel, on répète une épreuve de Bernoulli : succès $S$ « l'appel dure plus de 5 minutes », de probabilité $p=0{,}4$. On répète $n=5$ fois cette épreuve, de manière identique et indépendante (hypothèse de l'énoncé). $X$ compte le nombre de succès : donc $X$ suit la loi binomiale $\mathcal{B}(5\,;\,0{,}4)$.
\item $P(X=2)=\dbinom{5}{2}\times 0{,}4^2\times 0{,}6^3=10\times 0{,}16\times 0{,}216=0{,}3456\approx 0{,}346$.
\item On passe par l'événement contraire « aucun appel ne dure plus de 5 minutes » :
\[P(X\ge 1)=1-P(X=0)=1-0{,}6^5=1-0{,}07776\approx 0{,}922.\]
\item $E(X)=np=5\times 0{,}4=2$. Sur un grand nombre d'échantillons de 5 appels, on observe en moyenne 2 appels durant plus de 5 minutes.
\end{enumerate}
\begin{attention}[Point de vigilance]
La justification du caractère binomial (deux issues, $n$ répétitions identiques et indépendantes, $X$ compte les succès) est une rédaction attendue ; ne pas l'oublier.
\end{attention}
\end{corrige}

\begin{corrige}[Exercice 8 — Fonction de répartition]
\begin{enumerate}
\item On rappelle que $F(x)=P(X\le x)$ et que $P(X=k)=F(k)-F(k^-)$ (saut de $F$ en $k$). Les sauts ont lieu en $1$, $2$, $3$, $4$ :
\begin{align*}
P(X=1)&=0{,}1-0=0{,}1 ; & P(X=2)&=0{,}4-0{,}1=0{,}3 ;\\
P(X=3)&=0{,}8-0{,}4=0{,}4 ; & P(X=4)&=1-0{,}8=0{,}2.
\end{align*}
Vérification : $0{,}1+0{,}3+0{,}4+0{,}2=1$.
\item $E(X)=1\times 0{,}1+2\times 0{,}3+3\times 0{,}4+4\times 0{,}2=0{,}1+0{,}6+1{,}2+0{,}8=2{,}7$.
\[E(X^2)=1\times 0{,}1+4\times 0{,}3+9\times 0{,}4+16\times 0{,}2=0{,}1+1{,}2+3{,}6+3{,}2=8{,}1.\]
\[V(X)=E(X^2)-[E(X)]^2=8{,}1-2{,}7^2=8{,}1-7{,}29=0{,}81.\]
\item $P(X>2)=P(X=3)+P(X=4)=0{,}4+0{,}2=0{,}6$ (ou directement $1-F(2)=1-0{,}4=0{,}6$).
\[P(1\le X\le 3)=P(X=1)+P(X=2)+P(X=3)=0{,}1+0{,}3+0{,}4=0{,}8.\]
\end{enumerate}
\begin{attention}[Point de vigilance]
$P(X>2)=1-F(2)$ car $F(2)=P(X\le 2)$ ; attention à ne pas écrire $1-F(3)$.
\end{attention}
\end{corrige}

\begin{corrige}[Exercice 9 — Tirage dans une urne — type Bac]
\textbf{Partie A.}
\begin{enumerate}
\item Au premier tirage : $P(B_1)=\dfrac{4}{10}=\dfrac{2}{5}$ et $P(N_1)=\dfrac{6}{10}=\dfrac{3}{5}$. Le tirage étant \textbf{sans remise}, la composition de l'urne change au second tirage (il ne reste que 9 boules) : par exemple, si la première boule est blanche, il reste 3 blanches sur 9, d'où $P_{B_1}(B_2)=\dfrac{3}{9}$.
\begin{popfigure}
\begin{center}
\begin{tikzpicture}[grow=right, level distance=3.4cm, level 1/.style={sibling distance=2.4cm}, level 2/.style={sibling distance=1.3cm}, every node/.style={font=\small}]
\node {}
  child { node {$N_1$}
    child { node {$N_2$ \quad $\dfrac{3}{5}\times\dfrac{5}{9}=\dfrac{1}{3}$} edge from parent node[below] {$\dfrac{5}{9}$} }
    child { node {$B_2$ \quad $\dfrac{3}{5}\times\dfrac{4}{9}=\dfrac{4}{15}$} edge from parent node[above] {$\dfrac{4}{9}$} }
    edge from parent node[below] {$\dfrac{3}{5}$} }
  child { node {$B_1$}
    child { node {$N_2$ \quad $\dfrac{2}{5}\times\dfrac{6}{9}=\dfrac{4}{15}$} edge from parent node[below] {$\dfrac{6}{9}$} }
    child { node {$B_2$ \quad $\dfrac{2}{5}\times\dfrac{3}{9}=\dfrac{2}{15}$} edge from parent node[above] {$\dfrac{3}{9}$} }
    edge from parent node[above] {$\dfrac{2}{5}$} };
\end{tikzpicture}
\end{center}
\legende{Arbre du tirage sans remise : les probabilités du second niveau dépendent du résultat du premier tirage.}
\end{popfigure}
\item 
\begin{enumerate}
\item $P(A)=P(B_1\cap B_2)=\dfrac{2}{5}\times\dfrac{3}{9}=\dfrac{6}{45}=\dfrac{2}{15}$.
\item $B=(B_1\cap B_2)\cup(N_1\cap N_2)$, réunion d'événements incompatibles :
\[P(B)=\dfrac{2}{15}+\dfrac{3}{5}\times\dfrac{5}{9}=\dfrac{2}{15}+\dfrac{1}{3}=\dfrac{2}{15}+\dfrac{5}{15}=\dfrac{7}{15}.\]
\item $P(C)=P(B_1\cap N_2)+P(N_1\cap B_2)=\dfrac{4}{15}+\dfrac{4}{15}=\dfrac{8}{15}$.
\end{enumerate}
\item $1-P(B)=1-\dfrac{7}{15}=\dfrac{8}{15}=P(C)$. C'était prévisible : « même couleur » et « couleurs différentes » sont deux événements \cle{contraires} l'un de l'autre, donc $C=\bar{B}$ et $P(C)=1-P(B)$.
\end{enumerate}

\textbf{Partie B.}
\begin{enumerate}
\setcounter{enumi}{3}
\item $X$ prend les valeurs $0$, $1$, $2$ :
\[P(X=0)=P(N_1\cap N_2)=\dfrac{1}{3}=\dfrac{5}{15} ;\qquad P(X=1)=P(C)=\dfrac{8}{15} ;\qquad P(X=2)=P(A)=\dfrac{2}{15}.\]
\begin{center}
\begin{tabularx}{\linewidth}{|l|X|X|X|}
\hline
\rowcolor{popblueL} $x_i$ & $0$ & $1$ & $2$ \\
\hline
$P(X=x_i)$ & $\dfrac{5}{15}$ & $\dfrac{8}{15}$ & $\dfrac{2}{15}$ \\
\hline
\end{tabularx}
\end{center}
Vérification : $\dfrac{5+8+2}{15}=1$.
\item $E(X)=0\times\dfrac{5}{15}+1\times\dfrac{8}{15}+2\times\dfrac{2}{15}=\dfrac{12}{15}=\dfrac{4}{5}=0{,}8$.
\[E(X^2)=0+1\times\dfrac{8}{15}+4\times\dfrac{2}{15}=\dfrac{16}{15},\]
\[V(X)=\dfrac{16}{15}-\left(\dfrac{4}{5}\right)^2=\dfrac{16}{15}-\dfrac{16}{25}=\dfrac{80-48}{75}=\dfrac{32}{75}\approx 0{,}4267,\qquad \sigma(X)=\sqrt{\dfrac{32}{75}}\approx 0{,}65.\]
\item Le joueur gagne $1000X$ FCFA et a misé $500$ FCFA : $Y=1000X-500$.

Par linéarité de l'espérance : $E(Y)=1000\,E(X)-500=1000\times 0{,}8-500=300$ FCFA.

$E(Y)=300\ne 0$ : le jeu n'est \textbf{pas équitable} ; il est favorable au joueur, qui gagne en moyenne \SI{300}{FCFA} par partie.
\end{enumerate}
\textbf{Barème indicatif (sur 5 pts) :} arbre correct avec probabilités conditionnelles (1 pt) ; $P(A)$, $P(B)$, $P(C)$ (1 pt) ; événements contraires justifiés (0,5 pt) ; loi de $X$ avec vérification (1 pt) ; $E(X)$ et $\sigma(X)$ (1 pt) ; $Y$ en fonction de $X$ et conclusion (0,5 pt).

\begin{attention}[Points de vigilance]
Citer explicitement « sans remise » pour justifier les probabilités du second niveau ; un jeu est équitable si et seulement si $E(Y)=0$ ; simplifier les fractions jusqu'au bout.
\end{attention}
\end{corrige}

\begin{corrige}[Exercice 10 — Sachets mal scellés à Bonabéri — type Bac]
\textbf{Partie A.}
\begin{enumerate}
\item Pour chaque sachet prélevé, on réalise une épreuve de Bernoulli : succès $S$ « le sachet est mal scellé », de probabilité $p=0{,}03$. On répète $n=20$ fois cette épreuve ; le prélèvement étant assimilé à un tirage avec remise, les répétitions sont identiques et indépendantes. $X$ compte le nombre de succès : donc $X\sim\mathcal{B}(20\,;\,0{,}03)$.
\item $P(X=0)=\dbinom{20}{0}\times 0{,}03^0\times 0{,}97^{20}=0{,}97^{20}\approx 0{,}5438$.
\item $P(X=1)=\dbinom{20}{1}\times 0{,}03\times 0{,}97^{19}=20\times 0{,}03\times \dfrac{0{,}97^{20}}{0{,}97}\approx 0{,}6\times \dfrac{0{,}5438}{0{,}97}\approx 0{,}3364$.
\item $P(X\ge 2)=1-P(X=0)-P(X=1)\approx 1-0{,}5438-0{,}3364=0{,}1198$.
\item Le nombre moyen est l'espérance : $E(X)=np=20\times 0{,}03=0{,}6$ sachet mal scellé par prélèvement de 20.
\end{enumerate}

\textbf{Partie B.}
\begin{enumerate}
\setcounter{enumi}{5}
\item Sur \num{10000} sachets, le nombre de sachets mal scellés suit $\mathcal{B}(10000\,;\,0{,}03)$, d'espérance $10000\times 0{,}03=300$. L'affirmation est abusive : certes $300<400$, mais il s'agit d'une \emph{moyenne} ; le nombre réel est aléatoire et peut dépasser 400 certains jours. Le mot « forcément » est donc incorrect.
\item La probabilité qu'aucun des 20 sachets ne soit défectueux est $(1-p)^{20}$. La condition de l'usine s'écrit donc $(1-p)^{20}>0{,}9$. La fonction $t\mapsto t^{1/20}$ étant croissante sur $[0\,;\,+\infty[$ :
\[1-p>0{,}9^{1/20}\approx 0{,}9947 \quad\Longleftrightarrow\quad p<1-0{,}9947=0{,}0053.\]
Le taux de sachets défectueux doit donc être au maximum $p\approx 0{,}005$, soit $0{,}5\,\%$.
\end{enumerate}
\textbf{Barème indicatif (sur 5 pts) :} justification complète de la loi binomiale (1 pt) ; $P(X=0)$ et $P(X=1)$ (1 pt) ; $P(X\ge 2)$ par le complément (0,5 pt) ; $E(X)$ (0,5 pt) ; critique de l'affirmation (1 pt) ; inéquation en $p$ et résolution (1 pt).

\begin{attention}[Points de vigilance]
L'assimilation à un tirage avec remise est la condition qui autorise l'indépendance — elle doit être citée ; pour « au moins deux », passer par $1-P(X=0)-P(X=1)$ ; dans la question 7, justifier le sens de l'inégalité par la croissance de la fonction puissance.
\end{attention}
\end{corrige}

\begin{corrige}[Exercice 11 — QCM du concours d'entrée — type Bac]
\textbf{Partie A.}
\begin{enumerate}
\item Chaque question est une épreuve de Bernoulli : succès « réponse correcte » de probabilité $p=\dfrac{1}{4}=0{,}25$ (une seule réponse juste parmi 4, choix au hasard). Les 8 réponses sont indépendantes et $X$ compte les succès : donc $X\sim\mathcal{B}(8\,;\,0{,}25)$.
\item $P(X=4)=\dbinom{8}{4}\times 0{,}25^4\times 0{,}75^4=70\times 0{,}00390625\times 0{,}31640625\approx 0{,}0865$.
\item Les événements $(X=k)$ sont incompatibles deux à deux :
\[P(X\ge 4)=P(X=4)+P(X=5)+P(X=6)+P(X=7)+P(X=8).\]
Or $P(X=8)=0{,}25^8\approx 0{,}0000153\approx 0{,}0000$, d'où
\[P(X\ge 4)\approx 0{,}0865+0{,}0231+0{,}0038+0{,}0004+0{,}0000=0{,}1138.\]
En répondant au hasard, un candidat a environ $11{,}4\,\%$ de chances d'obtenir au moins la moitié des bonnes réponses.
\item $E(X)=np=8\times 0{,}25=2$ ; \quad $V(X)=np(1-p)=8\times 0{,}25\times 0{,}75=1{,}5$ ; \quad $\sigma(X)=\sqrt{1{,}5}\approx 1{,}22$.

En répondant au hasard, un candidat obtient en moyenne 2 bonnes réponses sur 8.
\end{enumerate}

\textbf{Partie B.}
\begin{enumerate}
\setcounter{enumi}{4}
\item Le candidat a $X$ bonnes réponses et $8-X$ mauvaises :
\[N=X-0{,}25(8-X)=X-2+0{,}25X=1{,}25X-2.\]
Par linéarité de l'espérance : $E(N)=1{,}25\,E(X)-2=1{,}25\times 2-2=0{,}5$.

Le barème à points négatifs ramène la note moyenne d'un candidat répondant au hasard à $0{,}5$ sur 8, proche de zéro : il dissuade efficacement les réponses au hasard.
\item Le coût annuel pour la compagnie est une variable aléatoire valant \SI{400000}{FCFA} avec probabilité $0{,}05$ et $0$ avec probabilité $0{,}95$. Le coût moyen (espérance) par client est :
\[E=0{,}05\times 400000+0{,}95\times 0=20000\ \text{FCFA}.\]
Pour dégager une marge de $20\,\%$ sur ce coût moyen, la prime doit valoir au moins :
\[20000\times 1{,}20=24000\ \text{FCFA}.\]
La prime annuelle minimale est donc de \SI{24000}{FCFA} (déjà un multiple de 1000). C'est exactement le principe du calcul des primes pures en assurance : prime $=$ espérance du coût des sinistres, majorée d'un chargement.
\end{enumerate}
\textbf{Barème indicatif (sur 5 pts) :} loi binomiale justifiée (1 pt) ; $P(X=4)$ avec le coefficient binomial (0,5 pt) ; $P(X\ge 4)$ en sommant les valeurs (1 pt) ; $E(X)$, $\sigma(X)$ et interprétation (1 pt) ; $N$ en fonction de $X$ et $E(N)$ (1 pt) ; coût moyen et prime (0,5 pt).

\begin{attention}[Points de vigilance]
Ne pas oublier $\dbinom{8}{4}=70$ dans $P(X=4)$ ; « au moins la moitié » signifie $X\ge 4$ (et non $X>4$) ; la linéarité de l'espérance $E(aX+b)=aE(X)+b$ doit être nommée ; attention : $V(aX+b)=a^2V(X)$, le $b$ disparaît.
\end{attention}
\end{corrige}

\newpage

\coursec{Fiche bilan}

\coursec{Probabilités, variables aléatoires et loi binomiale}

\courssub{Vocabulaire des probabilités et équiprobabilité}

\begin{definition}[Expérience aléatoire, univers et événements]
Une \cle{expérience aléatoire} est une expérience dont on ne peut pas prédire le résultat avec certitude, mais dont on connaît l'ensemble de tous les résultats possibles.
\begin{itemize}
    \item L'ensemble de tous les résultats possibles s'appelle l'\cle{univers} et se note $\Omega$.
    \item Un \cle{événement} est une partie de $\Omega$ (un ensemble de résultats). On note généralement les événements par des lettres majuscules $A, B, C, \dots$
    \item Un \cle{événement élémentaire} est un événement qui ne contient qu'un seul résultat : $\{\omega\}$ avec $\omega \in \Omega$.
    \item L'\cle{événement contraire} (ou complémentaire) de $A$ est l'ensemble des résultats de $\Omega$ qui n'appartiennent pas à $A$ : $\bar{A} = \Omega \setminus A = \{\omega \in \Omega \mid \omega \notin A\}$.
    \item L'\cle{événement impossible} est $\varnothing$ ; l'\cle{événement certain} est $\Omega$.
\end{itemize}
\end{definition}

\begin{exemplebox}[Exemples concrets au Cameroun]
\begin{enumerate}
    \item \textbf{Lancer un dé équilibré à 6 faces} : $\Omega = \{1,2,3,4,5,6\}$. L'événement « obtenir un nombre pair » est $A = \{2,4,6\}$. Son contraire est $\bar{A} = \{1,3,5\}$.
    \item \textbf{Code de recharge MTN (14 chiffres)} : On choisit un code au hasard. $\Omega$ = ensemble des $10^{14}$ codes possibles. $A$ = « le code commence par 237 » (indicatif pays). $\mathrm{Card}(A) = 10^{11}$.
    \item \textbf{Enquête sur la consommation d'électricité (ENEO)} : On interroge un ménage au hasard à Yaoundé. $\Omega = \{\text{Abonné}, \text{Non abonné}\}$. $A = \{\text{Abonné}\}$.
\end{enumerate}
\end{exemplebox}

\begin{definition}[Équiprobabilité]
On dit qu'il y a \cle{équiprobabilité} lorsque tous les résultats élémentaires de $\Omega$ ont la même chance de se produire. Dans ce cas, pour tout événement $A \subset \Omega$ :
\[
P(A) = \frac{\mathrm{Card}(A)}{\mathrm{Card}(\Omega)} = \frac{\text{nombre de cas favorables}}{\text{nombre de cas possibles}}.
\]
\end{definition}

\begin{attention}[Condition d'application]
\textbf{N'utilise jamais cette formule sans avoir justifié l'équiprobabilité.} Au Bac, écrire « Par équiprobabilité » ou « Les issues sont équiprobables car le dé est parfait / le tirage est uniforme » est obligatoire pour obtenir le point de méthode.
\end{attention}

\begin{propriete}[Propriétés immédiates]
Pour tout événement $A$ :
\begin{itemize}
    \item $0 \le P(A) \le 1$.
    \item $P(\Omega) = 1$ ; $P(\varnothing) = 0$.
    \item $P(\bar{A}) = 1 - P(A)$.
\end{itemize}
\end{propriete}

\begin{exoresolu}[Probabilité et code promo Orange]
\textbf{Énoncé :} Orange Cameroun lance un jeu. Un code gagnant est un nombre à 4 chiffres (de 0000 à 9999). Un client compose un code au hasard.
\begin{enumerate}
    \item Définir $\Omega$ et calculer son cardinal.
    \item Calculer la probabilité que le code soit un palindrome (ex: 1221, 3443).
    \item Calculer la probabilité que le code ne contienne aucun chiffre 0.
\end{enumerate}
\textbf{Solution :}
\begin{enumerate}
    \item $\Omega = \{0000, 0001, \dots, 9999\}$. $\mathrm{Card}(\Omega) = 10^4 = 10\,000$. Équiprobabilité justifiée (choix au hasard).
    \item Un palindrome à 4 chiffres a la forme $abba$. $a \in \{0,\dots,9\}$ (10 choix), $b \in \{0,\dots,9\}$ (10 choix). $\mathrm{Card}(A) = 10 \times 10 = 100$. $P(A) = \frac{100}{10\,000} = 0,01$.
    \item $B$ = « aucun 0 ». Chaque chiffre a 9 possibilités $\{1,\dots,9\}$. $\mathrm{Card}(B) = 9^4 = 6561$. $P(B) = \frac{6561}{10\,000} = 0,6561$.
\end{enumerate}
\end{exoresolu}

\courssub{Opérations sur les événements et arbres de probabilités}

\begin{propriete}[Probabilité de la réunion]
Pour quels que soient deux événements $A$ et $B$ :
\[
P(A \cup B) = P(A) + P(B) - P(A \cap B).
\]
Si $A$ et $B$ sont \cle{incompatibles} ($A \cap B = \varnothing$), alors $P(A \cup B) = P(A) + P(B)$.
\end{propriete}

\begin{methode}[Arbre de probabilités pondéré]
Pour modéliser une expérience à plusieurs étapes :
\begin{enumerate}
    \item Dessiner les branches pour la première étape, étiqueter chaque branche par sa probabilité.
    \item Pour chaque issue, dessiner les branches de l'étape suivante avec leurs probabilités (qui peuvent dépendre de l'issue précédente, par exemple dans un tirage sans remise).
    \item \textbf{Règle du produit} : La probabilité d'un chemin (issue finale) est le produit des probabilités sur ses branches.
    \item \textbf{Règle de la somme} : La probabilité d'un événement est la somme des probabilités des chemins qui le réalisent.
    \item Vérifier que la somme des probabilités des issues finales vaut 1.
\end{enumerate}
\end{methode}





\begin{popfigure}
\begin{tikzpicture}[
    grow=right, sloped, level distance=3.5cm, sibling distance=1.5cm,
    edge from parent/.style={draw, thick, -{Stealth[length=3mm]}},
    every node/.style={font=\small, inner sep=2pt},
    leaf/.style={font=\small\bfseries, text=popdark}
]
\node[draw, rounded corners, fill=popblueL, text=popblue] {$P(M)=0,02$}
    child { node[draw, rounded corners, fill=popgreenL, text=popgreen] {$P(T_+|M)=0,95$}
        child { node[leaf] {$P(M \cap T_+) = 0,019$} }
        child { node[leaf] {$P(M \cap T_-) = 0,001$} }
    }
    child { node[draw, rounded corners, fill=poporangeL, text=poporange] {$P(T_+|\bar{M})=0,05$}
        child { node[leaf] {$P(\bar{M} \cap T_+) = 0,049$} }
        child { node[leaf] {$P(\bar{M} \cap T_-) = 0,931$} }
    };
\node[align=center, anchor=west] at (8,1.5) {\textbf{Arbre pondéré} :\\Produit sur les branches,\\Somme pour $P(T_+)$.};
\end{tikzpicture}
\legende{Arbre de probabilités pour l'exemple du dépistage.}
\end{popfigure}

\courssub{Variables aléatoires et loi de probabilité}

\begin{definition}[Variable aléatoire discrète]
Une \cle{variable aléatoire} (v.a.) $X$ est une fonction définie sur $\Omega$ à valeurs dans $\mathbb{R}$. Elle associe un nombre réel $x$ à chaque issue $\omega$.
On note $X(\Omega) = \{x_1, x_2, \dots, x_n\}$ l'ensemble des valeurs prises par $X$.
La \cle{loi de probabilité} de $X$ est le tableau :
\[
\begin{array}{c|cccc}
X & x_1 & x_2 & \dots & x_n \\
\hline
P & p_1 & p_2 & \dots & p_n
\end{array}
\quad \text{avec } p_i = P(X=x_i) > 0 \text{ et } \sum_{i=1}^n p_i = 1.
\]
\end{definition}

\begin{definition}[Fonction de répartition]
La \cle{fonction de répartition} de $X$ est la fonction $F : \mathbb{R} \to [0,1]$ définie par $F(x) = P(X \le x)$.
\begin{itemize}
    \item $F$ est croissante, continue à droite, avec des sauts aux $x_i$.
    \item $\lim_{x\to -\infty} F(x) = 0$, $\lim_{x\to +\infty} F(x) = 1$.
    \item $P(a < X \le b) = F(b) - F(a)$.
    \item $P(X = x_i) = F(x_i) - \lim_{x \to x_i^-} F(x)$ (hauteur du saut).
\end{itemize}
\end{definition}

\begin{exemplebox}[Loi de probabilité et fonction de répartition]
\textbf{Contexte :} On lance deux dés équilibrés. $X$ = somme des résultats.
$\Omega$ a 36 issues équiprobables.
\begin{center}
\begin{tabular}{|c|ccccccccccc|}
\hline
$X$ & 2 & 3 & 4 & 5 & 6 & 7 & 8 & 9 & 10 & 11 & 12 \\
\hline
$P$ & 1/36 & 2/36 & 3/36 & 4/36 & 5/36 & 6/36 & 5/36 & 4/36 & 3/36 & 2/36 & 1/36 \\
\hline
\end{tabular}
\end{center}
$F(x)$ est un escalier. $F(5) = P(X\le 5) = \frac{1+2+3+4}{36} = \frac{10}{36}$.
$P(5 < X \le 8) = F(8) - F(5) = \frac{26}{36} - \frac{10}{36} = \frac{16}{36}$.
\end{exemplebox}

\begin{popfigure}
\begin{tikzpicture}
\begin{axis}[
    width=\linewidth, height=6cm, axis lines=middle,
    xmin=1, xmax=13, ymin=0, ymax=1.1,
    xtick={2,3,4,5,6,7,8,9,10,11,12}, ytick={0,0.2,0.4,0.6,0.8,1},
    xlabel=$x$, ylabel=$F(x)$, title style={font=\small}, title={Fonction de répartition $F(x)$ de la somme de 2 dés},
    jump mark left, mark=*, mark size=2, thick, popcurve,
    clip=false
]
\addplot[popblue, const plot] coordinates {
    (1,0) (2,1/36) (3,3/36) (4,6/36) (5,10/36) (6,15/36) (7,21/36) (8,26/36) (9,30/36) (10,33/36) (11,35/36) (12,1) (13,1)
};
% Open circles for left limits (optional visual aid, tikz can do it but const plot handles jumps)
\end{axis}
\end{tikzpicture}
\legende{Graphique de la fonction de répartition (escalier).}
\end{popfigure}

\courssub{Espérance, variance et écart-type}

\begin{definition}[Espérance mathématique]
L'\cle{espérance} (ou moyenne) de $X$ est :
\[
E(X) = \sum_{i=1}^n x_i\, p_i.
\]
C'est la moyenne théorique pondérée par les probabilités. \textbf{Unité :} même unité que $X$.
\end{definition}

\begin{propriete}[Propriétés de l'espérance (Linéarité)]
Pour toutes v.a. $X, Y$ et réels $a, b$ :
\begin{itemize}
    \item $E(aX + b) = aE(X) + b$.
    \item $E(X+Y) = E(X) + E(Y)$ (toujours vrai, même sans indépendance).
\end{itemize}
\end{propriete}

\begin{definition}[Variance et écart-type]
La \cle{variance} mesure la dispersion autour de la moyenne :
\[
V(X) = E\bigl[(X - E(X))^2\bigr] = \sum_{i=1}^n p_i \bigl(x_i - E(X)\bigr)^2.
\]
Formule de calcul pratique (dite \cle{formule de König-Huygens}) :
\[
V(X) = E(X^2) - [E(X)]^2 \quad \text{avec } E(X^2) = \sum x_i^2 p_i.
\]
L'\cle{écart-type} est $\sigma(X) = \sqrt{V(X)}$. Il a la même unité que $X$.
\end{definition}

\begin{propriete}[Propriétés de la variance]
Pour $a, b \in \mathbb{R}$ :
\begin{itemize}
    \item $V(aX + b) = a^2 V(X)$ (le $b$ ne change pas la dispersion).
    \item Si $X$ et $Y$ sont \textbf{indépendantes} : $V(X+Y) = V(X) + V(Y)$.
\end{itemize}
\end{propriete}

\begin{attention}[Pièges classiques]
\begin{itemize}
    \item Oublier de vérifier $\sum p_i = 1$ avant de calculer $E(X)$.
    \item Confondre $E(X^2)$ et $[E(X)]^2$. $E(X^2) = \sum x_i^2 p_i$ ; $[E(X)]^2 = (\sum x_i p_i)^2$.
    \item Écrire $V(X) = E(X^2) - E(X)^2$ sans parenthèses : $E(X)^2$ est ambigu, écrire $[E(X)]^2$.
    \item Appliquer $V(X+Y)=V(X)+V(Y)$ sans vérifier l'indépendance.
\end{itemize}
\end{attention}

\begin{exoresolu}[Gain à un jeu de hasard (Loterie Nationale)]
\textbf{Énoncé :} Un ticket coûte 500 FCFA. Gains possibles : 0 FCFA (prob. 0,9), 1\,000 FCFA (prob. 0,08), 50\,000 FCFA (prob. 0,02). $X$ = gain net (gain - prix ticket).
\begin{enumerate}
    \item Donner la loi de $X$.
    \item Calculer $E(X)$ et l'interpréter.
    \item Calculer $V(X)$ et $\sigma(X)$.
\end{enumerate}
\textbf{Solution :}
\begin{enumerate}
    \item $X$ prend les valeurs : $-500$ (p=0,9), $500$ (p=0,08), $49\,500$ (p=0,02). Somme proba = 1.
    \item $E(X) = (-500)\times0,9 + 500\times0,08 + 49500\times0,02 = -450 + 40 + 990 = 580$ FCFA.
    \textbf{Interprétation :} En moyenne, le joueur gagne 580 FCFA par ticket. Le jeu est favorable au joueur (rare !).
    \item $E(X^2) = 250000\times0,9 + 250000\times0,08 + 2450250000\times0,02 = 225000 + 20000 + 49005000 = 49250000$.
    $V(X) = 49250000 - 580^2 = 49250000 - 336400 = 48913600$.
    $\sigma(X) = \sqrt{48913600} \approx 6994$ FCFA. Grande dispersion : risque élevé.
\end{enumerate}
\end{exoresolu}

\courssub{Épreuve de Bernoulli et schéma de Bernoulli}

\begin{definition}[Épreuve de Bernoulli]
Une \cle{épreuve de Bernoulli} est une expérience aléatoire n'ayant que deux issues possibles :
\begin{itemize}
    \item \textbf{Succès} (souvent noté 1), de probabilité $p \in [0,1]$.
    \item \textbf{Échec} (souvent noté 0), de probabilité $1-p$ (souvent noté $q$).
\end{itemize}
La v.a. $X$ associée ($X=1$ si succès, $0$ sinon) suit une \cle{loi de Bernoulli} de paramètre $p$ : $\mathcal{B}e(p)$.
$E(X)=p$, $V(X)=p(1-p)$.
\end{definition}

\begin{definition}[Schéma de Bernoulli]
On répète $n$ fois ($n \in \mathbb{N}^*$) une même épreuve de Bernoulli de paramètre $p$, de manière \textbf{indépendante}.
C'est un \cle{schéma de Bernoulli} de paramètres $(n, p)$.
\textbf{Conditions à vérifier (check-list) :}
\begin{enumerate}
    \item Nombre fixé $n$ d'épreuves.
    \item Deux issues par épreuve (Succès/Échec).
    \item Probabilité de succès $p$ identique à chaque épreuve.
    \item \textbf{Indépendance} des épreuves (souvent : tirages \emph{avec remise}).
\end{enumerate}
\end{definition}

\begin{attention}[Indépendance et « Sans remise »]
Si on tire \textbf{sans remise} dans une population finie, ce n'est \textbf{PAS} un schéma de Bernoulli (les probabilités changent, pas d'indépendance). C'est une loi hypergéométrique (hors programme Terminale, mais à savoir pour ne pas confondre). Au Bac, si « sans remise » $\rightarrow$ Arbre de probabilités ou dénombrement direct.
\end{attention}

\courssub{Loi binomiale et applications}

\begin{propriete}[Loi binomiale $\mathcal{B}(n,p)$]
Soit $X$ le nombre de succès lors d'un schéma de Bernoulli $(n,p)$. $X$ suit la \cle{loi binomiale} de paramètres $n$ et $p$, notée $X \sim \mathcal{B}(n,p)$.
Pour tout $k \in \{0, 1, \dots, n\}$ :
\[
P(X = k) = \binom{n}{k} p^k (1-p)^{n-k}.
\]
\textbf{Démonstration (exigible au Bac) :}
Un chemin avec exactement $k$ succès a une probabilité $p^k(1-p)^{n-k}$ (produit car indépendance).
Il y a $\binom{n}{k}$ façons de placer ces $k$ succès parmi les $n$ épreuves (choix des positions).
Les chemins sont incompatibles $\rightarrow$ on somme les probabilités.
\end{propriete}

\begin{propriete}[Espérance et variance de la loi binomiale]
Si $X \sim \mathcal{B}(n,p)$ :
\[
E(X) = np, \qquad V(X) = np(1-p), \qquad \sigma(X) = \sqrt{np(1-p)}.
\]
\textbf{Justification :} $X = X_1 + \dots + X_n$ où $X_i \sim \mathcal{B}e(p)$ indépendantes.
$E(X) = \sum E(X_i) = np$.
$V(X) = \sum V(X_i) = n p(1-p)$ (indépendance).
\end{propriete}

\begin{methode}[Calculer $P(X=k)$ ou $P(X \le k)$ avec la loi binomiale]
\begin{enumerate}
    \item Identifier $n$ (nb d'épreuves) et $p$ (prob. succès). Vérifier les 4 conditions du schéma.
    \item Écrire $X \sim \mathcal{B}(n,p)$.
    \item \textbf{Cas $P(X=k)$ :} Appliquer la formule $\binom{n}{k}p^k(1-p)^{n-k}$. Calculer $\binom{n}{k} = \frac{n!}{k!(n-k)!}$.
    \item \textbf{Cas $P(X \le k)$ ou $P(X \ge k)$ :} Somme des termes. Astuce : $P(X \ge 1) = 1 - P(X=0) = 1 - (1-p)^n$.
    \item Arrondir le résultat final (souvent à $10^{-4}$ ou $10^{-3}$) et conclure par une phrase.
\end{enumerate}
\end{methode}

\begin{savaistu}[Tables de la loi binomiale et Calculatrice]
Au Cameroun, les tables de la loi binomiale (valeurs de $P(X \le k)$) sont fournies dans les sujets du Bac pour $n \le 20$ et certaines valeurs de $p$. Sur calculatrice (mode STAT/DIST) : \texttt{BinomialCD} (cumulée) et \texttt{BinomialPD} (ponctuelle). \textbf{Attention :} Même avec la calculatrice, il faut écrire la formule et les paramètres $n, p, k$ sur la copie pour avoir les points de méthode.
\end{savaistu}

\begin{exoresolu}[Contrôle qualité usine de bouteilles (Brasseries du Cameroun)]
\textbf{Énoncé :} Sur une chaîne de remplissage, 4\% des bouteilles sont sous-remplies. On prélève au hasard 20 bouteilles (population grande $\approx$ avec remise). $X$ = nombre de bouteilles sous-remplies.
\begin{enumerate}
    \item Justifier que $X$ suit une loi binomiale. Préciser ses paramètres.
    \item Calculer $P(X=0)$ et $P(X=1)$.
    \item Calculer la probabilité qu'il y ait au moins 2 bouteilles défectueuses.
    \item Donner l'espérance et l'écart-type de $X$.
\end{enumerate}
\textbf{Solution :}
\begin{enumerate}
    \item $n=20$ épreuves (bouteilles), 2 issues (conforme/défectueuse), $p=0,04$ constant, indépendance (prélèvement dans grande population ou avec remise). $X \sim \mathcal{B}(20; 0,04)$.
    \item $P(X=0) = (0,96)^{20} \approx 0,4420$.
    $P(X=1) = \binom{20}{1} 0,04^1 0,96^{19} = 20 \times 0,04 \times 0,96^{19} \approx 0,3683$.
    \item $P(X \ge 2) = 1 - P(X=0) - P(X=1) \approx 1 - 0,4420 - 0,3683 = 0,1897$.
    \item $E(X) = 20 \times 0,04 = 0,8$ bouteille.
    $V(X) = 20 \times 0,04 \times 0,96 = 0,768$.
    $\sigma(X) = \sqrt{0,768} \approx 0,876$ bouteille.
\end{enumerate}
\end{exoresolu}

\begin{exoresolu}[Réussite au Baccalauréat (Série C/D/E/TI)]
\textbf{Énoncé :} Dans un lycée, la probabilité qu'un élève de Terminale C obtienne son Bac est 0,72. On choisit au hasard un groupe de 15 élèves. $Y$ = nombre de reçus.
\begin{enumerate}
    \item Modéliser. Quelle est la probabilité qu'exactement 10 élèves soient reçus ?
    \item Probabilité que la majorité (au moins 8) soit reçue ?
    \item Combien d'élèves faut-il interroger pour être sûr à 99\% d'avoir au moins un reçu ? (Utiliser $P(X \ge 1) \ge 0,99$).
\end{enumerate}
\textbf{Solution :}
\begin{enumerate}
    \item $Y \sim \mathcal{B}(15; 0,72)$.
    $P(Y=10) = \binom{15}{10} 0,72^{10} 0,28^5 = 3003 \times 0,72^{10} \times 0,28^5 \approx 0,206$.
    \item $P(Y \ge 8) = 1 - P(Y \le 7)$. Utiliser table ou calculatrice : $P(Y \le 7) \approx 0,057$. Donc $P(Y \ge 8) \approx 0,943$.
    \item On cherche $n$ tel que $1 - (1-0,72)^n \ge 0,99 \iff 0,28^n \le 0,01$.
    $n \ln(0,28) \le \ln(0,01) \iff n \ge \frac{\ln(0,01)}{\ln(0,28)} \approx \frac{-4,605}{-1,273} \approx 3,62$.
    Il faut au moins \textbf{4 élèves}.
\end{enumerate}
\end{exoresolu}

\begin{exercice}[Entraînement : Forage d'eau (MINEE)]
On réalise des forages pour trouver de l'eau potable dans une zone rurale. La probabilité de succès (trouver de l'eau) pour un forage est de 0,35. On prévoit 10 forages indépendants. $Z$ = nombre de forages fructueux.
\begin{enumerate}
    \item Justifier la loi suivie par $Z$ et donner ses paramètres.
    \item Calculer la probabilité qu'exactement 3 forages soient fructueux.
    \item Calculer la probabilité qu'au moins un forage soit fructueux.
    \item Déterminer l'espérance et l'écart-type de $Z$.
\end{enumerate}
\end{exercice}

\begin{corrige}[Exercice : Forage d'eau]
\begin{enumerate}
    \item Schéma de Bernoulli : $n=10$ forages, 2 issues (eau/sec), $p=0,35$ constant, indépendance. $Z \sim \mathcal{B}(10; 0,35)$.
    \item $P(Z=3) = \binom{10}{3} 0,35^3 0,65^7 = 120 \times 0,042875 \times 0,049022 \approx 0,252$.
    \item $P(Z \ge 1) = 1 - P(Z=0) = 1 - 0,65^{10} \approx 1 - 0,0135 = 0,9865$.
    \item $E(Z) = 10 \times 0,35 = 3,5$ forages.
    $V(Z) = 10 \times 0,35 \times 0,65 = 2,275$.
    $\sigma(Z) = \sqrt{2,275} \approx 1,508$ forages.
\end{enumerate}
\end{corrige}

\begin{bilan}

\textbf{1. Vocabulaire et Bases}
\begin{itemize}
    \item \cle{Univers $\Omega$}, \cle{Événement} (partie de $\Omega$), \cle{Événement contraire $\bar{A}$}.
    \item \cle{Équiprobabilité} : $P(A) = \frac{\mathrm{Card}(A)}{\mathrm{Card}(\Omega)}$ \textbf{(justifier toujours)}.
    \item $P(\bar{A}) = 1 - P(A)$ ; $P(A \cup B) = P(A) + P(B) - P(A \cap B)$.
    \item \cle{Arbre pondéré} : Produit sur un chemin, Somme pour un événement.
\end{itemize}

\textbf{2. Variable Aléatoire (V.A.)}
\begin{itemize}
    \item \cle{Loi de probabilité} : Tableau valeurs $x_i$ / probabilités $p_i$ ($\sum p_i = 1$).
    \item \cle{Fonction de répartition} $F(x) = P(X \le x)$ : croissante, escalier, sauts aux $x_i$.
    \item \cle{Espérance} $E(X) = \sum x_i p_i$. Linéaire : $E(aX+b)=aE(X)+b$.
    \item \cle{Variance} $V(X) = E(X^2) - [E(X)]^2 = \sum p_i(x_i-E(X))^2$. $\sigma(X)=\sqrt{V(X)}$.
    \item $V(aX+b) = a^2 V(X)$. Si $X,Y$ indépendantes : $V(X+Y)=V(X)+V(Y)$.
\end{itemize}

\textbf{3. Loi Binomiale $\mathcal{B}(n,p)$}
\begin{itemize}
    \item \cle{Schéma de Bernoulli} : $n$ épreuves, 2 issues, $p$ constant, \textbf{indépendance} (avec remise).
    \item $P(X=k) = \binom{n}{k} p^k (1-p)^{n-k}$ pour $k=0,\dots,n$.
    \item $E(X) = np$ ; $V(X) = np(1-p)$ ; $\sigma(X) = \sqrt{np(1-p)}$.
    \item Astuce : $P(X \ge 1) = 1 - (1-p)^n$.
\end{itemize}

\end{bilan}

\begin{aretenir}[Checklist Bac — Suis-je prêt ?]
\begin{itemize}
    \item[$\square$] Je définis $\Omega$, événement, événement élémentaire, événement contraire sans confusion.
    \item[$\square$] J'écris \textbf{« Par équiprobabilité »} avant d'appliquer $P(A)=\frac{\mathrm{Card}(A)}{\mathrm{Card}(\Omega)}$.
    \item[$\square$] J'utilise $P(\bar{A})=1-P(A)$ pour les « au moins un » ou « pas tous ».
    \item[$\square$] J'applique $P(A\cup B)=P(A)+P(B)-P(A\cap B)$ sans oublier l'intersection.
    \item[$\square$] Je construis un arbre pondéré : probabilités sur les branches, produit le long du chemin, somme des chemins pour un événement.
    \item[$\square$] Je détermine la loi d'une v.a. $X$ (tableau valeurs/probabilités) et je vérifie $\sum p_i = 1$.
    \item[$\square$] Je trace et j'interprète la fonction de répartition $F(x)=P(X\le x)$.
    \item[$\square$] Je calcule $E(X)$, $E(X^2)$, $V(X)=E(X^2)-[E(X)]^2$, $\sigma(X)$.
    \item[$\square$] Je reconnais un schéma de Bernoulli (4 critères) et j'écris $X \sim \mathcal{B}(n,p)$.
    \item[$\square$] Je calcule $P(X=k)$ avec la formule binomiale (attention aux exposants $k$ et $n-k$).
    \item[$\square$] Je cite $E(X)=np$ et $V(X)=np(1-p)$ pour la binomiale.
    \item[$\square$] Je conclus par une phrase répondant à la question, avec arrondi si nécessaire.
\end{itemize}
\end{aretenir}

\findecours

\end{document}
